% Lesson 5 — Grammar: Expressing purpose/result/obligation/preference, etc., let's/let's not + bare infinitive (the present/past perfect) — Anglais 1ère A — Cours signé OnBuch+
% Programme officiel MINESEC. Compiler avec : tectonic ENA05-grammar-expressing-purpose-result-obligation-preference.tex
\newcommand{\DOCMATIERE}{Anglais}
\newcommand{\DOCNIVEAU}{1ère A}
\newcommand{\DOCMODULE}{Chapter 1 : Family and Social Life}
\newcommand{\DOCLECON}{ENA05}
\newcommand{\DOCTITRE}{Lesson 5 — Grammar: Expressing purpose/result/obligation/preference, etc., let's/let's not + bare infinitive (the present/past perfect)}
% OnBuch+ — préambule « cours signés » ENRICHI (compilé avec tectonic / XeLaTeX)
% Système visuel « Pop » d'OnBuch+ : fond crème (jamais blanc), bordures encre
% épaisses, ombres « dures » décalées, titres Archivo Black en capitales.
% Version MATHÉMATIQUES : graphes (pgfplots/tikz), boîtes pédagogiques
% (définition, propriété, méthode, attention, le savais-tu, exercice résolu,
% exercice, TP, bilan), page de garde et signature OnBuch+ ancrée en bas.
%
% Macros attendues dans main.tex AVANT \input{preamble} :
%   \DOCMATIERE  \DOCNIVEAU  \DOCMODULE  \DOCLECON (numéro)  \DOCTITRE
\documentclass[11pt]{article}

\usepackage{fontspec}
\usepackage[english,french]{babel} % français = langue principale ; l'anglais via \eng{..} / textebox
\usepackage[a4paper,margin=2cm,top=3.4cm,bottom=2.8cm,headheight=1.4cm,footskip=1.3cm]{geometry}
\usepackage{amsmath,amssymb,amsfonts}
\usepackage{mathtools} % \xmapsto, \coloneqq... souvent utilisés par les modèles
\usepackage{cancel} % simplifications barrées (bilans de forces)
% \abs{z} (module, valeur absolue) et \norm{u} (norme), utilisés par les modèles
\providecommand{\abs}{}\let\abs\relax\DeclarePairedDelimiter{\abs}{\lvert}{\rvert}
\providecommand{\norm}{}\let\norm\relax\DeclarePairedDelimiter{\norm}{\lVert}{\rVert}
\usepackage[table]{xcolor} % [table] : fournit aussi \rowcolors (alternance de lignes)
\usepackage{pagecolor}
\usepackage{tikz}
\usetikzlibrary{arrows.meta,positioning,calc,shapes.geometric,shapes.misc,shapes.arrows,decorations.pathmorphing,decorations.pathreplacing,decorations.markings,patterns,shadows,mindmap,trees,fit,backgrounds,matrix,intersections,angles,quotes}
\usepackage{pgfplots}
\usepackage[version=4]{mhchem} % équations nucléaires \ce{^{235}_{92}U}
\usepackage[european,siunitx]{circuitikz} % schémas électriques
\pgfplotsset{compat=1.18}
\usepgfplotslibrary{fillbetween,groupplots} % aires entre courbes (calcul intégral)
\usepackage{enumitem}
\usepackage{fancyhdr}
% [most] charge toutes les bibliothèques tcolorbox (listings, minted, raster,
% documentation...) et gonfle inutilement la mémoire TeX sur un document long
% (~300 boîtes) : on ne charge que ce qui est réellement utilisé.
\usepackage[skins,breakable]{tcolorbox}
\usepackage{graphicx}
\usepackage{colortbl}
\usepackage{booktabs}
\usepackage{tabularx}
\usepackage{diagbox} % case d'en-tête diagonale des tableaux à double entrée
% Colonnes extensibles centrée / gauche / droite (C, L, R), souvent utilisées par les modèles
\newcolumntype{C}{>{\centering\arraybackslash}X}
\newcolumntype{L}{>{\raggedright\arraybackslash}X}
\newcolumntype{R}{>{\raggedleft\arraybackslash}X}
\usepackage{array}
\usepackage{multicol}
\usepackage{multirow}
\usepackage{siunitx}
% parse-numbers=false : accepte aussi \SI{2,5 \times 10^{-3}}{...}
\sisetup{locale=FR,output-decimal-marker={,},group-minimum-digits=4,parse-numbers=false}
\DeclareSIUnit\ohm{\ensuremath{\symup{\Omega}}} % U+2126 absent de la police
\DeclareSIUnit\division{div} % réglages d'oscilloscope (ms/div, V/div)
\DeclareSIUnit\tr{tr} % tours (vitesse de rotation, ex. \SI{1500}{\tr\per\minute})
\DeclareSIUnit\molar{M} % concentration molaire (ex. \SI{3}{\molar}, \SI{1}{\micro\molar})
\DeclareSIUnit\an{an} % année (le modèle confond parfois avec \year, un compteur TeX) — ex. \SI{5}{\kilo\meter\per\an}
\DeclareSIUnit\FCFA{FCFA} % franc CFA (ex. \SI{500}{\FCFA\per\kilo\gram})
\DeclareSIUnit\degreeBrix{\textdegree Brix} % teneur en sucre d'un sirop/jus (réfractométrie)
\DeclareSIUnit\month{mois} % le modèle utilise parfois \month, un compteur TeX, comme unité de durée
\usepackage{qrcode}
% \degree : commande fréquemment utilisée par les modèles pour un angle ou une
% température en dehors de \SI{}{} (non fournie par les paquets chargés).
\providecommand{\degree}{^{\circ}}
% \qed (fin de démonstration), utilisé par les modèles sans amsthm
\providecommand{\qed}{\hfill$\square$}
% \barre : double barre des valeurs interdites dans un tableau de variations (tkz-tab)
\providecommand{\barre}{\ensuremath{\|}}
% environnement proof (amsthm non chargé) utilisé par les modèles
\ifdefined\proof\else\newenvironment{proof}[1][Démonstration]{\par\noindent\textit{#1.}\ }{\qed\par}\fi
\usepackage{lastpage}
\usepackage{hyperref}
\hypersetup{hidelinks}

% ============================================================
% Palette « Pop » OnBuch+ — mêmes hex que la classe P (plus_ui.dart)
% ============================================================
\definecolor{poporange}{HTML}{F59321}
\definecolor{popdark}{HTML}{E07A0C}
\definecolor{popcream}{HTML}{FFF6E8}
\definecolor{popink}{HTML}{1C1714}
\definecolor{poppurple}{HTML}{7A5AE0}
\definecolor{popgreen}{HTML}{2E9E6A}
\definecolor{poppink}{HTML}{E0457B}
\definecolor{popblue}{HTML}{2F7DE1}
\definecolor{popgold}{HTML}{C98A12}
\definecolor{popmuted}{HTML}{8A7F76}
\definecolor{popline}{HTML}{EADFD2}
\definecolor{popgray}{HTML}{8A7F76}
\colorlet{popgrey}{popgray}
% Alias tolérants (le modèle nomme parfois les couleurs différemment)
\colorlet{popwhite}{white}
\colorlet{popblack}{popink}
\colorlet{popred}{poppink}
\colorlet{popyellow}{popgold}
% Teintes claires (fonds de tableaux/encadrés) — le modèle nomme parfois
% les couleurs avec un suffixe L (ex. popblueL) sans les avoir définies.
\colorlet{poporangeL}{poporange!20}
\colorlet{popdarkL}{popdark!20}
\colorlet{poppurpleL}{poppurple!20}
\colorlet{popgreenL}{popgreen!20}
\colorlet{poppinkL}{poppink!20}
\colorlet{popblueL}{popblue!20}
\colorlet{popgoldL}{popgold!20}
\colorlet{popredL}{popred!20}
\colorlet{popgrayL}{popgray!20}
% Teintes claires pour fonds de tableaux / zones
\colorlet{popblueL}{popblue!10!white}
\colorlet{poporangeL}{poporange!14!white}
\colorlet{popgreenL}{popgreen!12!white}
\colorlet{poppurpleL}{poppurple!12!white}
\colorlet{poppinkL}{poppink!12!white}
\colorlet{popgoldL}{popgold!14!white}

\pagecolor{popcream}

% ============================================================
% Typographie
% ============================================================
\defaultfontfeatures{Ligatures=TeX}
\setmainfont{PlusJakartaSans-Medium.ttf}[Path=../../../fonts/,BoldFont=PlusJakartaSans-Bold.ttf]
% Symboles, API et emojis absents de la police principale : repli sur DejaVu Sans (caractères actifs)
\newfontface\symfont{DejaVuSans.ttf}[Path=../../../fonts/]
\makeatletter
\newcommand\symactive[1]{\catcode#1=13 \begingroup\lccode`\~=#1 \lowercase{\endgroup\def~}{{\symfont\char#1}}}
\newcommand\symblank[1]{\catcode#1=13 \begingroup\lccode`\~=#1 \lowercase{\endgroup\def~}{}}
\newcommand\symhyph[1]{\catcode#1=13 \begingroup\lccode`\~=#1 \lowercase{\endgroup\def~}{\mbox{-}}}
\newcount\symc
\newcommand\symrange[2]{\symc=#1 \@whilenum\symc<#2 \do{\expandafter\symactive\expandafter{\the\symc}\advance\symc by 1 }}
\newcommand\symblankrange[2]{\symc=#1 \@whilenum\symc<#2 \do{\expandafter\symblank\expandafter{\the\symc}\advance\symc by 1 }}
\makeatother
\symrange{"0250}{"02FF}\symactive{"2713}\symactive{"2714}\symactive{"2717}\symactive{"2718}\symactive{"2610}\symactive{"2611}\symactive{"2612}
\symhyph{"2011}\symhyph{"2E1A}\symblankrange{"1F300}{"1FAFF}\symblank{"FE0F}
% Maths en sans-serif (Fira Math) avec chiffres et lettres droites de Plus
% Jakarta, pour que les formules se fondent dans le texte.
\usepackage[math-style=ISO,bold-style=ISO]{unicode-math}
\setmathfont{FiraMath-Regular.otf}[Path=../../../fonts/]
\setmathfont{PlusJakartaSans-Medium.ttf}[Path=../../../fonts/,range={up/num,up/{Latin,latin},\mathup}]
\usepackage{icomma} % virgule décimale sans espace en mode math
% Caractères mathématiques Unicode absents des polices de texte (Plus Jakarta,
% Archivo Black) : rendus par leur équivalent mathématique, en texte comme en
% formule — sinon ils s'affichent en carré vide (ex. « Arithmétique dans ℤ »).
\usepackage{newunicodechar}
\newunicodechar{ℝ}{\ensuremath{\mathbb{R}}}
\newunicodechar{ℤ}{\ensuremath{\mathbb{Z}}}
\newunicodechar{ℂ}{\ensuremath{\mathbb{C}}}
\newunicodechar{ℕ}{\ensuremath{\mathbb{N}}}
\newunicodechar{ℚ}{\ensuremath{\mathbb{Q}}}
\newunicodechar{𝔻}{\ensuremath{\mathbb{D}}}
\newunicodechar{α}{\ensuremath{\alpha}}
\newunicodechar{β}{\ensuremath{\beta}}
\newunicodechar{γ}{\ensuremath{\gamma}}
\newunicodechar{δ}{\ensuremath{\delta}}
\newunicodechar{ε}{\ensuremath{\varepsilon}}
\newunicodechar{θ}{\ensuremath{\theta}}
\newunicodechar{λ}{\ensuremath{\lambda}}
\newunicodechar{μ}{\ensuremath{\mu}}
\newunicodechar{σ}{\ensuremath{\sigma}}
\newunicodechar{τ}{\ensuremath{\tau}}
\newunicodechar{φ}{\ensuremath{\varphi}}
\newunicodechar{ψ}{\ensuremath{\psi}}
\newunicodechar{ω}{\ensuremath{\omega}}
\newunicodechar{Σ}{\ensuremath{\Sigma}}
% majuscules produites par \MakeUppercase dans les titres (α-aminés -> Α)
\newunicodechar{Α}{\ensuremath{\alpha}}
\newunicodechar{Β}{\ensuremath{\beta}}
\newunicodechar{Γ}{\ensuremath{\Gamma}}
\newunicodechar{Δ}{\ensuremath{\Delta}}
\newunicodechar{Ε}{\ensuremath{\varepsilon}}
\newunicodechar{Θ}{\ensuremath{\Theta}}
\newunicodechar{Λ}{\ensuremath{\Lambda}}
\newunicodechar{Μ}{\ensuremath{\mu}}
\newunicodechar{Τ}{\ensuremath{\tau}}
\newunicodechar{Φ}{\ensuremath{\Phi}}
\newunicodechar{Ψ}{\ensuremath{\Psi}}
\newunicodechar{Ω}{\ensuremath{\Omega}}
\newunicodechar{↦}{\ensuremath{\mapsto}}
\newunicodechar{∈}{\ensuremath{\in}}
\newunicodechar{∉}{\ensuremath{\notin}}
\newunicodechar{∧}{\ensuremath{\wedge}}
\newunicodechar{⇔}{\ensuremath{\Leftrightarrow}}
\newunicodechar{⟺}{\ensuremath{\Longleftrightarrow}}
\newunicodechar{⇒}{\ensuremath{\Rightarrow}}
\newunicodechar{⟹}{\ensuremath{\Longrightarrow}}
\newunicodechar{∘}{\ensuremath{\circ}}
\newunicodechar{∪}{\ensuremath{\cup}}
\newunicodechar{∩}{\ensuremath{\cap}}
\newunicodechar{≡}{\ensuremath{\equiv}}
\newunicodechar{⊂}{\ensuremath{\subset}}
\newunicodechar{⊥}{\ensuremath{\perp}}
\newunicodechar{∃}{\ensuremath{\exists}}
\newunicodechar{∀}{\ensuremath{\forall}}
\newunicodechar{∛}{\ensuremath{\sqrt[3]{\,}}}
\newunicodechar{∜}{\ensuremath{\sqrt[4]{\,}}}
\newunicodechar{⃗}{\ensuremath{{}^{\to}}}
\newunicodechar{ⁿ}{\ifmmode^{n}\else\textsuperscript{\ensuremath{n}}\fi}
\newunicodechar{ˣ}{\ifmmode^{x}\else\textsuperscript{\ensuremath{x}}\fi}
\newunicodechar{ᵇ}{\ifmmode^{b}\else\textsuperscript{\ensuremath{b}}\fi}
\newunicodechar{ʳ}{\ifmmode^{r}\else\textsuperscript{\ensuremath{r}}\fi}
\newunicodechar{ᵘ}{\ifmmode^{u}\else\textsuperscript{\ensuremath{u}}\fi}
\newunicodechar{ʸ}{\ifmmode^{y}\else\textsuperscript{\ensuremath{y}}\fi}
\newunicodechar{ᵃ}{\ifmmode^{a}\else\textsuperscript{\ensuremath{a}}\fi}
\newunicodechar{ᵉ}{\ifmmode^{e}\else\textsuperscript{\ensuremath{e}}\fi}
\newunicodechar{ᵏ}{\ifmmode^{k}\else\textsuperscript{\ensuremath{k}}\fi}
\newunicodechar{ᵖ}{\ifmmode^{p}\else\textsuperscript{\ensuremath{p}}\fi}
\newunicodechar{ᴺ}{\ifmmode^{N}\else\textsuperscript{\ensuremath{N}}\fi}
\newunicodechar{ᶜ}{\ifmmode^{c}\else\textsuperscript{\ensuremath{c}}\fi}
\newunicodechar{ʰ}{\ifmmode^{h}\else\textsuperscript{\ensuremath{h}}\fi}
\newunicodechar{ᵠ}{\ifmmode^{q}\else\textsuperscript{\ensuremath{q}}\fi}
\newunicodechar{ₙ}{\ifmmode_{n}\else\textsubscript{\ensuremath{n}}\fi}
\newunicodechar{ₐ}{\ifmmode_{a}\else\textsubscript{\ensuremath{a}}\fi}
\newunicodechar{ᵢ}{\ifmmode_{i}\else\textsubscript{\ensuremath{i}}\fi}
\newunicodechar{ᵣ}{\ifmmode_{r}\else\textsubscript{\ensuremath{r}}\fi}
\newunicodechar{ₖ}{\ifmmode_{k}\else\textsubscript{\ensuremath{k}}\fi}
\newunicodechar{ᵦ}{\ifmmode_{\beta}\else\textsubscript{\ensuremath{\beta}}\fi}
\newunicodechar{ₓ}{\ifmmode_{x}\else\textsubscript{\ensuremath{x}}\fi}
\newfontfamily\popdisplay{ArchivoBlack-Regular.ttf}[Path=../../../fonts/]
\newfontfamily\poplogo{SpaceGrotesk-Bold.ttf}[Path=../../../fonts/]
\newfontfamily\popsemi{PlusJakartaSans-SemiBold.ttf}[Path=../../../fonts/]
\newfontfamily\popxbold{PlusJakartaSans-ExtraBold.ttf}[Path=../../../fonts/]
\newfontfamily\popxboldls{PlusJakartaSans-ExtraBold.ttf}[Path=../../../fonts/,LetterSpace=3.0]

\setlist[enumerate]{leftmargin=1.7em,itemsep=5pt,topsep=4pt}
\setlist[itemize]{leftmargin=1.7em,itemsep=4pt,topsep=4pt,label={\color{poporange}\textbullet}}
\setlength{\parindent}{0pt}
\setlength{\parskip}{5pt}
\renewcommand{\arraystretch}{1.25}

% pgfplots : style Pop par défaut
\pgfplotsset{
  every axis/.append style={axis line style={popink,line width=1pt},tick style={popink},
    grid style={popline,line width=0.5pt},label style={font=\small},tick label style={font=\footnotesize}},
  popcurve/.style={poporange,line width=1.8pt,no markers,smooth},
}
% Même style disponible dans un \draw TikZ simple (hors pgfplots)
\tikzset{popcurve/.style={poporange,line width=1.8pt}}
\tikzset{popgrid/.style={popline,very thin}}
\colorlet{popcurve}{poporange} % le nom du style est parfois utilisé comme couleur

% ============================================================
% Pill / squiggle
% ============================================================
\newcommand{\poppill}[3]{%
  \tikz[baseline=-0.55ex]\node[draw=popink,line width=1.2pt,rounded corners=2.6pt,%
    fill=#2,inner xsep=6.5pt,inner ysep=3pt]{%
    \popxboldls\fontsize{7}{7}\selectfont\color{#3}\MakeUppercase{#1}};%
}
\newcommand{\popsquiggle}[1]{%
  \tikz[baseline=-0.4ex]{
    \draw[#1,line width=2.6pt,line cap=round]
      plot [smooth] coordinates {(0,0.09) (0.34,0.01) (0.68,0.21) (1.02,0.01) (1.36,0.10)};
  }%
}
% Mot-clé surligné (marqueur orange sous le texte)
\newcommand{\cle}[1]{{\popxbold\color{popdark}#1}}

% ============================================================
% En-tête / pied
% ============================================================
\pagestyle{fancy}
\fancyhf{}
\renewcommand{\headrulewidth}{0pt}
\renewcommand{\footrulewidth}{0pt}
\lhead{\raisebox{-0.15ex}{\poplogo\fontsize{13}{13}\selectfont\color{popink}OnBuch\color{poporange}+}}
\rhead{\poppill{\DOCMATIERE\ \textbullet\ \DOCNIVEAU\ \textbullet\ Leçon \DOCLECON}{white}{popink}}
\lfoot{\popxbold\fontsize{7.6}{7.6}\selectfont\color{popmuted}\MakeUppercase{Cours signé OnBuch+}\ \textbullet\ {\popsemi\DOCTITRE}}
\rfoot{\tikz[baseline=-0.6ex]\node[rounded corners=8.5pt,fill=popink,minimum height=17pt,minimum width=17pt,inner xsep=5pt,inner sep=0]{\popxbold\fontsize{8}{8}\selectfont\color{popcream}\thepage\,/\,\pageref*{LastPage}};}

% ============================================================
% Titres
% ============================================================
\newcommand{\chaptitre}[2]{%
  \begin{tcolorbox}[enhanced,colback=poporange,colframe=popink,boxrule=2.6pt,arc=11pt,
      drop shadow=popink,left=15pt,right=15pt,top=12pt,bottom=11pt,before skip=3pt,after skip=18pt]
    \poppill{#2}{popink}{popcream}\par\vspace{9pt}
    {\raggedright\hyphenpenalty=10000\exhyphenpenalty=10000\popdisplay\fontsize{22}{24}\selectfont\color{popcream}\MakeUppercase{#1}\par}
  \end{tcolorbox}
}
% Section numérotée : pastille orange avec le numéro + titre Archivo
\newcounter{popsec}
\newcommand{\coursec}[1]{%
  \stepcounter{popsec}\setcounter{popsub}{0}%
  \par\vspace{16pt}\noindent
  \begin{minipage}{\linewidth}
  \tikz[baseline=-0.7ex]\node[draw=popink,line width=1.6pt,fill=poporange,rounded corners=4pt,minimum size=22pt,inner sep=0]{\popdisplay\fontsize{12}{12}\selectfont\color{popcream}\thepopsec};%
  \hspace{8pt}{\popdisplay\fontsize{15}{17}\selectfont\color{popink}\MakeUppercase{#1}}\par
  \vspace{4pt}\noindent{\color{poporange}\rule{36pt}{3.6pt}}
  \end{minipage}\par\nopagebreak\vspace{8pt}
}
\newcounter{popsub}
\newcommand{\courssub}[1]{%
  \stepcounter{popsub}%
  \par\vspace{9pt}\noindent{\popxbold\fontsize{11.5}{13}\selectfont\color{popink}{\color{poporange}\thepopsec.\thepopsub}\hspace{6pt}#1}\par\nopagebreak\vspace{4pt}
}

% ============================================================
% Boîtes pédagogiques — même recette : fond blanc, bordure épaisse colorée,
% ombre dure encre, badge plein. Titre optionnel [#1].
% ============================================================
\tcbset{popbase/.style={breakable,enhanced,colback=white,boxrule=2.3pt,arc=9pt,
  shadow={3pt}{-3pt}{0pt}{popink},left=14pt,right=14pt,top=11pt,bottom=11pt,before skip=11pt,after skip=15pt,
  fontupper=\popsemi\fontsize{10.6}{14.5}\selectfont\color{popink},
  fontlower=\popsemi\fontsize{10.6}{14.5}\selectfont\color{popink}}}
% Coche dessinée (le glyphe ✓ n'existe pas dans les polices du document)
\newcommand{\popcheck}[1]{\tikz[baseline=-0.5ex]\draw[#1,line width=1.6pt,line cap=round,line join=round](-0.13,0.0)--(-0.04,-0.09)--(0.13,0.1);}
\providecommand{\coche}{\popcheck{popgreen}}
\providecommand{\resultat}[1]{\par\smallskip\noindent\fcolorbox{popgreen}{popgreenL}{\parbox{\dimexpr\linewidth-2\fboxsep-2\fboxrule\relax}{\textbf{Résultat.} #1}}\par\smallskip}
\newcommand{\popboxhead}[3]{\poppill{#1}{#2}{white}\ifx\relax#3\relax\else\hspace{6pt}{\popxbold\fontsize{10.6}{12}\selectfont\color{popink}#3}\fi\par\vspace{7pt}}

\newtcolorbox{aretenir}[1][]{popbase,colframe=popgreen,before upper={\popboxhead{À retenir}{popgreen}{#1}}}
% Texte en anglais (document de lecture, dialogue, extrait) : cadre
% d'étude. Un titre optionnel (source, auteur) s'affiche dans l'en-tête.
\newtcolorbox{textebox}[1][]{popbase,colframe=popblue,before upper={\popboxhead{Text}{popblue}{#1}\begin{otherlanguage}{english}},after upper={\end{otherlanguage}}}
% Marques ✓ / ✗ (forme correcte / fautive) souvent utilisées par les modèles
\providecommand{\cross}{\textcolor{poppink}{\ensuremath{\times}}}
\providecommand{\xmark}{\cross}\providecommand{\crossmark}{\cross}\providecommand{\cmark}{\ensuremath{\checkmark}}
% Ligne de réponse / justification à compléter (inventée par les modèles)
\providecommand{\justification}[1]{\par\noindent\textit{Justification~:} \dotfill\par}
% \textipa (package tipa non chargé) : la police ne couvre pas l'API -> simple passage
\providecommand{\textipa}[1]{#1}
% Soulignement (ulem non chargé) : \uline, \ul
\providecommand{\uline}[1]{\underline{#1}}\providecommand{\ul}[1]{\underline{#1}}
% \glossaire{\item ...} inventée par les modèles : liste de glossaire
\providecommand{\glossaire}[1]{\begin{itemize}#1\end{itemize}}
% \alert (beamer) inventée par les modèles : texte mis en évidence
\providecommand{\alert}[1]{\textbf{\textcolor{poppink}{#1}}}
% variantes ASCII d'ordinaux écrites en macro
\providecommand{\iere}{\textsuperscript{re}}\providecommand{\ieme}{\textsuperscript{e}}\providecommand{\ere}{\textsuperscript{re}}\providecommand{\eme}{\textsuperscript{e}}\providecommand{\er}{\textsuperscript{er}}\providecommand{\ie}{\textsuperscript{e}}
% Ordinaux français écrits en macro par les modèles : 1\ière, 3\ième
\newcommand{\ière}{\textsuperscript{re}}\newcommand{\ième}{\textsuperscript{e}}
% \exemple (inventée par les modèles) : étiquette « Exemple : »
\providecommand{\exemple}{\textbf{Exemple~:}\ }
% \square utilisable aussi hors mode math (cases à cocher)
\AtBeginDocument{\let\squareMath\square\renewcommand{\square}{\ensuremath{\squareMath}}}
% Mot ou phrase en anglais à l'intérieur d'un paragraphe français
\newcommand{\eng}[1]{\ifmmode\text{\textit{#1}}\else\begingroup\selectlanguage{english}\itshape #1\endgroup\fi}
\let\esp\eng % alias toléré
% flèches d'intonation inventées par les modèles
\providecommand{\textsearrow}{}\renewcommand{\textsearrow}{\ensuremath{\searrow}}\providecommand{\textnearrow}{}\renewcommand{\textnearrow}{\ensuremath{\nearrow}}
% \fr{..} : mot français (inventé par les modèles)
\providecommand{\fr}[1]{\textit{#1}}
\newtcolorbox{exemplebox}[1][]{popbase,colframe=poporange,before upper={\popboxhead{Exemple}{poporange}{#1}}}
\newtcolorbox{definition}[1][]{popbase,colframe=popblue,before upper={\popboxhead{Définition}{popblue}{#1}}}
\newtcolorbox{propriete}[1][]{popbase,colframe=poppurple,before upper={\popboxhead{Propriété}{poppurple}{#1}}}
\newtcolorbox{methode}[1][]{popbase,colframe=popgold,before upper={\popboxhead{Méthode}{popgold}{#1}}}
\newtcolorbox{attention}[1][]{popbase,colframe=poppink,before upper={\popboxhead{Attention}{poppink}{#1}}}
\newtcolorbox{experience}[1][]{popbase,colframe=popblue,colback=popblueL,before upper={\popboxhead{Expérience}{popblue}{#1}}}
% « Le savais-tu ? » : carte violette pleine (contexte, culture, Cameroun)
\newtcolorbox{savaistu}[1][]{popbase,colback=poppurple,colframe=popink,
  fontupper=\popsemi\fontsize{10.4}{14.2}\selectfont\color{white},
  before upper={\poppill{Le savais-tu\,?}{white}{poppurple}\ifx\relax#1\relax\else\hspace{6pt}{\popxbold\color{white}#1}\fi\par\vspace{7pt}}}
% Exercice résolu : énoncé en haut, \tcblower puis la solution sur fond crème
\newcounter{popexo}\newcounter{popexr}
\newtcolorbox{exoresolu}[1][]{popbase,colframe=poporange,colbacklower=popcream,
  segmentation style={popink,line width=1.2pt,dashed},
  before upper={\stepcounter{popexr}\popboxhead{Exercice résolu \thepopexr}{poporange}{#1}},
  before lower={\poppill{Solution}{popink}{popcream}\par\vspace{6pt}}}
% Exercice (non résolu en place) — la correction est dans la partie Corrigés
\newtcolorbox{exercice}[1][]{popbase,colframe=popink,
  before upper={\stepcounter{popexo}\popboxhead{Exercice \thepopexo}{popink}{#1}}}
% Corrigé
\newtcolorbox{corrige}[1][]{popbase,colframe=popgreen,colback=popgreenL,
  before upper={\popboxhead{Corrigé}{popgreen}{#1}}}
% Objectifs / prérequis / bilan
\newtcolorbox{objectifs}{popbase,colframe=popink,colback=poporangeL,
  before upper={\popboxhead{Objectifs de la leçon}{popink}{}}}
\newtcolorbox{prerequis}{popbase,colframe=popink,colback=popblueL,
  before upper={\popboxhead{Prérequis}{popblue}{}}}
\newtcolorbox{bilan}{popbase,colframe=popink,colback=white,boxrule=2.8pt,
  before upper={\popboxhead{Fiche bilan}{poporange}{L'essentiel à connaître pour l'examen}}}
% Figure encadrée avec légende
\newtcolorbox{popfigure}[1][]{enhanced,breakable=false,colback=white,colframe=popline,boxrule=1.4pt,arc=8pt,
  left=10pt,right=10pt,top=10pt,bottom=8pt,before skip=10pt,after skip=12pt,halign=center,
  lower separated=false,halign lower=center,
  fontlower=\popsemi\fontsize{9}{11.5}\selectfont\color{popmuted},
  #1}
\newcommand{\legende}[1]{\par\vspace{6pt}{\popsemi\fontsize{9}{11.5}\selectfont\color{popmuted}#1}}

% ============================================================
% Signature OnBuch+
% ============================================================
\newcommand{\poplogoinline}[1]{{\poplogo\fontsize{#1}{#1}\selectfont\color{popink}OnBuch\color{poporange}+}}
\newcommand{\signatureonbuch}{%
  \begin{tcolorbox}[enhanced,colback=popink,colframe=popink,boxrule=2.2pt,arc=10pt,
      drop shadow=poporange,left=14pt,right=14pt,top=10pt,bottom=10pt,before skip=0pt,after skip=0pt]
    \tikz[baseline=-0.7ex]\node[circle,fill=popgreen,draw=popcream,line width=1.3pt,minimum size=22pt,inner sep=0]{\popcheck{white}};%
    \hspace{9pt}\begin{minipage}[c]{0.62\linewidth}
      {\popxbold\fontsize{12}{13}\selectfont\color{popcream}Signé\ \textbullet\ {\poplogo OnBuch\color{poporange}+}}\par\vspace{2pt}
      {\raggedright\popsemi\fontsize{8}{10.5}\selectfont\color{popline}\MakeUppercase{Équipe pédagogique OnBuch+}\\\MakeUppercase{Conforme au programme officiel MINESEC}\par}
    \end{minipage}\hfill
    {\poplogo\fontsize{22}{22}\selectfont\color{popcream}OnBuch\color{poporange}+}
  \end{tcolorbox}%
}

% ============================================================
% Page de garde
% #1 titre · #2 matière · #3 niveau · #4 module · #5 n° leçon · #6 durée ·
% #7 description / objectifs (texte ou itemize)
% ============================================================
\newcommand{\pagegarde}[7]{%
  \thispagestyle{empty}
  \newgeometry{a4paper,margin=1.7cm,top=1.6cm,bottom=1.4cm}%
  \begin{tikzpicture}[remember picture,overlay]
    \fill[poporange!18] ([xshift=-3.2cm,yshift=-3.6cm]current page.north east) circle (4.6cm);
    \fill[poppurple!14] ([xshift=2.4cm,yshift=7.6cm]current page.south west) circle (3.8cm);
  \end{tikzpicture}%
  {\poplogo\fontsize{30}{30}\selectfont\color{popink}OnBuch\color{poporange}+}\hfill
  \raisebox{0.6ex}{\poppill{Cours signé \textbullet\ Premium}{popink}{popcream}}\par
  \vspace{1pt}\popsquiggle{poppurple}\par
  \vspace{8pt}
  \begin{tcolorbox}[enhanced,colback=poporange,colframe=popink,boxrule=2.8pt,arc=13pt,
      drop shadow=popink,left=18pt,right=18pt,top=12pt,bottom=13pt,after skip=10pt]
    \poppill{#2 \textbullet\ #3}{popink}{popcream}\hspace{5pt}\poppill{#4}{white}{popink}\par\vspace{8pt}
    {\popdisplay\fontsize{14}{14}\selectfont\color{popink}LEÇON #5}\par\vspace{4pt}
    {\raggedright\hyphenpenalty=10000\exhyphenpenalty=10000\popdisplay\fontsize{25}{27}\selectfont\color{popcream}\MakeUppercase{#1}\par}
    \vspace{8pt}
    {\popxbold\fontsize{9.5}{9.5}\selectfont\color{popink}\MakeUppercase{Durée conseillée : #6}}
  \end{tcolorbox}
  \begin{tcolorbox}[enhanced,colback=white,colframe=popink,boxrule=2.2pt,arc=10pt,
      drop shadow=popink,left=16pt,right=16pt,top=10pt,bottom=10pt,after skip=8pt]
    \poppill{Dans cette leçon}{poppurple}{white}\par\vspace{5pt}
    \popsemi\fontsize{9.3}{12.6}\selectfont\color{popink}#7
  \end{tcolorbox}
  \noindent\begin{minipage}[t]{0.487\linewidth}
    \begin{tcolorbox}[enhanced,colback=popgreen,colframe=popink,boxrule=2.2pt,arc=10pt,
        drop shadow=popink,left=11pt,right=11pt,top=10pt,bottom=10pt,height=3.1cm,valign=center]
      \tcbox[colback=white,colframe=popink,boxrule=1.3pt,arc=5pt,boxsep=0pt,nobeforeafter,
        left=3pt,right=3pt,top=3pt,bottom=3pt,valign=center]{\qrcode[height=1.9cm]{https://play.google.com/store/apps/details?id=com.luvvix.onbuch&hl=fr}}%
      \hspace{8pt}\begin{minipage}[c]{0.5\linewidth}
        {\popdisplay\fontsize{11}{12.5}\selectfont\color{white}\MakeUppercase{Télécharge OnBuch}}\par\vspace{4pt}
        {\raggedright\popsemi\fontsize{8.4}{11}\selectfont\color{white}Gratuit \textbullet\ Google Play\\Tuteur IA, annales, concours, résultats\par}
      \end{minipage}
    \end{tcolorbox}
  \end{minipage}\hfill
  \begin{minipage}[t]{0.487\linewidth}
    \begin{tcolorbox}[enhanced,colback=poppurple,colframe=popink,boxrule=2.2pt,arc=10pt,
        drop shadow=popink,left=11pt,right=11pt,top=10pt,bottom=10pt,height=3.1cm,valign=center]
      \tikz[baseline=-0.7ex]\node[circle,fill=white,draw=popink,line width=1.3pt,minimum size=1.6cm,inner sep=0]{\popdisplay\fontsize{24}{24}\selectfont\color{poppurple}+};%
      \hspace{8pt}\begin{minipage}[c]{0.6\linewidth}
        {\popdisplay\fontsize{11}{12.5}\selectfont\color{white}\MakeUppercase{Passe à OnBuch+}}\par\vspace{4pt}
        {\raggedright\popsemi\fontsize{8.4}{11}\selectfont\color{white}Tous les cours signés, Tuteur IA illimité, fiches PDF — onglet OnBuch+ de l'app\par}
      \end{minipage}
    \end{tcolorbox}
  \end{minipage}
  \vspace{8pt}
  \popcontenu
  \vfill
  \signatureonbuch
  \restoregeometry
  \newpage
}

% Rangée « Ce cours contient » (page de garde)
\newcommand{\poptile}[3]{% #1 couleur #2 gros texte #3 légende
  \begin{tcolorbox}[enhanced,colback=#1,colframe=popink,boxrule=2pt,arc=9pt,shadow={2.5pt}{-2.5pt}{0pt}{popink},
      left=6pt,right=6pt,top=9pt,bottom=9pt,halign=center,valign=center,height=1.75cm,width=\linewidth,nobeforeafter]
    {\popdisplay\fontsize{12.5}{13}\selectfont\color{white}\MakeUppercase{#2}}\par\vspace{3pt}
    {\centering\popsemi\fontsize{7.8}{9.5}\selectfont\color{white}#3\par}
  \end{tcolorbox}}
\newcommand{\popcontenu}{%
  \par\noindent{\popxboldls\fontsize{7.5}{7.5}\selectfont\color{popmuted}CE COURS SIGNÉ CONTIENT}\par\vspace{7pt}
  \noindent\begin{minipage}[t]{0.235\linewidth}\poptile{popblue}{Cours}{complet, illustré, pas à pas}\end{minipage}\hfill
  \begin{minipage}[t]{0.235\linewidth}\poptile{poporange}{Méthodes}{et exercices résolus}\end{minipage}\hfill
  \begin{minipage}[t]{0.235\linewidth}\poptile{poppink}{Exercices}{type examen + corrigés}\end{minipage}\hfill
  \begin{minipage}[t]{0.235\linewidth}\poptile{popgreen}{Fiche bilan}{l'essentiel à retenir}\end{minipage}\par}

% Sommaire
\newtcolorbox{sommaire}{enhanced,colback=white,colframe=popink,boxrule=2.2pt,arc=10pt,
  drop shadow=popink,left=18pt,right=18pt,top=13pt,bottom=13pt,before skip=6pt,after skip=20pt,
  before upper={\poppill{Sommaire}{poppurple}{white}\par\vspace{9pt}\popsemi\fontsize{10.6}{14.5}\selectfont\color{popink}}}

% Signature de fin de cours — poussée tout en bas de la dernière page
\newcommand{\findecours}{%
  \par\vfill\nopagebreak
  \begin{center}\popsquiggle{poporange}\end{center}\vspace{4pt}
  \signatureonbuch
}
\AtBeginDocument{\def\llbracket{\mathopen{[\mkern-3.5mu[}}\def\rrbracket{\mathclose{]\mkern-3.5mu]}}} % intervalles d'entiers (glyphe ⟦ absent de la police)
% Glyphes absents de Fira Math : redéfinis à partir de glyphes disponibles
\AtBeginDocument{%
  \def\setminus{\mathbin{\backslash}}\let\smallsetminus\setminus
  \def\vdots{\mathord{\vbox{\baselineskip3.2pt\lineskiplimit0pt\kern4pt\hbox{.}\hbox{.}\hbox{.}}}}%
  \def\ddots{\mathinner{\mkern1mu\raise6.5pt\hbox{.}\mkern2mu\raise3.6pt\hbox{.}\mkern2mu\raise0.7pt\hbox{.}\mkern1mu}}%
  \def\triangle{\mathord{\scalebox{0.9}{$\symup{\Delta}$}}}%
  \def\nabla{\mathord{\raisebox{\depth}{\scalebox{1}[-1]{$\symup{\Delta}$}}}}%
  \def\checkmark{\popcheck{popdark}}%
  \def\star{\mathbin{\ast}}\def\bigstar{\mathord{\ast}}%
  \def\diamond{\mathbin{\scalebox{0.6}{$\Diamond$}}}%
  \def\bigcup{\mathop{\mathchoice{\raisebox{-0.3em}{\scalebox{1.7}{$\cup$}}}{\scalebox{1.25}{$\cup$}}{\cup}{\cup}}\displaylimits}%
  \def\bigcap{\mathop{\mathchoice{\raisebox{-0.3em}{\scalebox{1.7}{$\cap$}}}{\scalebox{1.25}{$\cap$}}{\cap}{\cap}}\displaylimits}%
  \def\lesseqqgtr{\mathrel{\lessgtr}}\def\bigoplus{\mathop{\oplus}\displaylimits}%
}
\providecommand{\ssi}{si et seulement si} % abréviation fréquente
\AtBeginDocument{\providecommand{\wideparen}{\overparen}} % arc de cercle

%% FIGFIX-BEGIN v3 — correctifs globaux des figures (docs/onbuch-generation/tools/figfix)
\usepackage{adjustbox}\usepackage{etoolbox}
% toute figure TikZ/pgfplots plus large que la ligne est réduite à la largeur disponible
\BeforeBeginEnvironment{tikzpicture}{\begin{adjustbox}{max width=\linewidth}}
\AfterEndEnvironment{tikzpicture}{\end{adjustbox}}
% légendes pgfplots lisibles par-dessus les courbes
\pgfplotsset{every axis/.append style={legend style={fill=white,fill opacity=0.92,text opacity=1,draw=black!15,inner sep=3pt}}}
% étiquettes d'arêtes (syntaxe edge["..."]) avec fond blanc
\tikzset{every edge quotes/.append style={fill=white,inner sep=1.2pt}}
% bulles de cartes mentales : texte à la ligne dans la bulle, bulle un peu plus grande
\tikzset{every concept/.append style={text width=2.0cm,align=center,minimum size=2.3cm,inner sep=2pt}}
%% FIGFIX-END
\begin{document}

\pagegarde{Lesson 5 — Grammar: Expressing purpose/result/obligation/preference, etc., let's/let's not + bare infinitive (the present/past perfect)}{Anglais}{1ère A}{Chapter 1 : Family and Social Life}{ENA05}{2 h}{Cette leçon de grammaire étudie les moyens d'exprimer le but, la conséquence, l'obligation et la préférence en anglais, ainsi que la structure let's / let's not + base verbale et les temps present perfect et past perfect. Observation guidée, règles illustrées d'exemples traduits, comparaison avec le français et entraînement aux erreurs typiques des francophones.
\vspace{4pt}
\begin{multicols}{2}\small
\begin{itemize}
\item À la fin de cette leçon, je sais exprimer le but avec to + infinitive, in order to, so as to et so that
\item À la fin de cette leçon, je sais exprimer la conséquence avec so, so... that, such... that et therefore
\item À la fin de cette leçon, je sais exprimer l'obligation avec must, have to et l'interdiction avec mustn't
\item À la fin de cette leçon, je sais exprimer la préférence avec prefer, would rather et would rather... than
\item À la fin de cette leçon, je sais proposer une action avec let's / let's not + base verbale et réagir à une proposition
\item À la fin de cette leçon, je sais conjuguer et employer le present perfect et le past perfect avec leurs marqueurs (already, just, yet, since, for, before, after)
\end{itemize}\end{multicols}}

\begin{sommaire}
\begin{multicols}{2}
\begin{enumerate}
\item Observation : découvrir les structures en contexte
\item Exprimer le but (purpose)
\item Exprimer la conséquence (result)
\item Obligation, interdiction et préférence
\item Let's / Let's not + base verbale
\item Present perfect et past perfect
\item Entraînement et production
\item Activité d'intégration
\item Méthodes et exercices résolus
\item Exercices
\item Corrigés des exercices
\item Fiche bilan
\end{enumerate}
\end{multicols}
\end{sommaire}

\begin{objectifs}
\begin{itemize}
\item À la fin de cette leçon, je sais exprimer le but avec to + infinitive, in order to, so as to et so that
\item À la fin de cette leçon, je sais exprimer la conséquence avec so, so... that, such... that et therefore
\item À la fin de cette leçon, je sais exprimer l'obligation avec must, have to et l'interdiction avec mustn't
\item À la fin de cette leçon, je sais exprimer la préférence avec prefer, would rather et would rather... than
\item À la fin de cette leçon, je sais proposer une action avec let's / let's not + base verbale et réagir à une proposition
\item À la fin de cette leçon, je sais conjuguer et employer le present perfect et le past perfect avec leurs marqueurs (already, just, yet, since, for, before, after)
\item À la fin de cette leçon, je sais éviter les erreurs fréquentes des francophones (for to, must to, let's to go, I have seen him yesterday)
\item À la fin de cette leçon, je sais utiliser ces structures dans des phrases et un court texte sur la vie familiale et sociale
\end{itemize}
\end{objectifs}

\begin{prerequis}
\begin{itemize}
\item Les auxiliaires be, have, do et la base verbale (Seconde)
\item Le prétérit simple et ses emplois (Seconde)
\item Les modaux de base : can, must, should (Seconde)
\item La formation de l'infinitif avec et sans to (début d'année)
\end{itemize}
\end{prerequis}

\newpage

\coursec{Observation : découvrir les structures en contexte}

\courssub{Accroche : une matinée de samedi à Bamenda}

Imaginez la scène. Il est samedi matin à \eng{Bamenda}. \eng{Ngong} dit à sa sœur~: \eng{“I have finished my homework, so let's go to the market to buy some fish for Mum.”} Sa sœur répond~: \eng{“I'd rather stay at home — I have to clean the house first. We had already eaten when you called yesterday!”} En seulement trois phrases, les enfants expriment le \cle{but} (\eng{purpose}), la \cle{conséquence} (\eng{result}), l'\cle{obligation}, la \cle{préférence} et utilisent le \cle{present perfect} et le \cle{past perfect}. Aujourd'hui, vous allez apprendre à faire de même — correctement et naturellement.

\courssub{Un dialogue familial à Bamenda}

Lisez attentivement le dialogue suivant entre \eng{Manka} et \eng{Tabi}, sœur et frère. Ce texte contient \textbf{toutes} les structures que nous allons étudier : but, conséquence, obligation, préférence, \eng{let's/let's not}, \eng{present perfect} et \eng{past perfect}.

\begin{textebox}[Dialogue : Saturday Morning Chores]
Manka: I have just finished washing the plates, so let's visit Grandma.
Tabi: I'd rather rest a little. I had already swept the yard before you woke up.
Manka: OK, but we have to leave early in order to catch the first bus.
Tabi: Fine. Let's not forget the plantains she likes so much!
\end{textebox}

\textbf{Glossaire}\par
\begin{itemize}
    \item \eng{just} (adv.) : juste, à l'instant
    \item \eng{wash the plates} (v. phr.) : laver les assiettes
    \item \eng{visit} (v.) : rendre visite à
    \item \eng{rather} (adv.) : plutôt (marque la préférence)
    \item \eng{rest} (v.) : se reposer
    \item \eng{sweep / swept / swept} (v. irr.) : balayer
    \item \eng{yard} (n.) : cour, jardin
    \item \eng{wake up / woke up / woken up} (v. irr.) : se réveiller
    \item \eng{have to} (semi-aux.) : devoir (obligation externe)
    \item \eng{leave / left / left} (v. irr.) : partir
    \item \eng{in order to} (loc. conj.) : afin de, pour (but)
    \item \eng{catch / caught / caught} (v. irr.) : attraper, prendre (un bus)
    \item \eng{plantains} (n. pl.) : plantains (bananes plantains)
    \item \eng{forget / forgot / forgotten} (v. irr.) : oublier
\end{itemize}

\courssub{Méthode d'observation guidée}

Avant d'énoncer les règles, suivez cette démarche en quatre étapes pour devenir acteur de votre apprentissage.

\begin{methode}[Comment observer un texte grammatical]
\begin{enumerate}
    \item \textbf{Lisez le texte} une première fois pour la compréhension globale~: qui parle ? de quoi ? quel ton ?
    \item \textbf{Soulignez} (ou surlignez) toutes les structures qui expriment~:
    \begin{itemize}
        \item le \textbf{but} (\eng{purpose}) : \eng{to, in order to, so as to, so that}…
        \item la \textbf{conséquence} (\eng{result}) : \eng{so, so... that, such... that}…
        \item l'\textbf{obligation} / l'interdiction : \eng{have to, must, let's not}…
        \item la \textbf{préférence} : \eng{would rather, prefer}…
    \end{itemize}
    \item \textbf{Repérez les verbes} au \eng{present perfect} (\eng{have/has + participe passé}) et au \eng{past perfect} (\eng{had + participe passé}). Notez leurs \textbf{marqueurs temporels} (\eng{just, already, before, when, yesterday}…).
    \item \textbf{Faites des hypothèses} sur le sens de chaque structure grâce au contexte, \textbf{avant} de lire la règle formelle.
\end{enumerate}
\end{methode}

\courssub{Repérage et classification des exemples}

Appliquez la méthode. Voici le tableau de classement à quatre colonnes (but, conséquence, obligation, préférence) complété à partir du dialogue. Observez comment une même phrase peut combiner plusieurs notions (ex. \eng{so} lie une action achevée au \eng{present perfect} à une proposition introduite par \eng{let's}).

\begin{center}
\begin{tabularx}{\linewidth}{|X|X|X|X|}
\hline
\rowcolor{popblueL}
\textbf{Purpose (But)} & \textbf{Result (Conséquence)} & \textbf{Obligation / Interdiction} & \textbf{Preference (Préférence)} \\
\hline
\eng{in order to catch the first bus} & \eng{so let's visit Grandma} & \eng{we have to leave early} & \eng{I'd rather rest a little} \\
\hline
\eng{(let's visit Grandma) to see her} — infinitif de but implicite &  & \eng{Let's not forget the plantains} (suggestion négative = interdiction douce) &  \\
\hline
\end{tabularx}
\end{center}

\textbf{Repérage des temps parfaits (Perfect Tenses)}\par

\begin{center}
\begin{tabularx}{\linewidth}{|X|X|X|}
\hline
\rowcolor{popgreenL}
\textbf{Phrase extraite} & \textbf{Temps} & \textbf{Marqueur / Contexte} \\
\hline
\eng{I \textbf{have just finished} washing the plates} & Present Perfect & \eng{just} (action toute récente, conséquence présente) \\
\hline
\eng{I \textbf{had already swept} the yard \textbf{before} you \textbf{woke} up} & Past Perfect & \eng{already, before} (antériorité par rapport à un passé : \eng{woke up}) \\
\hline
\end{tabularx}
\end{center}

\courssub{Hypothèses de sens (avant la règle)}

À partir du contexte, essayez de formuler vous-même la fonction de chaque structure. Ne lisez pas encore les sections 2 à 6~!

\begin{exemplebox}[Vos hypothèses de travail]
\begin{itemize}
    \item \eng{so + proposition} $\rightarrow$ \emph{« Donc / alors » : lie une cause (present perfect) à une décision immédiate (let's).}
    \item \eng{in order to + base verbale} $\rightarrow$ \emph{« Afin de » : but précis, forme soutenue.}
    \item \eng{have to + base verbale} $\rightarrow$ \emph{« Devoir » : obligation venue de l'extérieur (horaire du bus).}
    \item \eng{I'd rather + base verbale} $\rightarrow$ \emph{« Je préférerais » : préférence personnelle, alternative.}
    \item \eng{Let's / Let's not + base verbale} $\rightarrow$ \emph{« Allons / N'allons pas » : suggestion inclusive (locuteur inclus).}
    \item \eng{have/has + participe passé (just/already)} $\rightarrow$ \emph{« Venir de / déjà » : passé récent lié au présent.}
    \item \eng{had + participe passé (already/before)} $\rightarrow$ \emph{« Avait déjà » : antériorité dans le passé (retour en arrière).}
\end{itemize}
\end{exemplebox}

\courssub{Carte mentale : les objectifs grammaticaux de cette leçon}

La figure ci-dessous résume visuellement les cinq piliers que nous allons maîtriser. Reproduisez-la dans votre cahier comme page de garde du chapitre.

\begin{popfigure}
\begin{tikzpicture}[
    mindmap,
    grow cyclic,
    every node/.style={concept, execute at begin node=\hskip0pt},
    concept color=popblue!70,
    level 1/.style={level distance=4.5cm, sibling angle=72},
    text=white,
    font=\small\bfseries
]
\node[concept, scale=1.2, align=center]{Grammar\\Goals}
    child[concept color=poporange] { node[concept, align=center]{Purpose\\But} }
    child[concept color=popgreen] { node[concept, align=center]{Result\\Conséquence} }
    child[concept color=poppurple] { node[concept, align=center]{Obligation\\Interdiction} }
    child[concept color=poppink] { node[concept, align=center]{Preference\\Préférence} }
    child[concept color=popdark] { node[concept, align=center]{Perfect\\Tenses} };
\end{tikzpicture}
\legende{Carte mentale des cinq notions grammaticales de la Leçon 5.}
\end{popfigure}

\courssub{Exercice résolu : analyse guidée du dialogue}

Appliquons la méthode sur une phrase complexe du dialogue.

\begin{exoresolu}[Analyse de « I had already swept the yard before you woke up »]
\textbf{Énoncé~:} Identifiez les temps, les formes, les deux actions, leur ordre chronologique réel et le marqueur de relation temporelle.
\tcblower
\textbf{Solution détaillée~:}
\begin{enumerate}
    \item \textbf{Temps~:} \eng{Past Perfect} (\eng{had swept}) pour la première action ; \eng{Past Simple} (\eng{woke up}) pour la seconde.
    \item \textbf{Formes~:} \eng{had + participe passé (swept)} ; verbe irrégulier au prétérit (\eng{woke up}).
    \item \textbf{Actions~:} A1 = \eng{swept the yard} (balayer la cour) ; A2 = \eng{woke up} (te réveiller).
    \item \textbf{Chronologie réelle~:} A1 se produit \textbf{AVANT} A2. Tabi a balayé \emph{avant} que Manka ne se réveille.
    \item \textbf{Marqueurs~:} \eng{before} (conjonction de subordination) + \eng{already} (adverbe insistant sur l'achèvement antérieur).
    \item \textbf{Traduction~:} « J'avais déjà balayé la cour avant que tu ne te réveilles. »
    \item \textbf{Piège francophone~:} Ne dites pas \eng{*I have already swept the yard before you woke up}. Le \eng{present perfect} ne s'utilise pas avec un repère de temps passé révolu (\eng{yesterday, before you woke up}).
\end{enumerate}
\end{exoresolu}

\courssub{Attention : les pièges à éviter dès maintenant}

\begin{attention}[Erreurs fréquentes des francophones — Observation]
\begin{itemize}
    \item \textbf{Traduction mot à mot de « pour »~:} « pour » = \eng{to} (but) \textbf{OU} \eng{for} (destinataire/bénéfice) \textbf{OU} \eng{in order to} (but soutenu) \textbf{OU} \eng{so that} (but avec sujet différent). \textit{Ex.~:} \eng{I go to the market to buy fish} (but) vs \eng{This fish is for Mum} (destinataire).
    \item \textbf{Confusion \eng{so} (conséquence) / \eng{so that} (but)~:} \eng{So} + proposition = conséquence (\eng{I am tired, so I am resting}) ; \eng{so that} + \eng{can/could} = but (\eng{I work hard so that I can succeed}).
    \item \textbf{\eng{Have to} vs \eng{Must}~:} \eng{have to} = obligation externe (loi, horaire, parent) — \eng{We have to leave early} (l'horaire du bus). \eng{Must} = obligation ressentie par le locuteur — \eng{I must call Grandma} (je sens que je devrais).
    \item \textbf{\eng{Let's} = \eng{Let us} (inclusif)~:} Ne pas confondre avec l'impératif \eng{Go!} (ordre). \eng{Let's go} = « Allons-y » (moi + toi). \eng{Let's not go} = « N'y allons pas ».
    \item \textbf{Present Perfect vs Past Simple~:} le \eng{present perfect} = passé lié au \textbf{présent} (\eng{just, already, yet, ever, never, since, for}) ; le \eng{past simple} = passé \textbf{révolu}, daté (\eng{yesterday, last week, in 2020, when…}).
    \item \textbf{Past Perfect = « le passé du passé »~:} il ne s'emploie \textbf{jamais} seul ; il nécessite un point de repère au \eng{past simple} (explicite ou implicite). \eng{*I had eaten} (incorrect seul) $\rightarrow$ \eng{I had eaten when you arrived}.
\end{itemize}
\end{attention}

\courssub{Exercice d'application immédiate}

\begin{exercice}[Repérage dans un nouveau micro-texte]
Lisez le court échange ci-dessous. Recopiez le tableau à quatre colonnes (Purpose, Result, Obligation, Preference) et classez chaque segment souligné. Identifiez aussi les temps parfaits.
\begin{quote}
\eng{Father: You have to finish your homework before you go out.}\\
\eng{Son: I'd rather play football. I have already done it, so can I go now?}\\
\eng{Father: OK, but don't forget to buy bread on your way.}
\end{quote}
\end{exercice}

\begin{corrige}[Exercice n°1]
\begin{center}
\begin{tabularx}{\linewidth}{|X|X|X|X|}
\hline
\rowcolor{popblueL}
\textbf{Purpose} & \textbf{Result} & \textbf{Obligation} & \textbf{Preference} \\
\hline
\eng{to buy bread} (infinitif de but : acheter du pain) & \eng{so can I go now?} & \eng{You have to finish your homework} & \eng{I'd rather play football} \\
\hline
 &  & \eng{don't forget} (interdiction / rappel ferme) &  \\
\hline
\end{tabularx}
\end{center}
\textbf{Remarque~:} \eng{before you go out} est une \textbf{subordonnée de temps} (et non de but) : elle indique \emph{quand} le devoir doit être accompli. \eng{On your way} est un complément de lieu (« sur ton chemin »).
\textbf{Temps parfaits~:} \eng{I \textbf{have already done} it} = \eng{Present Perfect} (marqueur \eng{already}, lien avec le présent : \eng{now}). Il n'y a pas de \eng{Past Perfect} dans cet échange.
\end{corrige}

\coursec{Exprimer le but (purpose)}

Dans la section précédente, nous avons observé des phrases tirées d'un texte sur la vie familiale et sociale. Vous avez remarqué des structures comme \eng{She works two jobs to pay the school fees} ou \eng{He speaks slowly so that his grandmother can understand him}. Toutes ces phrases répondent à la même question : \emph{pourquoi ?} — elles expriment le \cle{but} (en anglais \eng{purpose}). C'est la première grande fonction grammaticale de cette leçon. Nous allons maintenant l'étudier de façon systématique : d'abord en contexte, puis la règle, puis les pièges.

\courssub{Observation : le but dans un texte}

Lisez ce court texte et repérez toutes les expressions qui indiquent \emph{pourquoi} les personnages agissent.

\begin{textebox}[A family meeting in Yaoundé]
Last Saturday, the Mbappe family held a meeting in their living room in Yaoundé. Mr Mbappe, the father, had called everyone together \textbf{to discuss} the children's future. ``I work very hard \textbf{to give} you a good education,'' he said, ``and your mother sells vegetables at the market \textbf{in order to pay} for your school books. We do all this \textbf{so that} you \textbf{can have} a better life than ours.''

Amina, the eldest daughter, wants to become a doctor. She studies until midnight every night \textbf{so as to pass} her baccalauréat with good marks. She goes to bed early on Fridays \textbf{in order not to be} tired during Saturday classes. Her little brother Junior prefers football, but he has promised to work harder \textbf{so that} his parents \textbf{will not be} disappointed.

At the end of the meeting, Mrs Mbappe spoke softly \textbf{so as not to wake} the baby who was sleeping in the next room. Then she said: ``We are a family. We help each other \textbf{to succeed}. That is our purpose.''

\emph{Glossaire :} to hold a meeting (v.) : tenir une réunion ; school fees (n.) : frais de scolarité ; marks (n.) : notes ; disappointed (adj.) : déçu ; to whisper / to speak softly (v.) : chuchoter ; purpose (n.) : but, objectif.
\end{textebox}

\textbf{Questions de compréhension :}
\begin{enumerate}
\item Why did Mr Mbappe call a family meeting?
\item Why does Amina study until midnight?
\item Why does she go to bed early on Fridays?
\item Why did Mrs Mbappe speak softly?
\end{enumerate}

Vous avez repéré quatre constructions différentes : \eng{to discuss}, \eng{in order to pay}, \eng{so as to pass}, \eng{in order not to be}, \eng{so as not to wake} et \eng{so that you can have}. Toutes expriment le but. La différence entre elles est une question de \textbf{forme} et de \textbf{sujet}. Voyons cela en détail.

\courssub{Même sujet : to + base verbale}

\begin{propriete}[Le but avec le même sujet : to + base verbale]
Quand la personne qui agit et la personne visée par le but sont \textbf{la même}, on exprime le but avec :

\begin{center}
\textbf{to + base verbale} \quad (forme la plus courante, à l'oral comme à l'écrit)
\end{center}

\eng{I study hard to pass my exam.} « J'étudie sérieusement pour réussir mon examen. »

Variantes plus formelles, très utiles à l'écrit (rédactions, lettres officielles) :

\begin{center}
\textbf{in order to + base verbale} \qquad \textbf{so as to + base verbale}
\end{center}

\eng{She works two jobs in order to pay the school fees.} « Elle fait deux métiers afin de payer les frais de scolarité. »
\end{propriete}

\begin{exemplebox}[Exemples traduits — même sujet]
\begin{itemize}
\item \eng{He works hard to feed his family.} = Il travaille dur pour nourrir sa famille.
\item \eng{We left early to arrive on time.} = Nous sommes partis tôt pour arriver à l'heure.
\item \eng{Amina studies at night in order to pass her baccalauréat.} = Amina étudie la nuit afin de réussir son baccalauréat.
\item \eng{They saved money so as to buy a plot of land in Douala.} = Ils ont économisé de l'argent afin d'acheter un terrain à Douala.
\item \eng{I am learning English to get a better job.} = J'apprends l'anglais pour obtenir un meilleur emploi.
\end{itemize}
\end{exemplebox}

\textbf{Remarque de traduction.} Le français \emph{pour + infinitif} se traduit presque toujours par \eng{to + base verbale} en anglais, jamais par \eng{for to}. Nous y reviendrons dans les pièges.

\courssub{La forme négative : in order not to / so as not to}

\begin{propriete}[But négatif : « pour ne pas »]
Pour exprimer un but négatif (« pour ne pas… », « afin de ne pas… »), on place \textbf{not} devant \textbf{to}, mais uniquement avec les formes longues :

\begin{center}
\textbf{in order not to + BV} \qquad \textbf{so as not to + BV}
\end{center}

\eng{She whispered so as not to wake the baby.} « Elle a chuchoté pour ne pas réveiller le bébé. »

\eng{He left home early in order not to miss the bus.} « Il est parti tôt de la maison pour ne pas rater le bus. »

\textbf{Attention :} on ne dit pratiquement jamais \eng{not to} seul pour exprimer un but. \eng{I whispered not to wake the baby} est incorrect ou très maladroit ; dites \eng{so as not to wake} ou \eng{in order not to wake}.
\end{propriete}

\begin{exemplebox}[Exemples traduits — but négatif]
\begin{itemize}
\item \eng{She whispered so as not to wake the baby.} = Elle a chuchoté pour ne pas réveiller le bébé.
\item \eng{He drove slowly in order not to damage the car.} = Il a roulé lentement pour ne pas abîmer la voiture.
\item \eng{They hid the money so as not to attract thieves.} = Ils ont caché l'argent pour ne pas attirer les voleurs.
\end{itemize}
\end{exemplebox}

\courssub{Sujets différents : so that + proposition}

Jusqu'ici, le sujet de l'action et celui du but étaient identiques : \emph{Amina} étudie et c'est \emph{Amina} qui réussira. Mais que faire quand les sujets sont différents ? Reprenons une phrase du texte : \eng{We do all this so that you can have a better life.} Ici, « nous » faisons des efforts pour que « vous » ayez une vie meilleure. Deux sujets différents : \eng{to} seul devient impossible.

\begin{propriete}[Le but avec deux sujets différents : so that]
Quand le sujet de la phrase principale et celui du but sont \textbf{différents}, on utilise obligatoirement :

\begin{center}
\textbf{so that + sujet + can / could / will / would + base verbale}
\end{center}

\begin{itemize}
\item Présent ou futur : \eng{so that + sujet + can / will + BV}
\item Passé : \eng{so that + sujet + could / would + BV}
\end{itemize}

\eng{She speaks slowly so that everyone can understand.} « Elle parle lentement pour que tout le monde puisse comprendre. »

\eng{We left early so that we could arrive on time.} « Nous sommes partis tôt afin d'arriver à l'heure. » (ou : « afin que nous puissions arriver à l'heure »)
\end{propriete}

\begin{exemplebox}[Exemples traduits — sujets différents]
\begin{itemize}
\item \eng{We left early so that we could arrive on time.} = Nous sommes partis tôt afin d'arriver à l'heure.
\item \eng{The teacher repeats the lesson so that the students can take notes.} = Le professeur répète la leçon pour que les élèves puissent prendre des notes.
\item \eng{My father gave me money so that I would not walk to school.} = Mon père m'a donné de l'argent pour que je n'aille pas à l'école à pied.
\item \eng{Open the window so that fresh air can come in.} = Ouvre la fenêtre pour que l'air frais puisse entrer.
\end{itemize}
\end{exemplebox}

\textbf{Cas particulier.} Quand les deux sujets sont identiques, \eng{so that} reste possible mais est plus lourd : \eng{We left early so that we could arrive on time} = \eng{We left early to arrive on time}. À l'écrit, préférez la forme courte avec \eng{to} quand le sujet est le même.

\courssub{Schéma récapitulatif : choisir la bonne structure}

\begin{popfigure}
\begin{center}
\begin{tikzpicture}[
  node distance=0.9cm and 1.6cm,
  box/.style={draw=popblue, thick, rounded corners, fill=popblueL, align=center, inner sep=6pt},
  leaf/.style={draw=popgreen, thick, rounded corners, fill=popgreenL, align=center, inner sep=6pt},
  neg/.style={draw=poporange, thick, rounded corners, fill=poporangeL, align=center, inner sep=6pt},
  arr/.style={-{Stealth[length=3mm]}, thick, popdark}
]
\node[box, font=\bfseries, align=center] (but){Exprimer un BUT\\(pourquoi ?)};
\node[box, below left=1.1cm and 2.2cm of but] (same) {Même sujet};
\node[box, below right=1.1cm and 2.2cm of but] (diff) {Sujets différents};
\node[leaf, below=0.9cm of same, align=center] (to){\eng{to + BV}\\\eng{in order to + BV}\\\eng{so as to + BV}};
\node[neg, below=0.6cm of to, align=center] (not){négatif :\\\eng{in order not to + BV}\\\eng{so as not to + BV}};
\node[leaf, below=0.9cm of diff, align=center] (sothat){\eng{so that + sujet}\\\eng{+ can / will + BV}\\passé : \eng{could / would}};
\draw[arr] (but) -- (same);
\draw[arr] (but) -- (diff);
\draw[arr] (same) -- (to);
\draw[arr] (to) -- (not);
\draw[arr] (diff) -- (sothat);
\end{tikzpicture}
\end{center}
\legende{Arbre de décision : comment exprimer le but en anglais.}
\end{popfigure}

\courssub{Comparaison avec le français}

\begin{center}
\begin{tabularx}{\linewidth}{|X|X|X|}
\hline
\rowcolor{popblueL} \textbf{Français} & \textbf{English} & \textbf{Remarque} \\
\hline
pour + infinitif (même sujet) & to + BV & forme de base, la plus fréquente \\
\hline
afin de + infinitif & in order to / so as to + BV & plus formel, idéal à l'écrit \\
\hline
pour ne pas + infinitif & in order not to / so as not to + BV & jamais \eng{not to} seul \\
\hline
pour que / afin que + subjonctif & so that + sujet + can/could + BV & pas de subjonctif en anglais : on utilise un modal \\
\hline
\end{tabularx}
\end{center}

\begin{aretenir}[L'essentiel]
\begin{itemize}
\item Même sujet : \eng{to + BV} (ou \eng{in order to}, \eng{so as to}).
\item But négatif : \eng{in order not to} / \eng{so as not to}.
\item Sujets différents : \eng{so that + sujet + can/could/will/would + BV}.
\item Le français « pour que + subjonctif » devient \eng{so that + modal} : l'anglais n'a pas de subjonctif ici.
\end{itemize}
\end{aretenir}

\courssub{Exercice résolu et pièges à éviter}

\begin{exoresolu}[Choisir la bonne structure]
Énoncé : complétez chaque phrase avec la bonne expression du but : \eng{to}, \eng{in order not to}, \eng{so as to} ou \eng{so that}.

1. My mother cooks early \dots\ the children can eat before school.

2. I went to the market \dots\ buy some tomatoes.

3. He spoke quietly \dots\ disturb his sleeping brother.

4. We saved money \dots\ buy a new television.

5. The headmaster rang the bell \dots\ the students would know the break was over.

\tcblower
Solution détaillée :

1. \eng{so that} — les sujets sont différents (\emph{my mother} / \emph{the children}) : \eng{My mother cooks early so that the children can eat before school.}

2. \eng{to} — même sujet (\emph{I}) : \eng{I went to the market to buy some tomatoes.}

3. \eng{in order not to} (ou \eng{so as not to}) — même sujet, but négatif : \eng{He spoke quietly in order not to disturb his sleeping brother.}

4. \eng{so as to} (ou \eng{to} / \eng{in order to}) — même sujet (\emph{we}) : \eng{We saved money so as to buy a new television.}

5. \eng{so that} — sujets différents (\emph{the headmaster} / \emph{the students}) et contexte au passé, donc \eng{would} : \eng{The headmaster rang the bell so that the students would know the break was over.}
\end{exoresolu}

\begin{attention}[Erreurs fréquentes des francophones]
\textbf{1. Le faux « for to ».} Sous l'influence du français « pour », beaucoup d'élèves écrivent \eng{I went to the market for to buy tomatoes}. C'est \textbf{faux}. Dites \eng{to buy} ou \eng{in order to buy}.

\textbf{2. Quand « for » est correct.} \eng{for + nom} ou \eng{for + V-ing} exprime l'\emph{usage} d'un objet, pas le but d'une action : \eng{a knife for cutting bread} « un couteau pour couper le pain » ; \eng{a school for girls} « une école pour les filles ». Mais pour le but d'une action : \eng{I use a knife to cut bread}.

\textbf{3. « Not to » seul.} Ne dites pas \eng{She whispered not to wake the baby}. Dites \eng{so as not to wake} ou \eng{in order not to wake}.

\textbf{4. Oublier le sujet après « so that ».} \eng{so that} est toujours suivi d'une \emph{proposition complète} (sujet + verbe) : \eng{so that everyone can understand}, et non \eng{so that to understand}.

\textbf{5. Le temps du modal.} Au passé, utilisez \eng{could/would} : \eng{He hurried so that he could catch the bus}, et non \eng{so that he can catch}.
\end{attention}

\begin{exercice}[À vous de jouer]
Traduisez en anglais en utilisant la structure du but appropriée.

1. Je téléphone à mon frère pour lui souhaiter bon anniversaire.

2. Elle est partie tôt afin de ne pas rater le car pour Buea.

3. Le professeur parle fort pour que tous les élèves puissent l'entendre.

4. Nous économisons de l'argent afin d'acheter un ordinateur.

5. Il a fermé la porte doucement pour ne pas réveiller sa mère.

6. Mon père travaille le week-end pour que nous puissions aller à l'école.
\end{exercice}

\begin{corrige}[Exercice « À vous de jouer »]
1. \eng{I am phoning my brother to wish him a happy birthday.}

2. \eng{She left early so as not to miss the bus to Buea.} (ou \eng{in order not to miss})

3. \eng{The teacher speaks loudly so that all the students can hear him.}

4. \eng{We are saving money in order to buy a computer.} (ou \eng{to buy} / \eng{so as to buy})

5. \eng{He closed the door softly so as not to wake his mother.} (ou \eng{in order not to wake})

6. \eng{My father works at weekends so that we can go to school.}
\end{corrige}

\begin{savaistu}[Pourquoi « in order to » ?]
L'expression \eng{in order to} vient de l'idée de « mettre les choses en ordre » pour atteindre un objectif. À l'oral, les anglophones la raccourcissent souvent : on entend « in \emph{order} to » prononcé rapidement « in-or-deu ». Dans les examens camerounais, utiliser \eng{in order to} ou \eng{so as to} dans une rédaction montre au correcteur que vous maîtrisez un anglais soigné et formel : c'est une excellente stratégie pour valoriser votre copie !
\end{savaistu}

Vous savez désormais exprimer \emph{pourquoi} on agit. Mais une action a aussi des \emph{conséquences} : c'est ce que nous verrons dans la section suivante, consacrée à l'expression du résultat.

\coursec{Exprimer la conséquence (result)}

Dans la section précédente, nous avons appris à exprimer le \cle{but} (\eng{purpose}) : \eng{She works hard in order to pass the Bac.} Nous allons maintenant faire le chemin inverse : partir d'une cause et exprimer sa \cle{conséquence} (\eng{result}). Le but regarde vers l'avenir ; la conséquence regarde vers l'effet produit. Comparez :

\eng{He studied hard so as to succeed.} (Il a travaillé dur \textbf{pour} réussir — but.)

\eng{He studied hard, so he succeeded.} (Il a travaillé dur, \textbf{donc} il a réussi — conséquence.)

\courssub{Observation : la conséquence en contexte}

\begin{textebox}[Reading text : A Terrible Day in Douala]
Last Saturday, the Mbappé family planned a picnic at Limbe beach. But everything went wrong. It was raining so heavily that the bendskin drivers refused to leave Douala. The road was so muddy that even the taxis could not pass. Mrs Mbappé had prepared such delicious food that the children cried when they heard the picnic was cancelled. The rain was so violent that the roof of the kitchen started to leak. It was such a disappointment that little Junior refused to eat his lunch. Finally, the rain stopped at four o'clock, so the family decided to have the picnic in the sitting room. The children were so happy that they danced around the table. It was such a funny afternoon that they took many photos and sent them to their grandmother in Buea. As a result, the picnic was saved, and everybody agreed it was the best “indoor picnic” ever. Therefore, Mrs Mbappé decided that rain would never again stop their family projects.
\end{textebox}

\textbf{Glossaire :}

\begin{center}
\begin{tabularx}{\linewidth}{|X|X|X|}\hline
\rowcolor{popblueL} Mot anglais & Nature & Traduction \\ \hline
heavily & adverbe & fort, abondamment (« é-vi-li ») \\ \hline
muddy & adjectif & boueux (« me-di ») \\ \hline
to leak & verbe & fuir (toit) (« li-que ») \\ \hline
disappointment & nom & déception (« di-sa-poïnt-ment ») \\ \hline
sitting room & nom & salon (« si-ting roum ») \\ \hline
to save & verbe & sauver (« seïve ») \\ \hline
\end{tabularx}
\end{center}

\textbf{Questions de compréhension :}

\begin{enumerate}
\item Why did the bendskin drivers refuse to leave Douala?
\item What had Mrs Mbappé prepared for the picnic?
\item Where did the family finally have the picnic? Why?
\end{enumerate}

\textbf{Observation grammaticale.} Soulignez dans le texte toutes les structures avec \eng{so} et \eng{such}. Que remarquez-vous ? Après \eng{so}, on trouve un \textbf{adjectif} ou un \textbf{adverbe} (\eng{so heavily}, \eng{so muddy}, \eng{so happy}) ; après \eng{such}, on trouve un \textbf{nom} (\eng{such delicious food}, \eng{such a disappointment}). C'est toute la logique de la conséquence en anglais.

\courssub{So + proposition : la conséquence simple}

\begin{propriete}[So + proposition = « donc, alors, si bien que »]
\eng{So} placé entre deux propositions introduit la conséquence d'un fait. Il correspond au français « donc », « alors », « si bien que », « c'est pourquoi ».

\textbf{Forme :} proposition 1 (cause) + \cle{so} + proposition 2 (conséquence), souvent avec une virgule avant \eng{so}.
\end{propriete}

\begin{exemplebox}[So + proposition]
\eng{It was raining, so we stayed at home.} « Il pleuvait, donc nous sommes restés à la maison. »

\eng{The bus was full, so we took a bendskin.} « Le bus était plein, alors nous avons pris un bendskin. »

\eng{She missed the train, so she arrived late at school.} « Elle a raté le train, si bien qu'elle est arrivée en retard à l'école. »
\end{exemplebox}

\begin{attention}[Ne confondez pas les deux « so » !]
\eng{So} + proposition = conséquence (« donc »). Mais \eng{so that} + proposition = but (« pour que, afin que »), vu dans la section précédente.

\eng{He spoke slowly, so we understood.} (Il parlait lentement, \textbf{donc} nous avons compris — conséquence.)

\eng{He spoke slowly so that we could understand.} (Il parlait lentement \textbf{pour que} nous puissions comprendre — but.)
\end{attention}

\courssub{So + adjectif/adverbe + that : « si... que, tellement... que »}

\begin{propriete}[So + adjectif/adverbe + that]
Pour insister sur le degré d'une qualité et exprimer sa conséquence, l'anglais utilise :

\textbf{Forme :} \cle{so} + adjectif ou adverbe + \cle{that} + proposition.

Elle correspond au français « si... que » ou « tellement... que ». Dans la langue parlée, \eng{that} est souvent omis : \eng{The food was so hot (that) nobody could eat it.}
\end{propriete}

\begin{exemplebox}[So... that]
\eng{The food was so hot that nobody could eat it.} « La nourriture était si chaude que personne n'a pu la manger. »

\eng{The film was so boring that I slept.} « Le film était si ennuyeux que je me suis endormi. »

\eng{He ran so fast that nobody could follow him.} « Il courait si vite que personne ne pouvait le suivre. » (adverbe)

\eng{The exam was so difficult that half the class failed.} « L'examen était tellement difficile que la moitié de la classe a échoué. »
\end{exemplebox}

\courssub{Such (a/an) + (adjectif) + nom + that}

\begin{propriete}[Such + nom + that]
Quand l'idée principale est un \textbf{nom} (et non un adjectif), on utilise \cle{such} :

\textbf{Forme :} \cle{such} + (a/an) + (adjectif) + nom + \cle{that} + proposition.

\begin{itemize}
\item Nom \textbf{comptable singulier} : \eng{such a} + adjectif + nom — \eng{It was such a good meal that we thanked the cook twice.}
\item Nom \textbf{pluriel} ou \textbf{indénombrable} : \eng{such} + adjectif + nom (sans article) — \eng{They are such good students that the teachers love them.} ; \eng{It was such cold weather that we stayed inside.}
\end{itemize}
\end{propriete}

\begin{exemplebox}[Such... that]
\eng{She is such a kind woman that everybody loves her.} « C'est une femme si gentille que tout le monde l'aime. »

\eng{It was such a good meal that we thanked the cook twice.} « C'était un si bon repas que nous avons remercié le cuisinier deux fois. »

\eng{They made such a noise that the neighbours complained.} « Ils ont fait un tel bruit que les voisins se sont plaints. »

\eng{Yaoundé has such heavy traffic that I leave home at six.} « Yaoundé a une circulation si dense que je pars de la maison à six heures. » (indénombrable : pas d'article)
\end{exemplebox}

\begin{attention}[Le piège n°1 des francophones : l'ordre des mots]
En français on dit « une \textbf{si} gentille femme » : le mot du degré vient \emph{avant} l'adjectif. En anglais, on ne dit JAMAIS \eng{a so kind woman}. Les deux seules constructions correctes sont :

\begin{itemize}
\item \eng{such a kind woman} (construction normale) ;
\item \eng{so kind a woman} (construction littéraire, rare, à reconnaître mais pas à produire).
\end{itemize}

De même : \eng{such a tall boy} et non \eng{a such tall boy} ni \eng{a so tall boy}.
\end{attention}

\begin{popfigure}
\begin{center}
\begin{tikzpicture}[
box/.style={draw=popblue, thick, rounded corners, fill=popblueL, align=center, text width=5.6cm, inner sep=6pt},
exbox/.style={draw=poporange, thick, rounded corners, fill=poporangeL, align=center, text width=5.6cm, inner sep=6pt},
lab/.style={font=\bfseries, text=popdark}]
\node[lab] (t1) at (0,0) {so + ADJ / ADV + that};
\node[box, below=0.15cm of t1, align=center] (b1){The film was \textbf{so boring} that I slept.\\ He ran \textbf{so fast} that we lost him.};
\node[lab] (t2) at (8.6,0) {such (a) + ADJ + NOUN + that};
\node[exbox, below=0.15cm of t2, align=center] (b2){It was \textbf{such a boring film} that I slept.\\ She is \textbf{such a fast runner} that she wins every race.};
\draw[-{Stealth}, very thick, poppurple] (b1.east) -- (b2.west) node[midway, above, align=center, text=poppurple, font=\small\bfseries]{transformation\\ (Bac)};
\draw[-{Stealth}, very thick, poppurple] (b2.west) -- (b1.east);
\end{tikzpicture}
\end{center}
\legende{Les deux constructions de la conséquence : \eng{so} devant un adjectif/adverbe, \eng{such} devant un nom.}
\end{popfigure}

\courssub{Therefore et as a result : les connecteurs formels}

Pour l'expression écrite (composition, lettre formelle, essai au Bac), \eng{so} est un peu familier. On préfère des connecteurs plus soutenus, placés en tête de phrase et suivis d'une virgule.

\begin{definition}[Connecteurs formels de conséquence]
\begin{itemize}
\item \cle{therefore} = « par conséquent, donc » (très formel, se prononce « zé-fo ») ;
\item \cle{as a result} = « par conséquent, en conséquence » (« èze a ri-zeult ») ;
\item \cle{consequently} = « en conséquence » (formel, « kon-sé-kwen-tli ») ;
\item \cle{that is why} = « c'est pourquoi » (registre courant).
\end{itemize}
\end{definition}

\begin{exemplebox}[Connecteurs formels]
\eng{The bridge was destroyed by the floods. Therefore, the village was cut off for a week.} « Le pont a été détruit par les inondations. Par conséquent, le village a été isolé pendant une semaine. »

\eng{He did not revise his lessons. As a result, he failed the test.} « Il n'a pas révisé ses leçons. En conséquence, il a échoué au contrôle. »

\eng{Prices have increased; consequently, many families buy less meat.} « Les prix ont augmenté ; en conséquence, beaucoup de familles achètent moins de viande. »
\end{exemplebox}

\begin{attention}[Therefore n'est pas « donc » en début de phrase parlée]
\eng{Therefore} est un adverbe de liaison, pas une conjonction : on ne l'utilise pas pour relier deux propositions à l'intérieur d'une même phrase comme le « donc » français. Écrivez deux phrases : \eng{I was tired. Therefore, I went to bed early.} et non \eng{I was tired therefore I went to bed early} sans ponctuation forte.
\end{attention}

\courssub{Comparaison avec le français et tableau récapitulatif}

\begin{center}
\begin{tabularx}{\linewidth}{|X|X|X|}\hline
\rowcolor{popblueL} Français & English & Remarque \\ \hline
Il pleuvait, donc nous sommes restés. & It was raining, so we stayed. & \eng{so} + proposition \\ \hline
Le film était si ennuyeux que... & The film was so boring that... & \eng{so} + adjectif \\ \hline
Il courait si vite que... & He ran so fast that... & \eng{so} + adverbe \\ \hline
C'est une si gentille femme que... & She is such a kind woman that... & \eng{such a} + adjectif + nom singulier \\ \hline
Ce sont de si bons élèves que... & They are such good students that... & pluriel : pas d'article \\ \hline
C'était un temps si froid que... & It was such cold weather that... & indénombrable : pas d'article \\ \hline
Par conséquent, il a échoué. & Therefore, he failed. / As a result, he failed. & connecteurs formels \\ \hline
\end{tabularx}
\end{center}

\begin{aretenir}[L'essentiel de la conséquence]
\begin{enumerate}
\item \eng{So} + proposition = « donc » : \eng{I was ill, so I stayed at home.}
\item \eng{So} + adjectif/adverbe + \eng{that} = « si/tellement... que » : \eng{so tired that...}
\item \eng{Such (a/an)} + (adjectif) + nom + \eng{that} = « un tel... que » : \eng{such a kind woman that...}
\item Question-test : si le mot central est un \textbf{nom}, utilisez \eng{such} ; si c'est un \textbf{adjectif ou adverbe}, utilisez \eng{so}.
\item À l'écrit formel : \eng{therefore}, \eng{as a result}, \eng{consequently}.
\end{enumerate}
\end{aretenir}

\begin{savaistu}[Un grand classique du Bac]
Au baccalauréat camerounais, la transformation \eng{so... that} $\leftrightarrow$ \eng{such... that} est un exercice de réécriture très fréquent dans la partie « Grammar ». Exemple de consigne : \emph{Rewrite the sentence using “such”.} — \eng{The boy was so intelligent that he solved the problem.} devient \eng{He was such an intelligent boy that he solved the problem.} La méthode : repérez le nom (ici \eng{boy}), déplacez l'adjectif devant ce nom, ajoutez l'article \eng{a/an} si le nom est singulier comptable, et remplacez \eng{so} par \eng{such}. Entraînez-vous dans les deux sens !
\end{savaistu}

\begin{exoresolu}[Transformation so / such]
Énoncé : réécrivez la phrase suivante avec \eng{such} : \eng{The match was so exciting that the crowd shouted for an hour.}
\tcblower
Solution détaillée :
\begin{enumerate}
\item On identifie le nom à créer : la qualité \eng{exciting} porte sur \eng{match} (nom comptable singulier).
\item On construit le groupe nominal : \eng{an exciting match} (article \eng{an} devant voyelle).
\item On place \eng{such} devant ce groupe : \eng{such an exciting match}.
\item On adapte le début de la phrase : \eng{It was such an exciting match that the crowd shouted for an hour.}
\end{enumerate}
Traduction : « C'était un match si passionnant que la foule a crié pendant une heure. »
\end{exoresolu}

\begin{exercice}[À vous de jouer]
\begin{enumerate}
\item Complétez avec \eng{so} ou \eng{such (a/an)} : \\
a) The soup was \ldots\ hot that I could not eat it. \\
b) It was \ldots\ a long journey that we slept in the car. \\
c) The students made \ldots\ a noise that the teacher got angry. \\
d) She speaks \ldots\ quickly that I cannot follow her. \\
e) They are \ldots\ good footballers that everybody wants them in the team.
\item Réécrivez avec \eng{such} : \eng{The house was so big that we got lost inside.}
\item Réécrivez avec \eng{so} : \eng{It was such a boring lesson that two students slept.}
\item Traduisez : « Il faisait tellement chaud que nous sommes allés à la piscine. Par conséquent, nous avons raté le match. »
\end{enumerate}
\end{exercice}

\begin{corrige}[Exercice « À vous de jouer »]
\begin{enumerate}
\item a) \eng{so} (adjectif \eng{hot}) ; b) \eng{such} (nom singulier \eng{journey}) ; c) \eng{such} (nom \eng{noise}) ; d) \eng{so} (adverbe \eng{quickly}) ; e) \eng{such} (nom pluriel \eng{footballers}, sans article).
\item \eng{It was such a big house that we got lost inside.}
\item \eng{The lesson was so boring that two students slept.}
\item \eng{It was so hot that we went to the swimming pool. As a result, we missed the match.}
\end{enumerate}
\end{corrige}

\coursec{Obligation, interdiction et préférence}

Dans la section précédente, nous avons appris à exprimer la conséquence avec \eng{so}, \eng{so... that} et \eng{as a result}. Mais la vie familiale et sociale, c'est aussi un ensemble de \cle{règles} : ce qu'il faut faire, ce qu'il est interdit de faire, et ce que l'on préfère faire. Un règlement scolaire, un père qui donne des instructions à ses enfants, un ami qui propose une sortie : dans tous ces cas, l'anglais utilise des structures précises qu'il faut maîtriser sans faute. Observons d'abord ces structures en contexte.

\courssub{Observation : les structures en contexte}

\begin{textebox}[School rules and family choices — a dialogue]
\textbf{Mrs Etoundi} (a teacher at GBHS Yaoundé): Good morning, students. Before the new term starts, let me remind you of the school rules. You \textbf{must wear} the school uniform every day, and you \textbf{mustn't use} your phones during lessons. These are the headmaster's rules, so you \textbf{have to respect} them.

\textbf{Ngono}: Excuse me, Madam. Do we \textbf{have to come} to school on Saturdays?

\textbf{Mrs Etoundi}: No, you \textbf{don't have to come} on Saturdays. There are no classes. But last year, students \textbf{had to attend} extra lessons on Saturday mornings because of the exams. Next term, you \textbf{will have to work} very hard.

\textbf{Ali}: Madam, I \textbf{prefer} studying in the morning \textbf{to} studying in the evening. Can we choose our timetable?

\textbf{Mrs Etoundi}: I'm afraid not, Ali. Personally, I \textbf{would rather teach} in the morning too, but the timetable is fixed.

\textbf{Ngono}: I' \textbf{d rather not have} classes on Friday afternoons!

\textbf{Mrs Etoundi}: Well, Ngono, I understand you, but rules are rules!
\end{textebox}

\textbf{Glossaire :}

\begin{center}
\begin{tabularx}{\linewidth}{|l|l|X|}
\hline
\rowcolor{popblueL} \textbf{English} & \textbf{Nature} & \textbf{Français} \\
\hline
a rule & nom & une règle, un règlement \\
\hline
to remind & verbe & rappeler \\
\hline
a headmaster & nom & un proviseur, un directeur \\
\hline
to attend & verbe & assister à, suivre (des cours) \\
\hline
extra lessons & nom & des cours supplémentaires \\
\hline
a timetable & nom & un emploi du temps \\
\hline
fixed & adjectif & fixe, immuable \\
\hline
\end{tabularx}
\end{center}

\textbf{Questions de compréhension :}
\begin{enumerate}
\item What two rules does Mrs Etoundi mention?
\item Do students have to come to school on Saturdays this year?
\item Why did students have to attend extra lessons last year?
\item What does Ali prefer?
\end{enumerate}

\textbf{Observons.} Dans ce dialogue, on trouve \eng{must wear}, \eng{have to respect}, \eng{mustn't use}, \eng{don't have to come}, \eng{had to attend}, \eng{will have to work}, puis \eng{prefer... to} et \eng{would rather teach}. Toutes ces structures expriment l'\cle{obligation}, l'\cle{interdiction}, l'\cle{absence d'obligation} ou la \cle{préférence}. Étudions-les une par une.

\courssub{Must et have to : deux types d'obligation}

\begin{propriete}[Must + base verbale : l'obligation ressentie]
\eng{Must} est un \cle{auxiliaire modal}. Il est suivi de la \textbf{base verbale} (l'infinitif \textbf{sans} \eng{to}), il est \textbf{invariable} (jamais de \eng{-s} à la 3\up{e} personne) et il exprime une obligation \textbf{ressentie par le locuteur lui-même} : c'est lui qui juge que c'est nécessaire.

\eng{You must listen to your parents.} « Tu dois écouter tes parents. » (C'est moi qui te le dis, c'est mon avis.)

\eng{I must phone my grandmother today.} « Il faut absolument que je téléphone à ma grand-mère aujourd'hui. »
\end{propriete}

\begin{propriete}[Have to + base verbale : l'obligation extérieure]
\eng{Have to} + base verbale exprime une obligation qui vient de \textbf{l'extérieur} : une loi, un règlement, une circonstance. Ce n'est pas forcément l'avis du locuteur. Attention : \eng{have to} se conjugue comme un verbe ordinaire (\eng{has to} à la 3\up{e} personne, auxiliaire \eng{do/does} à la négation et à l'interrogation).

\eng{Students have to wear uniforms at GBHS Buea.} « Les élèves doivent porter l'uniforme. » (C'est le règlement de l'école.)

\eng{In Cameroon, drivers have to drive on the right.} « Au Cameroun, les conducteurs doivent rouler à droite. » (C'est la loi.)

\eng{Does she have to work on Sundays?} « Est-ce qu'elle est obligée de travailler le dimanche ? »
\end{propriete}

\begin{aretenir}[Must ou have to ?]
La différence est souvent subtile et, dans la conversation, les deux sont fréquemment interchangeables. Mais pour un examen, retenez :
\begin{itemize}
\item \eng{must} = obligation personnelle, ressentie par celui qui parle (\eng{I must go now}, « je dois partir », c'est ma décision) ;
\item \eng{have to} = obligation imposée par une règle, une loi, une situation extérieure (\eng{I have to wear a uniform}, « le règlement m'y oblige »).
\end{itemize}
\end{aretenir}

\courssub{Mustn't et don't have to : le piège capital}

Voici le point le plus dangereux de la leçon pour un francophone. En français, « ne pas devoir » peut vouloir dire deux choses très différentes ; en anglais, les deux sens s'expriment avec des structures \textbf{différentes}.

\begin{propriete}[Mustn't et don't have to : deux sens opposés]
\begin{itemize}
\item \eng{Mustn't} + base verbale = \textbf{interdiction} : « il est interdit de », « il ne faut surtout pas ».
\item \eng{Don't/Doesn't have to} + base verbale = \textbf{absence d'obligation} : « ce n'est pas la peine de », « tu n'es pas obligé de ».
\end{itemize}

\eng{You mustn't smoke here.} « Il est interdit de fumer ici. » (Si tu fumes, tu enfreins la règle.)

\eng{You don't have to come.} « Tu n'es pas obligé de venir. » (Tu peux venir si tu veux, mais ce n'est pas nécessaire.)
\end{propriete}

\begin{exemplebox}[Exemples traduits]
\begin{itemize}
\item \eng{Students must wear uniforms.} « Les élèves doivent porter l'uniforme. »
\item \eng{You mustn't smoke here.} « Il est interdit de fumer ici. »
\item \eng{You don't have to come.} « Tu n'es pas obligé de venir. »
\item \eng{We mustn't be late for the family meeting.} « Nous ne devons surtout pas être en retard à la réunion de famille. »
\item \eng{She doesn't have to cook tonight; her husband is doing it.} « Elle n'est pas obligée de cuisiner ce soir ; c'est son mari qui le fait. »
\item \eng{Passengers mustn't lean out of the bendskin.} « Les passagers ne doivent pas se pencher hors de la moto-taxi. »
\end{itemize}
\end{exemplebox}

\begin{attention}[Erreurs fréquentes des francophones]
\begin{itemize}
\item \eng{You must to go.} est \textbf{FAUX} : après \eng{must}, jamais de \eng{to}. Dites \eng{You must go.}
\item \eng{He musts work.} est \textbf{FAUX} : \eng{must} est invariable. Dites \eng{He must work.}
\item Ne traduisez jamais « tu ne dois pas venir » (au sens de « ce n'est pas la peine ») par \eng{you mustn't come} : cela signifie « il t'est interdit de venir » ! Dites \eng{You don't have to come.}
\item \eng{Mustn't} se prononce « meu-ssent » : le premier \eng{t} ne se prononce pas.
\end{itemize}
\end{attention}

\courssub{Passé et futur : had to et will have to}

\begin{propriete}[Must n'a ni passé ni futur]
Le modal \eng{must} ne possède \textbf{aucune forme de passé ni de futur}. Pour exprimer l'obligation à ces temps, on le remplace obligatoirement par \eng{have to} conjugué :
\begin{itemize}
\item passé : \eng{had to} + base verbale ;
\item futur : \eng{will have to} + base verbale.
\end{itemize}

\eng{Last year, we had to attend extra lessons.} « L'année dernière, nous avons dû suivre des cours supplémentaires. »

\eng{Next term, you will have to work harder.} « Le trimestre prochain, vous devrez travailler davantage. »

\eng{Yesterday my father had to go to Douala for a family meeting.} « Hier, mon père a dû aller à Douala pour une réunion de famille. »
\end{propriete}

\begin{attention}
\eng{I musted go} et \eng{I will must go} sont \textbf{FAUX} et n'existent pas en anglais. Au passé : \eng{I had to go.} Au futur : \eng{I will have to go.}
\end{attention}

\begin{popfigure}
\begin{center}
\begin{tikzpicture}[scale=1, every node/.style={font=\small}]
\node[draw=popblue, fill=popblueL, rounded corners, align=center, minimum width=3.4cm, minimum height=1cm] (pres) at (0,2.2) {\textbf{Présent}\\ must / have to};
\node[draw=popgreen, fill=popgreenL, rounded corners, align=center, minimum width=3.4cm, minimum height=1cm] (past) at (-3.6,0) {\textbf{Passé}\\ had to};
\node[draw=poporange, fill=poporangeL, rounded corners, align=center, minimum width=3.4cm, minimum height=1cm] (fut) at (3.6,0) {\textbf{Futur}\\ will have to};
\draw[-{Stealth}, thick, popgreen] (pres.south west) -- (past.north) node[midway, above left, popdark]{pas de passé de \eng{must}};
\draw[-{Stealth}, thick, poporange] (pres.south east) -- (fut.north) node[midway, above right, popdark]{pas de futur de \eng{must}};
\node[draw=poppurple, fill=poppinkL, rounded corners, align=center, minimum width=4.2cm, minimum height=1cm] (inter) at (0,-1.6) {\textbf{Interdiction :} mustn't\\ \textbf{Pas obligé :} don't have to};
\draw[-{Stealth}, thick, poppurple] (pres.south) -- (inter.north);
\end{tikzpicture}
\end{center}
\legende{Carte mentale : les formes de l'obligation et de l'interdiction en anglais.}
\end{popfigure}

\begin{center}
\begin{tabularx}{\linewidth}{|l|X|X|}
\hline
\rowcolor{popblueL} \textbf{Structure} & \textbf{Sens} & \textbf{Exemple et traduction} \\
\hline
must + BV & obligation ressentie & \eng{I must help my mother.} « Je dois aider ma mère. » \\
\hline
have to + BV & règle extérieure & \eng{We have to wear uniforms.} « Nous devons porter l'uniforme (règlement). » \\
\hline
mustn't + BV & interdiction & \eng{You mustn't smoke here.} « Il est interdit de fumer ici. » \\
\hline
don't have to + BV & absence d'obligation & \eng{You don't have to come.} « Tu n'es pas obligé de venir. » \\
\hline
had to + BV & obligation au passé & \eng{She had to stay home.} « Elle a dû rester à la maison. » \\
\hline
will have to + BV & obligation au futur & \eng{They will have to pay.} « Ils devront payer. » \\
\hline
\end{tabularx}
\end{center}

\courssub{Exprimer la préférence : prefer et would rather}

\begin{propriete}[Prefer + nom ou V-ing ... to ...]
Pour exprimer une préférence \textbf{générale, habituelle}, on utilise \eng{prefer} suivi d'un \textbf{nom} ou d'un \textbf{verbe en -ing}. L'élément rejeté est introduit par la préposition \eng{to} (jamais \eng{than} !).

\eng{I prefer fufu to rice.} « Je préfère le fufu au riz. »

\eng{I prefer walking to running.} « Je préfère marcher plutôt que courir. »

\eng{She prefers tea to coffee.} « Elle préfère le thé au café. »

\eng{They prefer living in Bamenda to living in Douala.} « Ils préfèrent vivre à Bamenda plutôt qu'à Douala. »
\end{propriete}

\begin{propriete}[Would rather + base verbale (than + base verbale)]
Pour exprimer une préférence \textbf{ponctuelle, dans une situation précise}, on utilise \eng{would rather} (souvent contracté en \eng{'d rather}) suivi de la \textbf{base verbale sans to}. L'alternative est introduite par \eng{than} + base verbale.

\eng{I would rather stay home than go out.} « Je préférerais rester à la maison plutôt que de sortir. »

\eng{I'd rather drink water than wine.} « Je préférerais boire de l'eau plutôt que du vin. »

\textbf{Forme négative :} \eng{would rather not} + BV. \quad \eng{I'd rather not talk about it.} « Je préférerais ne pas en parler. »

\textbf{Forme interrogative :} \eng{Would you rather...?} \quad \eng{Would you rather eat now?} « Tu préférerais manger maintenant ? »

\textbf{Réponse courte :} \eng{I'd rather not.} « Je préférerais ne pas (le faire). »
\end{propriete}

\begin{attention}[Les trois pièges de la préférence]
\begin{itemize}
\item \eng{I prefer rice than fufu.} est \textbf{FAUX} : avec \eng{prefer}, la préposition est \eng{to}. Dites \eng{I prefer rice to fufu.}
\item \eng{I'd rather to stay.} est \textbf{FAUX} : après \eng{would rather}, la base verbale \textbf{sans to}. Dites \eng{I'd rather stay.}
\item \eng{I prefer to walk than to run} est lourd et souvent fautif à ce niveau ; préférez \eng{I prefer walking to running.}
\item Ne confondez pas \eng{would rather} (préférence) et \eng{would like} (désir) : \eng{I'd like some water} = « je voudrais de l'eau » ; \eng{I'd rather drink water} = « je préférerais boire de l'eau ».
\end{itemize}
\end{attention}

\begin{center}
\begin{tabularx}{\linewidth}{|l|X|X|}
\hline
\rowcolor{popblueL} \textbf{Structure} & \textbf{Construction} & \textbf{Exemple et traduction} \\
\hline
prefer (général) & prefer + nom / V-ing + \textbf{to} & \eng{I prefer fufu to rice.} « Je préfère le fufu au riz. » \\
\hline
would rather (ponctuel) & would rather + BV + \textbf{than} + BV & \eng{I'd rather stay than go.} « Je préférerais rester plutôt que partir. » \\
\hline
négation & would rather \textbf{not} + BV & \eng{I'd rather not go out.} « Je préférerais ne pas sortir. » \\
\hline
question & \textbf{Would} + sujet + \textbf{rather} + BV ? & \eng{Would you rather eat now?} « Tu préférerais manger maintenant ? » \\
\hline
\end{tabularx}
\end{center}

\courssub{Entraînement et production}

\begin{exoresolu}[Transformez selon le modèle]
\textbf{Énoncé.} Réécrivez chaque phrase avec l'expression entre parenthèses, sans changer le sens.

1. It is forbidden to park here. (\eng{mustn't})

2. It is not necessary for you to bring food; we have everything. (\eng{don't have to})

3. Last Saturday, I was obliged to visit my uncle in Bafoussam. (\eng{had to})

4. I like football more than volleyball. (\eng{prefer})

5. I want to stay home instead of going to the market. (\eng{would rather})

\tcblower
\textbf{Solution détaillée.}

1. \eng{You mustn't park here.} — L'interdiction se rend par \eng{mustn't} + base verbale (\eng{park}), sans \eng{to}.

2. \eng{You don't have to bring food; we have everything.} — L'absence d'obligation se rend par \eng{don't have to} + BV. Attention : \eng{you mustn't bring food} signifierait « il est interdit d'apporter de la nourriture », ce qui est le contraire du sens voulu !

3. \eng{Last Saturday, I had to visit my uncle in Bafoussam.} — \eng{Must} n'a pas de passé ; on utilise \eng{had to} + BV.

4. \eng{I prefer football to volleyball.} — Avec \eng{prefer} + nom, l'élément rejeté est introduit par \eng{to}, jamais par \eng{than}.

5. \eng{I would rather stay home than go to the market.} — \eng{Would rather} + base verbale (\eng{stay}), puis \eng{than} + base verbale (\eng{go}), sans \eng{to}.
\end{exoresolu}

\begin{exercice}[À vous de jouer]
\begin{enumerate}
\item Complétez avec \eng{must}, \eng{mustn't}, \eng{have to} ou \eng{don't have to} :
\begin{enumerate}
\item You \ldots\ cross the road when the light is red. It's dangerous.
\item Tomorrow is a holiday, so we \ldots\ go to school.
\item I \ldots\ finish this letter tonight; I promised my sister.
\item Students \ldots\ respect the teachers; it's the first school rule.
\end{enumerate}
\item Mettez les phrases au passé : a) \eng{She must sell her tomatoes early.} b) \eng{They have to travel to Yaoundé.}
\item Traduisez en anglais : a) « Je préfère le ndolé au poulet. » b) « Tu préférerais partir maintenant ? » c) « Je préférerais ne pas répondre. » d) « Il est interdit de jeter des ordures ici. »
\end{enumerate}
\end{exercice}

\begin{corrige}[Exercice « À vous de jouer »]
\textbf{1.} a) \eng{mustn't} (interdiction, danger) ; b) \eng{don't have to} (absence d'obligation : jour férié) ; c) \eng{must} (obligation ressentie par le locuteur : une promesse) ; d) \eng{have to} (règle extérieure : le règlement scolaire).

\textbf{2.} a) \eng{She had to sell her tomatoes early.} b) \eng{They had to travel to Yaoundé.} (\eng{must} et \eng{have to} deviennent tous deux \eng{had to} au passé.)

\textbf{3.} a) \eng{I prefer ndolé to chicken.} b) \eng{Would you rather leave now?} c) \eng{I'd rather not answer.} d) \eng{You mustn't throw litter here.}
\end{corrige}

\begin{aretenir}[L'essentiel de la section]
\begin{itemize}
\item \eng{must} + BV = obligation ressentie ; \eng{have to} + BV = règle extérieure.
\item \eng{mustn't} = interdiction ; \eng{don't have to} = pas obligé. Ne jamais les confondre !
\item Passé : \eng{had to} ; futur : \eng{will have to} (\eng{must} n'a ni passé ni futur).
\item Préférence générale : \eng{prefer} + nom/V-ing + \eng{to} ; préférence ponctuelle : \eng{would rather} + BV (+ \eng{than} + BV), négation \eng{would rather not}.
\end{itemize}
\end{aretenir}

Nous savons désormais exprimer ce qui est obligatoire, interdit ou simplement préféré. Dans la section suivante, nous verrons comment proposer une activité en groupe avec \eng{let's} et \eng{let's not} suivis de la base verbale.

\coursec{Let's / Let's not + base verbale}

Dans la section précédente, nous avons appris à exprimer l'obligation (\eng{must}, \eng{have to}), l'interdiction (\eng{mustn't}) et la préférence (\eng{would rather}, \eng{prefer}). Toutes ces structures servent à dire ce que quelqu'un \emph{doit} faire ou \emph{préfère} faire. Mais comment faire lorsque l'on veut \emph{proposer} une activité à faire \emph{ensemble}, en incluant soi-même dans le groupe ? En français, nous utilisons l'impératif à la première personne du pluriel : « Allons-y ! », « Ne perdons pas de temps ! ». En anglais, il existe une structure très simple et très fréquente pour cela : \eng{let's} suivi de la base verbale. C'est l'objet de cette section.

\courssub{Observation : découvrir la structure en contexte}

Lisons d'abord ce court dialogue entre deux camarades de classe, Amina et Peter, un samedi matin à Yaoundé. Soulignez mentalement toutes les phrases qui contiennent \eng{let's}.

\begin{textebox}[A Saturday plan — dialogue]
\textbf{Amina :} Hi Peter! The weather is lovely today. \textbf{Let's} do something together this afternoon.

\textbf{Peter :} Good idea! \textbf{Let's} go to the Mvog-Betsi zoo. I have never been there.

\textbf{Amina :} Hmm, it's quite far. \textbf{Let's not} take a taxi, it's too expensive. \textbf{Let's} take a bendskin instead.

\textbf{Peter :} OK. And \textbf{let's not} forget our student cards — the entrance is cheaper with them.

\textbf{Amina :} Right. \textbf{Let's} meet at the post office at two o'clock, shall we?

\textbf{Peter :} Yes, \textbf{let's}! Oh, and \textbf{let me} carry the bag with the food. You carried it last time.

\textbf{Amina :} That's kind of you. And \textbf{let's not} be late this time!
\end{textebox}

\begin{definition}[Glossaire du dialogue]
\begin{itemize}
\item \eng{bendskin} (n.) : motocyclette-taxi (mot courant au Cameroun anglophone).
\item \eng{student card} (n.) : carte d'étudiant.
\item \eng{entrance} (n.) : entrée (prix d'entrée).
\item \eng{to carry} (v.) : porter.
\item \eng{kind} (adj.) : gentil, aimable.
\end{itemize}
\end{definition}

\textbf{Questions d'observation :}
\begin{enumerate}
\item Combien de propositions d'actions communes Amina et Peter font-ils ? Par quelle structure commencent-elles ?
\item Que remarquez-vous après \eng{let's} : y a-t-il \eng{to} ? Y a-t-il un verbe en \eng{-ing} ?
\item Comment les deux amis expriment-ils une proposition \emph{négative} ?
\item Dans la phrase de Peter \eng{let me carry the bag}, s'agit-il d'une proposition collective ? Pourquoi ?
\end{enumerate}

\courssub{La règle : let's + base verbale}

\begin{propriete}[Règle — let's / let's not + base verbale]
\begin{itemize}
\item \eng{Let's} est la contraction de \eng{let us}. On l'emploie pour \cle{proposer une action qui inclut le locuteur} et au moins une autre personne.
\item \textbf{Forme affirmative :} \eng{let's + base verbale} (l'infinitif \emph{sans} \eng{to}, jamais de \eng{-ing}).\\
\eng{Let's play football.} = Jouons au football.
\item \textbf{Forme négative :} \eng{let's not + base verbale} (le \eng{not} se place \emph{avant} le verbe, sans auxiliaire \eng{do}).\\
\eng{Let's not be late.} = Ne soyons pas en retard.
\item Le verbe reste invariable : pas de \eng{-s}, pas de temps marqué.
\item Traduction française habituelle : l'impératif à la première personne du pluriel (« -ons ») : \eng{Let's go!} = « Allons-y ! »
\end{itemize}
\end{propriete}

\begin{exemplebox}[Exemples traduits]
\begin{itemize}
\item \eng{Let's visit our grandparents this weekend.} = Visitons nos grands-parents ce week-end.
\item \eng{Let's not argue about money.} = Ne nous disputons pas pour l'argent.
\item \eng{Let's help our mother in the kitchen.} = Aidons notre mère dans la cuisine.
\item \eng{Let's speak English in class.} = Parlons anglais en classe.
\item \eng{Let's not waste water.} = Ne gaspillons pas l'eau.
\item \eng{Let's revise our lessons before the test.} = Révisons nos leçons avant le contrôle.
\end{itemize}
\end{exemplebox}

\begin{attention}[Erreurs fréquentes des francophones]
\begin{itemize}
\item \eng{Let's to go.} est \textbf{faux} : jamais de \eng{to} après \eng{let's}. Correct : \eng{Let's go.}
\item \eng{Let's going.} est \textbf{faux} : jamais de \eng{-ing} après \eng{let's}. Correct : \eng{Let's go.}
\item \eng{Let's don't argue.} est \textbf{faux} : la négation se fait avec \eng{not} seul, sans \eng{do}. Correct : \eng{Let's not argue.}
\item Ne traduisez pas mot à mot : \eng{let's} ne signifie pas « laissons ». \eng{Let's eat} = « Mangeons », pas « Laissons manger ».
\end{itemize}
\end{attention}

\courssub{Réagir à une proposition}

Quand quelqu'un vous fait une proposition avec \eng{let's}, voici les réponses naturelles :

\begin{center}
\begin{tabularx}{\linewidth}{|X|X|X|}\hline
\rowcolor{popblueL} English & Français & Emploi \\ \hline
\eng{Yes, let's!} & Oui, allons-y ! & acceptation enthousiaste \\ \hline
\eng{That's a good idea!} & C'est une bonne idée ! & acceptation \\ \hline
\eng{OK, why not?} & D'accord, pourquoi pas ? & acceptation modérée \\ \hline
\eng{I'd rather not.} & Je préférerais ne pas. & refus poli \\ \hline
\eng{No, let's not.} & Non, ne le faisons pas. & refus direct \\ \hline
\end{tabularx}
\end{center}

\begin{exemplebox}[Mini-dialogue]
\eng{— Let's watch the football match tonight.}\\
\eng{— That's a good idea! Let's watch it at my house.}\\
« — Regardons le match de foot ce soir. — Bonne idée ! Regardons-le chez moi. »
\end{exemplebox}

\courssub{Let's et la question tag : shall we ?}

\begin{propriete}[La question tag de let's]
Une phrase commençant par \eng{let's} prend la question tag \eng{shall we ?} :\\
\eng{Let's go, shall we?} = Allons-y, d'accord ? / On y va ?\\
\eng{Let's not tell anyone, shall we?} = Ne le disons à personne, d'accord ?\\
C'est la seule question tag possible après \eng{let's} : jamais \eng{do we} ni \eng{will we}.
\end{propriete}

\begin{savaistu}[Un classique des exercices de tags]
\eng{Let's go, shall we?} est un grand classique des exercices de question tags au lycée, car c'est un cas \emph{irrégulier} : alors que la plupart des tags reprennent l'auxiliaire de la phrase (\eng{You are tired, aren't you?}), \eng{let's} impose toujours \eng{shall we}. Retenez aussi que \eng{shall we + base verbale ?} peut à lui seul exprimer une proposition : \eng{Shall we start?} = « On commence ? ». La différence est subtile : \eng{shall we...?} \emph{demande} l'avis de l'autre, tandis que \eng{let's...} \emph{affirme} la proposition.
\end{savaistu}

\courssub{Attention au piège : let's ≠ let + objet}

\begin{attention}[Let's (proposition) ≠ let + me/him/them (permission)]
Ne confondez pas les deux structures :
\begin{itemize}
\item \eng{Let's + BV} : proposition \emph{collective} incluant le locuteur. \eng{Let's dance!} = Dansons ! (toi et moi)
\item \eng{Let + complément d'objet + BV} : \emph{permission} ou suggestion concernant quelqu'un d'autre. \eng{Let me help you.} = Laisse-moi t'aider. \eng{Let him speak.} = Laisse-le parler. \eng{Let them play.} = Laisse-les jouer.
\end{itemize}
Dans les deux cas, le verbe reste à la base verbale, mais le sens est totalement différent : \eng{Let's go home} (rentrons ensemble) ≠ \eng{Let us go home} (laissez-nous rentrer — une demande de permission adressée à quelqu'un d'autre !).
\end{attention}

\begin{popfigure}
\begin{center}
\begin{tikzpicture}[node distance=0.6cm, every node/.style={font=\small}]
\node[draw=popblue, fill=popblueL, rounded corners, align=center, minimum width=4.2cm, minimum height=1.1cm] (lets) {\textbf{LET'S + BV}\\ \eng{Let's go!} = Allons-y !};
\node[draw=poporange, fill=poporangeL, rounded corners, align=center, minimum width=4.2cm, minimum height=1.1cm, right=1.2cm of lets] (letsnot) {\textbf{LET'S NOT + BV}\\ \eng{Let's not go.} = N'y allons pas.};
\node[draw=popgreen, fill=popgreenL, rounded corners, align=center, minimum width=4.6cm, minimum height=1.1cm, below=1.1cm of $(lets)!0.5!(letsnot)$] (letobj) {\textbf{LET + me/him/them + BV}\\ \eng{Let me help you.} = Laisse-moi t'aider.};
\node[below=0.15cm of lets, font=\footnotesize\itshape, text=popdark] {proposition collective};
\node[below=0.15cm of letsnot, font=\footnotesize\itshape, text=popdark] {proposition négative};
\node[below=0.15cm of letobj, font=\footnotesize\itshape, text=popdark] {permission / suggestion individuelle};
\draw[-{Stealth}, thick, popblue] (lets.south) -- (letobj.north west);
\draw[-{Stealth}, thick, poporange] (letsnot.south) -- (letobj.north east);
\end{tikzpicture}
\end{center}
\legende{Trois structures avec \eng{let} : deux propositions collectives et une permission individuelle.}
\end{popfigure}

\courssub{Comparaison avec le français}

\begin{center}
\begin{tabularx}{\linewidth}{|X|X|X|}\hline
\rowcolor{popblueL} Français & English & Remarque \\ \hline
Allons-y ! & \eng{Let's go!} & impératif en « -ons » = \eng{let's + BV} \\ \hline
Ne perdons pas de temps. & \eng{Let's not waste time.} & négation : \eng{not} avant le verbe \\ \hline
On commence ? & \eng{Shall we start?} & proposition sous forme de question \\ \hline
Laisse-moi t'aider. & \eng{Let me help you.} & permission individuelle, pas de proposition \\ \hline
Mangeons, d'accord ? & \eng{Let's eat, shall we?} & tag obligatoire : \eng{shall we} \\ \hline
\end{tabularx}
\end{center}

\courssub{Exercice résolu et entraînement}

\begin{exoresolu}[Complétez avec la bonne forme]
Énoncé — Complétez chaque phrase avec la forme correcte du verbe entre parenthèses, puis traduisez :
\begin{enumerate}
\item Let's \dots\dots\ (visit) the museum in Foumban.
\item Let's not \dots\dots\ (be) late for school.
\item Let me \dots\dots\ (carry) your bag.
\item Let's go to the market, \dots\dots\ we?
\end{enumerate}
\tcblower
Solution détaillée :
\begin{enumerate}
\item \eng{Let's visit the museum in Foumban.} — Après \eng{let's}, on met la base verbale, sans \eng{to} ni \eng{-ing}. « Visitons le musée de Foumban. »
\item \eng{Let's not be late for school.} — La négation se construit avec \eng{not} placé avant la base verbale. « Ne soyons pas en retard à l'école. »
\item \eng{Let me carry your bag.} — Après \eng{let + complément d'objet}, on met aussi la base verbale ; ici c'est une permission individuelle, pas une proposition. « Laisse-moi porter ton sac. »
\item \eng{Let's go to the market, shall we?} — La question tag de \eng{let's} est toujours \eng{shall we}. « Allons au marché, d'accord ? »
\end{enumerate}
\end{exoresolu}

\begin{exercice}[À vous de jouer]
\begin{enumerate}
\item Traduisez en anglais : a) « Jouons au football ! » b) « Ne nous disputons pas. » c) « Aidons nos parents. » d) « Rentrons à la maison, d'accord ? »
\item Corrigez les erreurs : a) \eng{Let's to watch TV.} b) \eng{Let's going home.} c) \eng{Let's don't be sad.}
\item Réagissez à chaque proposition par une acceptation ou un refus poli : a) \eng{Let's have a break!} b) \eng{Let's not do our homework.}
\end{enumerate}
\end{exercice}

\begin{corrige}[Exercice « À vous de jouer »]
\begin{enumerate}
\item a) \eng{Let's play football!} b) \eng{Let's not argue.} c) \eng{Let's help our parents.} d) \eng{Let's go home, shall we?}
\item a) \eng{Let's watch TV.} (pas de \eng{to}) b) \eng{Let's go home.} (pas de \eng{-ing}) c) \eng{Let's not be sad.} (négation sans \eng{do})
\item Réponses possibles : a) \eng{Yes, let's! / That's a good idea!} b) \eng{I'd rather not. We must do our homework.}
\end{enumerate}
\end{corrige}

\begin{aretenir}[L'essentiel]
\begin{itemize}
\item \eng{Let's + base verbale} = proposition collective (« -ons » en français) ; négation : \eng{let's not + BV}.
\item Jamais de \eng{to}, jamais de \eng{-ing}, jamais de \eng{do} après \eng{let's}.
\item Question tag : \eng{Let's go, shall we?}
\item \eng{Let me/him/them + BV} = permission individuelle, à ne pas confondre avec \eng{let's}.
\end{itemize}
\end{aretenir}

\coursec{Present perfect et past perfect}

Dans la section précédente, nous avons appris à formuler des propositions et des suggestions avec \eng{let's} et \eng{let's not} suivis de la base verbale : \eng{Let's visit grandmother this weekend!} Nous allons maintenant nous intéresser à deux temps essentiels pour raconter la vie de famille et les événements sociaux : le \cle{present perfect} et le \cle{past perfect}. Ces deux temps permettent de situer les actions dans le temps avec une précision que le français exprime différemment — et c'est là que se cachent les erreurs les plus fréquentes des francophones.

\courssub{Observation : découvrir les deux temps en contexte}

Lisons d'abord ce court texte sur une réunion de famille à Yaoundé. Repérez les verbes conjugués avec \eng{have/has} et ceux conjugués avec \eng{had}.

\begin{textebox}[A family reunion in Yaoundé]
The Mbarga family has always been very close. Every December, relatives travel from Douala, Bafoussam and even from abroad to spend a few days together in Yaoundé. This year, the reunion has been particularly special because Grandmother Ekotto has just celebrated her eightieth birthday.

Clarisse, the eldest granddaughter, has organized everything. She has booked a hall near the Mokolo market and she has already invited more than sixty guests. “I have never seen Grandmother so happy,” she says. “She has talked about this party for months.”

When the guests arrived on Saturday morning, the cooks had already prepared the eru and the ndolé. Clarisse had decorated the hall the night before, so everything was ready. By the time the pastor came to bless the meal, the children had eaten all the puff-puff! Nobody was angry: the family has always believed that a party without hungry children is not a real party.

Since that day, Grandmother Ekotto has shown the photos to everyone in the neighbourhood. “I have lived in Yaoundé since 1975,” she told her friends, “and I have never had such a wonderful birthday.”
\end{textebox}

\textbf{Glossaire}\par

\begin{center}
\begin{tabularx}{\linewidth}{|l|l|X|}
\hline
\rowcolor{popblueL} \textbf{Mot anglais} & \textbf{Nature} & \textbf{Traduction} \\
\hline
relatives & nom pluriel & les membres de la famille, les parents \\
\hline
abroad & adverbe & à l'étranger \\
\hline
eldest & adjectif & aîné(e) \\
\hline
to book & verbe & réserver \\
\hline
a hall & nom & une salle (des fêtes) \\
\hline
cooks & nom pluriel & cuisiniers, cuisinières \\
\hline
to bless & verbe & bénir \\
\hline
puff-puff & nom & beignets (pâte frite) \\
\hline
neighbourhood & nom & quartier, voisinage \\
\hline
\end{tabularx}
\end{center}

\textbf{Questions de compréhension}\par

\begin{enumerate}
\item How old is Grandmother Ekotto?
\item What had the cooks done before the guests arrived?
\item Why was nobody angry with the children?
\item How long has Grandmother Ekotto lived in Yaoundé?
\end{enumerate}

\textbf{Observation grammaticale.}\par Relevons quelques phrases du texte :

\begin{itemize}
\item \eng{Clarisse has organized everything.} — le verbe est formé de \eng{has} + \eng{organized}.
\item \eng{She has never had such a wonderful birthday.} — \eng{has} + \eng{had} (participe passé de \eng{have}).
\item \eng{The cooks had already prepared the eru.} — le verbe est formé de \eng{had} + \eng{prepared}.
\item \eng{Clarisse had decorated the hall the night before.} — \eng{had} + \eng{decorated}.
\end{itemize}

Deux structures apparaissent clairement : \eng{have/has + participe passé} et \eng{had + participe passé}. Voyons maintenant les règles en détail.

\courssub{Formation des deux temps}

\begin{propriete}[Formation du present perfect et du past perfect]
\textbf{Present perfect} = \cle{have/has} + participe passé du verbe.
\begin{itemize}
\item \eng{have} avec I, you, we, they ; \eng{has} avec he, she, it.
\item Négation : \eng{have not / has not} (formes contractées : \eng{haven't}, \eng{hasn't}).
\item Interrogation : inversion de \eng{have/has} et du sujet : \eng{Have you finished?}
\end{itemize}
\textbf{Past perfect} (ou pluperfect) = \cle{had} + participe passé du verbe, pour TOUTES les personnes.
\begin{itemize}
\item Négation : \eng{had not} (forme contractée : \eng{hadn't}).
\item Interrogation : \eng{Had they left?}
\end{itemize}
Le participe passé des verbes réguliers se forme avec \textbf{-ed} (\eng{organized}, \eng{decorated}) ; celui des verbes irréguliers est la troisième colonne de la liste (\eng{go – went – gone}, \eng{see – saw – seen}, \eng{eat – ate – eaten}).
\end{propriete}

\begin{center}
\begin{tabular}{|l|l|l|}
\hline
\rowcolor{popgreenL} \textbf{Personne} & \textbf{Present perfect} & \textbf{Past perfect} \\
\hline
I / you / we / they & have eaten & had eaten \\
\hline
he / she / it & has eaten & had eaten \\
\hline
négation & haven't / hasn't eaten & hadn't eaten \\
\hline
question & Have you eaten? / Has she eaten? & Had you eaten? \\
\hline
\end{tabular}
\end{center}

\courssub{Emplois du present perfect}

Le present perfect établit toujours un \textbf{lien entre le passé et le présent}. On l'emploie dans trois situations principales.

\begin{propriete}[Les trois emplois du present perfect]
\begin{enumerate}
\item \textbf{Action passée avec résultat présent} : l'action est terminée, mais sa conséquence est visible maintenant. \eng{I have lost my keys.} (résultat : je ne peux pas entrer chez moi maintenant).
\item \textbf{Expérience de vie} : on parle de ce qu'on a fait ou non au cours de sa vie, sans date précise. \eng{Have you ever tasted eru?}
\item \textbf{Action qui continue jusqu'au présent} : avec \eng{since} (point de départ) et \eng{for} (durée). \eng{She has lived in Douala since 2015.}
\end{enumerate}
\end{propriete}

\begin{exemplebox}[Exemples traduits]
\begin{itemize}
\item \eng{I have just eaten.} « Je viens de manger. » (résultat présent : je n'ai plus faim)
\item \eng{She has lived in Douala since 2015.} « Elle vit à Douala depuis 2015. » (elle y vit encore)
\item \eng{We have known each other for ten years.} « Nous nous connaissons depuis dix ans. »
\item \eng{Have you ever visited the Bamenda market?} « As-tu déjà visité le marché de Bamenda ? »
\item \eng{My father has never travelled abroad.} « Mon père n'a jamais voyagé à l'étranger. »
\item \eng{They have recently moved to Buea.} « Ils ont récemment déménagé à Buea. »
\end{itemize}
\end{exemplebox}

\begin{aretenir}[Les marqueurs du present perfect]
Certains mots signalent presque toujours le present perfect :
\begin{itemize}
\item \eng{just} = venir de ; \eng{already} = déjà ; \eng{yet} = déjà/encore (phrases négatives et interrogatives, en fin de phrase) ; \eng{ever} = déjà (questions) ; \eng{never} = jamais ; \eng{recently} = récemment ; \eng{since} = depuis (point de départ) ; \eng{for} = depuis (durée).
\end{itemize}
\eng{Have you finished your homework yet?} « As-tu déjà fini tes devoirs ? » — \eng{No, I haven't finished yet.} « Non, je n'ai pas encore fini. »
\end{aretenir}

\courssub{Emploi du past perfect : l'antériorité dans le passé}

\begin{propriete}[Le past perfect : une action AVANT une autre action passée]
Le past perfect exprime une action \textbf{achevée avant} une autre action passée (exprimée au prétérit). C'est l'équivalent du plus-que-parfait français.
\begin{itemize}
\item Combinaison fréquente : \eng{after} + past perfect, \eng{before} + prétérit.
\item \eng{After the guests had left, we cleaned the house.} « Après que les invités furent partis, nous avons nettoyé la maison. »
\item \eng{The bus had already left when we reached the station.} « Le bus était déjà parti quand nous sommes arrivés à la gare. »
\end{itemize}
\end{propriete}

\begin{exemplebox}[Exemples traduits]
\begin{itemize}
\item \eng{When I arrived, the guests had left.} « Quand je suis arrivé, les invités étaient partis. »
\item \eng{She had finished her studies before she got married.} « Elle avait terminé ses études avant de se marier. »
\item \eng{By the time we woke up, Mother had prepared breakfast.} « Quand nous nous sommes réveillés, maman avait préparé le petit-déjeuner. »
\item \eng{I had never seen the sea before I went to Kribi.} « Je n'avais jamais vu la mer avant d'aller à Kribi. »
\end{itemize}
\end{exemplebox}

\begin{popfigure}
\begin{center}
\begin{tikzpicture}[font=\small]
\draw[-{Stealth}, thick, popdark] (0,0) -- (12.5,0);
% points
\fill[poporange] (2.5,0) circle (0.12);
\fill[popblue] (6.5,0) circle (0.12);
\fill[poppink] (10.5,0) circle (0.12);
% labels
\node[above=0.25cm of {(2.5,0)}, align=center, text=poporange] {\textbf{Past perfect}\\ action 1 (la plus ancienne)\\ \eng{the guests had left}};
\node[above=0.25cm of {(6.5,0)}, align=center, text=popblue] {\textbf{Prétérit}\\ action 2 (passé)\\ \eng{when I arrived}};
\node[above=0.25cm of {(10.5,0)}, align=center, text=poppink] {\textbf{NOW}\\ present perfect = bilan\\ \eng{I have missed them}};
% flèche temps
\node[below=0.3cm of {(6.5,0)}, text=popmuted] {passé $\longrightarrow$ présent};
\end{tikzpicture}
\end{center}
\legende{Frise des temps : le past perfect précède le prétérit ; le present perfect relie le passé au moment présent.}
\end{popfigure}

\courssub{Comparaison avec le français et pièges}

\begin{center}
\begin{tabularx}{\linewidth}{|X|X|X|}
\hline
\rowcolor{popblueL} \textbf{Français} & \textbf{English} & \textbf{Remarque} \\
\hline
Je viens de manger. & I have just eaten. & \eng{just} + present perfect \\
\hline
Elle vit à Douala depuis 2015. & She has lived in Douala since 2015. & le français utilise le présent, l'anglais le present perfect \\
\hline
Je l'ai vu hier. & I saw him yesterday. & repère passé précis $\Rightarrow$ prétérit obligatoire \\
\hline
Quand je suis arrivé, ils étaient partis. & When I arrived, they had left. & plus-que-parfait = past perfect \\
\hline
\end{tabularx}
\end{center}

\begin{attention}[L'erreur capitale des francophones]
En français, le passé composé accepte un repère de temps passé : « Je l'ai vu \textbf{hier} ». En anglais, c'est \textbf{impossible} avec le present perfect !
\begin{itemize}
\item FAUX : \eng{I have seen him yesterday.}
\item CORRECT : \eng{I saw him yesterday.}
\end{itemize}
Dès qu'un repère passé précis apparaît (\eng{yesterday}, \eng{last week}, \eng{in 2019}, \eng{two days ago}, \eng{when I was a child}), il faut le \textbf{prétérit}. Le present perfect exige un lien avec le présent : pas de date passée précise !
Autre piège : ne dites pas \eng{I am living here since 2015} ; dites \eng{I have lived here since 2015.}
\end{attention}

\begin{methode}[Quel temps choisir ? Trois questions]
\begin{enumerate}
\item Y a-t-il un \textbf{repère passé précis} (yesterday, last year, ago, when...) ? $\Rightarrow$ \textbf{prétérit}. \eng{I met her last year.}
\item S'agit-il d'un \textbf{bilan, d'un résultat présent, d'une expérience ou d'une action qui continue} ? $\Rightarrow$ \textbf{present perfect}. \eng{I have just met her.}
\item L'action est-elle \textbf{antérieure à une autre action passée} ? $\Rightarrow$ \textbf{past perfect}. \eng{I had met her before the party.}
\end{enumerate}
\end{methode}

\begin{savaistu}[For ou since ?]
\eng{for} + une \textbf{durée} (\eng{for ten years}, \eng{for a long time}) et \eng{since} + un \textbf{point de départ} (\eng{since 2015}, \eng{since Monday}, \eng{since I was born}) sont presque toujours accompagnés du present perfect. Astuce : si vous pouvez répondre à la question « pendant combien de temps ? », utilisez \eng{for} ; si vous répondez « à partir de quel moment ? », utilisez \eng{since}. \eng{Grandmother has lived in Yaoundé since 1975 — for almost fifty years!}
\end{savaistu}

\courssub{Entraînement}

\begin{exoresolu}[Complète avec le temps correct]
Complète les phrases avec le verbe entre parenthèses au present perfect, au past perfect ou au prétérit.
\begin{enumerate}
\item I \_\_\_\_\_\_\_\_ (never / see) a bendskin accident before I came to Douala.
\item My sister \_\_\_\_\_\_\_\_ (just / pass) her Probatoire.
\item We \_\_\_\_\_\_\_\_ (visit) our grandparents last Sunday.
\item When the teacher arrived, the students \_\_\_\_\_\_\_\_ (already / clean) the classroom.
\item They \_\_\_\_\_\_\_\_ (know) each other since primary school.
\end{enumerate}
\tcblower
\textbf{Solution détaillée}
\begin{enumerate}
\item \eng{I had never seen a bendskin accident before I came to Douala.} — action antérieure à une autre action passée (\eng{before I came}) $\Rightarrow$ past perfect.
\item \eng{My sister has just passed her Probatoire.} — \eng{just} + résultat présent $\Rightarrow$ present perfect.
\item \eng{We visited our grandparents last Sunday.} — repère passé précis (\eng{last Sunday}) $\Rightarrow$ prétérit.
\item \eng{When the teacher arrived, the students had already cleaned the classroom.} — le nettoyage est antérieur à l'arrivée $\Rightarrow$ past perfect.
\item \eng{They have known each other since primary school.} — \eng{since} + action qui continue $\Rightarrow$ present perfect.
\end{enumerate}
\end{exoresolu}

\begin{exercice}[À toi de jouer]
Mets les verbes entre parenthèses au temps correct (prétérit, present perfect ou past perfect).
\begin{enumerate}
\item Amina \_\_\_\_\_\_\_\_ (live) in Bamenda since 2018.
\item I \_\_\_\_\_\_\_\_ (not / finish) my composition yet.
\item The match \_\_\_\_\_\_\_\_ (already / start) when we reached the stadium.
\item My uncle \_\_\_\_\_\_\_\_ (buy) a new car two months ago.
\item \_\_\_\_\_\_\_\_ you ever \_\_\_\_\_\_\_\_ (eat) ndolé?
\item After the rain \_\_\_\_\_\_\_\_ (stop), the children went out to play football.
\end{enumerate}
\end{exercice}

\begin{corrige}[Exercice « À toi de jouer »]
\begin{enumerate}
\item \eng{has lived} — \eng{since} + action qui continue $\Rightarrow$ present perfect.
\item \eng{have not finished} (haven't finished) — \eng{yet} en phrase négative $\Rightarrow$ present perfect.
\item \eng{had already started} — le début du match est antérieur à notre arrivée $\Rightarrow$ past perfect.
\item \eng{bought} — repère passé précis (\eng{two months ago}) $\Rightarrow$ prétérit.
\item \eng{Have you ever eaten} — expérience de vie $\Rightarrow$ present perfect.
\item \eng{had stopped} — \eng{after} + action antérieure $\Rightarrow$ past perfect.
\end{enumerate}
\end{corrige}

\begin{aretenir}[L'essentiel de la section]
\begin{itemize}
\item \textbf{Present perfect} = \eng{have/has + participe passé} : résultat présent, expérience de vie, action qui continue (\eng{since/for}).
\item \textbf{Past perfect} = \eng{had + participe passé} : action achevée avant une autre action passée.
\item Jamais de present perfect avec un repère passé précis : \eng{yesterday}, \eng{ago}, \eng{last week} $\Rightarrow$ prétérit.
\item Marqueurs : \eng{just, already, yet, ever, never, recently, since, for}.
\end{itemize}
\end{aretenir}

\coursec{Entraînement et production}

Nous venons d'étudier en détail le \cle{present perfect} et le \cle{past perfect}, ainsi que toutes les structures permettant d'exprimer le but, la conséquence, l'obligation et la préférence. Il est temps de vérifier que vous savez les \textbf{mobiliser ensemble}, comme on vous le demandera au contrôle et à l'examen. Cette dernière partie vous propose d'abord une méthode de travail, puis cinq exercices d'application corrigés, et enfin une production écrite guidée.

\courssub{La méthode pour réussir les exercices de transformation}

\begin{methode}[Réussir les transformations et les complétions]
Face à n'importe quel exercice de grammaire de cette leçon, suivez toujours les trois étapes suivantes :
\begin{enumerate}
\item \textbf{Identifier la fonction} : que veut dire la phrase ? Un but (pourquoi ?), une conséquence (avec quel résultat ?), une obligation (est-ce obligatoire, interdit ou inutile ?), une préférence, ou une action située dans le temps ?
\item \textbf{Choisir la structure} correspondant à cette fonction : \eng{to / in order to / so as to} pour le but, \eng{so... that / such... that} pour la conséquence, \eng{must / mustn't / don't have to} pour l'obligation, \eng{would rather / prefer... to} pour la préférence.
\item \textbf{Vérifier la construction} : base verbale après les modaux et après \eng{let's} ; \eng{to} + base verbale après \eng{prefer} ; une proposition avec \eng{that} + sujet + verbe après \eng{so} et \eng{such} ; participe passé après \eng{have/has} au present perfect.
\end{enumerate}
\end{methode}

\begin{popfigure}
\begin{center}
\begin{tikzpicture}[
  node distance=0.7cm and 0.5cm,
  start/.style={rectangle, rounded corners, draw=popblue, fill=popblueL, text=popdark, align=center, font=\small\bfseries, inner sep=6pt},
  fonc/.style={rectangle, rounded corners, draw=poppurple, fill=poppurpleL, text=popdark, align=center, font=\small, inner sep=5pt},
  sol/.style={rectangle, draw=popgreen, fill=popgreenL, text=popdark, align=center, font=\small, inner sep=5pt},
  fleche/.style={-{Stealth[length=2.5mm]}, thick, popmuted}
]
\node[start, align=center] (q){Quelle fonction\\ veux-je exprimer ?};
\node[fonc, below left=0.8cm and 2.6cm of q, align=center] (but){But\\ (pourquoi ?)};
\node[fonc, below=0.8cm of q, align=center] (cons){Conséquence\\ (résultat)};
\node[fonc, below right=0.8cm and 2.6cm of q, align=center] (obl){Obligation /\\ préférence};
\node[fonc, below=1.6cm of cons, xshift=3.4cm, align=center] (temps){Temps\\ (action passée)};
\node[sol, below=0.6cm of but, align=center] (sbut){to / in order to + BV\\ so that + sujet + verbe};
\node[sol, below=0.6cm of cons, align=center] (scons){so + adj + that\\ such + (a) + nom + that};
\node[sol, below=0.6cm of obl, align=center] (sobl){must / mustn't / don't have to\\ would rather + BV};
\node[sol, below=0.6cm of temps, align=center] (stemps){present perfect : have + p.p.\\ past perfect : had + p.p.};
\draw[fleche] (q) -- (but);
\draw[fleche] (q) -- (cons);
\draw[fleche] (q) -- (obl);
\draw[fleche] (q.south) |- (temps.west);
\draw[fleche] (but) -- (sbut);
\draw[fleche] (cons) -- (scons);
\draw[fleche] (obl) -- (sobl);
\draw[fleche] (temps) -- (stemps);
\end{tikzpicture}
\end{center}
\legende{Organigramme de choix : de la fonction visée à la structure correcte (BV = base verbale, p.p. = participe passé).}
\end{popfigure}

\begin{attention}[Les 5 erreurs à bannir absolument]
\begin{enumerate}
\item \eng{for to} : jamais \emph{for to} pour exprimer le but. On dit \eng{I study to pass my exam.} (« j'étudie pour réussir »), pas \emph{for to pass}.
\item \eng{must to} : après \eng{must}, la base verbale \textbf{sans} to. \eng{You must work hard.}, pas \emph{must to work}.
\item \eng{let's to} : après \eng{let's}, la base verbale seule. \eng{Let's go home.}, pas \emph{let's to go}.
\item \eng{prefer... than} : la préférence se construit avec \eng{to}, jamais \emph{than}. \eng{I prefer tea to coffee.} (« je préfère le thé au café »).
\item \eng{present perfect + yesterday} : avec un repère de temps terminé (\eng{yesterday}, \eng{last week}, \eng{in 2019}), on utilise le prétérit. \eng{I saw him yesterday.}, pas \emph{I have seen him yesterday}.
\end{enumerate}
\end{attention}

\courssub{Exercice 1 : compléter avec to, in order to ou so that}

\begin{exoresolu}[Exercice 1 — But : to, in order to ou so that]
Complétez les phrases avec \eng{to}, \eng{in order to} ou \eng{so that}.
\begin{enumerate}
\item My father works two jobs \ldots\ feed the family.
\item She left home early \ldots\ she would not miss the bus to Douala.
\item We went to the market \ldots\ buy some vegetables.
\item He speaks slowly \ldots\ his little sister can understand him.
\item I am saving money \ldots\ buy a new bicycle.
\end{enumerate}
\tcblower
\textbf{Solution détaillée.}
\begin{enumerate}
\item \eng{to} (ou \eng{in order to}) : but suivi d'un verbe, même sujet. « Mon père fait deux métiers pour nourrir la famille. »
\item \eng{so that} : le sujet change (\eng{she} / \eng{she would not miss}) ; on a besoin d'une proposition complète. « Elle est partie tôt de la maison afin de ne pas rater le car pour Douala. »
\item \eng{to} (ou \eng{in order to}) : même sujet, but + verbe. « Nous sommes allés au marché pour acheter des légumes. »
\item \eng{so that} : deux sujets différents (\eng{he} / \eng{his sister}), donc proposition avec \eng{that}. « Il parle lentement pour que sa petite sœur puisse le comprendre. »
\item \eng{in order to} (ou \eng{to}) : même sujet. « J'économise de l'argent pour acheter un nouveau vélo. »
\end{enumerate}
\textbf{Règle à retenir} : même sujet $\rightarrow$ \eng{to} + base verbale (ou \eng{in order to}) ; sujets différents $\rightarrow$ \eng{so that} + sujet + verbe (souvent avec un modal \eng{can/could/will/would}).
\end{exoresolu}

\courssub{Exercice 2 : transformer so... that en such... that}

Rappel de la règle : \eng{so} est suivi d'un \textbf{adjectif} (\eng{so tired that...}), tandis que \eng{such} est suivi d'un \textbf{nom} (avec \eng{a/an} si le nom est comptable au singulier : \eng{such a nice day that...}). Le sens reste identique.

\begin{exoresolu}[Exercice 2 — So... that / Such... that]
Transformez les phrases sans changer le sens.
\begin{enumerate}
\item The food was so good that we ate everything. $\rightarrow$ It was \ldots
\item It was such a hot day that we stayed indoors. $\rightarrow$ The day was \ldots
\item The bendskin was so fast that I arrived in ten minutes. $\rightarrow$ It was \ldots
\end{enumerate}
\tcblower
\textbf{Solution détaillée.}
\begin{enumerate}
\item \eng{It was such good food that we ate everything.} — \eng{food} est un nom (indénombrable, donc sans \eng{a}). « La nourriture était si bonne que nous avons tout mangé. »
\item \eng{The day was so hot that we stayed indoors.} — \eng{hot} est un adjectif, donc \eng{so}. « Il faisait si chaud que nous sommes restés à l'intérieur. »
\item \eng{It was such a fast bendskin that I arrived in ten minutes.} — nom comptable singulier, donc \eng{such a}. « La moto-taxi était si rapide que je suis arrivé en dix minutes. »
\end{enumerate}
\textbf{Astuce} : repérez le mot qui suit — adjectif $\rightarrow$ \eng{so} ; nom $\rightarrow$ \eng{such}.
\end{exoresolu}

\courssub{Exercice 3 : must, mustn't ou don't have to}

Rappel : \eng{must} = obligation (« il faut ») ; \eng{mustn't} = interdiction (« il ne faut pas », « c'est défendu ») ; \eng{don't have to} = absence d'obligation (« ce n'est pas la peine », « tu n'es pas obligé »). Attention : \eng{mustn't} et \eng{don't have to} ne sont \textbf{pas} synonymes ! Rappel d'accord : à la 3\up{e} personne du singulier, on écrit \eng{doesn't have to}.

\begin{exercice}[Exercice 3 — Obligation, interdiction ou absence d'obligation]
Complétez avec \eng{must}, \eng{mustn't}, \eng{don't have to} ou \eng{doesn't have to}.
\begin{enumerate}
\item You \ldots\ wear a uniform at our school; it is compulsory.
\item You \ldots\ bring food to the party; the school will provide everything.
\item Students \ldots\ use their phones during the examination; it is forbidden.
\item She \ldots\ hurry; we still have plenty of time.
\item We \ldots\ respect our parents and teachers.
\end{enumerate}
\end{exercice}

\begin{corrige}[Exercice 3]
\begin{enumerate}
\item \eng{must} — obligation (l'uniforme est obligatoire). « Tu dois porter l'uniforme. »
\item \eng{don't have to} — inutile, pas obligé. « Tu n'es pas obligé d'apporter de la nourriture. »
\item \eng{mustn't} — interdiction formelle. « Les élèves ne doivent pas utiliser leur téléphone pendant l'examen. »
\item \eng{doesn't have to} — pas nécessaire (3\up{e} personne du singulier : \eng{she}). « Elle n'a pas besoin de se dépêcher. »
\item \eng{must} — devoir moral. « Nous devons respecter nos parents et nos professeurs. »
\end{enumerate}
\end{corrige}

\courssub{Exercice 4 : present perfect, prétérit ou past perfect}

Rappel de la règle de choix :
\begin{itemize}
\item \textbf{Prétérit} : action passée, datée, terminée (\eng{yesterday}, \eng{last year}, \eng{in 2020}).
\item \textbf{Present perfect} : action passée sans date précise, avec un lien avec le présent (\eng{already}, \eng{just}, \eng{never}, \eng{ever}, \eng{since}, \eng{for}).
\item \textbf{Past perfect} : action \emph{antérieure} à une autre action passée (souvent avec \eng{before}, \eng{after}, \eng{when}, \eng{because}).
\end{itemize}

\begin{exoresolu}[Exercice 4 — Conjuguer au bon temps]
Mettez le verbe entre parenthèses au temps correct.
\begin{enumerate}
\item I \ldots\ (never / visit) Bamenda, but I would love to go there.
\item My uncle \ldots\ (leave) for Buea last Monday.
\item When we arrived at the station, the bus \ldots\ (already / depart).
\item She \ldots\ (just / finish) her homework, so she can play now.
\item They \ldots\ (live) in Yaoundé since 2018.
\end{enumerate}
\tcblower
\textbf{Solution détaillée.}
\begin{enumerate}
\item \eng{have never visited} — present perfect avec \eng{never} : expérience de vie sans date. « Je n'ai jamais visité Bamenda. »
\item \eng{left} — prétérit : repère daté \eng{last Monday}. « Mon oncle est parti pour Buea lundi dernier. »
\item \eng{had already departed} — past perfect : le départ du car est \emph{antérieur} à notre arrivée. « Quand nous sommes arrivés à la gare, le car était déjà parti. »
\item \eng{has just finished} — present perfect avec \eng{just} : résultat visible maintenant. « Elle vient de finir ses devoirs. »
\item \eng{have lived} — present perfect avec \eng{since} : action qui continue jusqu'à aujourd'hui. « Ils habitent à Yaoundé depuis 2018. »
\end{enumerate}
\end{exoresolu}

\courssub{Exercice 5 : corriger les erreurs typiques des francophones}

\begin{exercice}[Exercice 5 — À chacun son erreur]
Chaque phrase contient une erreur vue dans cette leçon. Corrigez-la.
\begin{enumerate}
\item I went to the library for to borrow a book.
\item You must to respect the school rules.
\item Let's to visit our grandmother this weekend.
\item I prefer football than volleyball.
\item I have seen that film yesterday.
\end{enumerate}
\end{exercice}

\begin{corrige}[Exercice 5]
\begin{enumerate}
\item \eng{I went to the library to borrow a book.} — le but s'exprime avec \eng{to}, jamais \emph{for to}. « Je suis allé à la bibliothèque pour emprunter un livre. »
\item \eng{You must respect the school rules.} — base verbale sans \eng{to} après \eng{must}. « Tu dois respecter le règlement de l'école. »
\item \eng{Let's visit our grandmother this weekend.} — base verbale seule après \eng{let's}. « Rendons visite à notre grand-mère ce week-end. »
\item \eng{I prefer football to volleyball.} — la préférence se construit avec \eng{to}. « Je préfère le football au volley-ball. »
\item \eng{I saw that film yesterday.} — avec \eng{yesterday}, on utilise le prétérit. « J'ai vu ce film hier. »
\end{enumerate}
\end{corrige}

\begin{center}
\begin{tabularx}{\linewidth}{|X|X|X|}
\hline
\rowcolor{popblueL} \textbf{Erreur du francophone} & \textbf{Forme correcte} & \textbf{Pourquoi} \\
\hline
for to + verbe & to + BV & le but = \eng{to} seul \\
\hline
must to + verbe & must + BV & modal suivi de la base verbale \\
\hline
let's to + verbe & let's + BV & impératif à la 1\textsuperscript{re} personne du pluriel \\
\hline
prefer... than & prefer... to & comparaison de préférence avec \eng{to} \\
\hline
present perfect + yesterday & prétérit + yesterday & repère de temps terminé \\
\hline
\end{tabularx}
\end{center}

\courssub{Vocabulaire utile pour la production}

\begin{definition}[Mots et expressions du contexte camerounais]
\begin{itemize}
\item \eng{bendskin} (n.) : moto-taxi, moyen de transport très courant.
\item \eng{traffic jam} / \eng{hold-up} (n.) : embouteillage.
\item \eng{quarter} (n.) : quartier (résidentiel).
\item \eng{compound} (n.) : concession, cour familiale clôturée.
\item \eng{provisions} (n. pl.) : vivres, courses alimentaires.
\item \eng{revising} (n.) : révision (scolaire).
\item \eng{to chill out} (phr. v., fam.) : se détendre, se reposer.
\end{itemize}
\end{definition}

\courssub{Production guidée : My family plans for the weekend}

\begin{exemplebox}[Phrase modèle de production]
\eng{We have decided to visit our uncle in Buea this weekend. Let's leave early so that we can avoid the traffic. I would rather travel by bus because it is cheaper.}

« Nous avons décidé de rendre visite à notre oncle à Buea ce week-end. Partons tôt afin d'éviter les embouteillages. Je préférerais voyager en car parce que c'est moins cher. »

On y reconnaît : present perfect (\eng{have decided}), suggestion (\eng{let's} + BV), but avec \eng{so that}, préférence avec \eng{would rather} + BV.
\end{exemplebox}

\begin{methode}[Rédiger votre paragraphe en 4 étapes]
\begin{enumerate}
\item \textbf{Préparez} : notez 5 ou 6 activités de week-end en famille (visiter un parent, faire le marché, aller au stade, réviser, se reposer).
\item \textbf{Planifiez les structures} : choisissez au moins 5 structures de la leçon — un but (\eng{to / so that}), une conséquence (\eng{so... that}), une obligation (\eng{must}), une absence d'obligation (\eng{don't have to}), une suggestion (\eng{let's}), une préférence (\eng{would rather / prefer... to}), un present perfect.
\item \textbf{Rédigez} 8 à 10 phrases liées par des connecteurs (\eng{first}, \eng{then}, \eng{because}, \eng{finally}).
\item \textbf{Relisez} avec la liste des 5 erreurs à bannir : pas de \emph{for to}, pas de \emph{must to}, pas de \emph{let's to}, pas de \emph{prefer... than}, pas de present perfect avec \eng{yesterday}.
\end{enumerate}
\end{methode}

\begin{exercice}[Production écrite — My family plans for the weekend]
Rédigez un paragraphe de 8 à 10 phrases intitulé \eng{My family plans for the weekend}. Utilisez \textbf{au moins cinq} des structures suivantes : \eng{to / in order to / so that} ; \eng{so... that / such... that} ; \eng{must / mustn't / don't have to} ; \eng{let's / let's not} ; \eng{would rather / prefer... to} ; present perfect ou past perfect. Situez votre texte dans un contexte camerounais (Yaoundé, Douala, Buea, le village, le marché\ldots).
\end{exercice}

\begin{experience}[Correction collective et jeu de rôle]
Par groupes de quatre : chaque élève lit son paragraphe à voix haute. Les autres écoutent et lèvent la main dès qu'ils entendent une des cinq structures de la leçon, qu'ils doivent nommer (\eng{purpose!}, \eng{result!}, \eng{obligation!}\ldots). Le groupe choisit ensuite le meilleur texte et le présente à la classe. Le professeur relève au tableau les erreurs fréquentes entendues et fait reformuler les phrases fautives par les élèves eux-mêmes.
\end{experience}

\begin{aretenir}[Ce que vous devez savoir faire après cette leçon]
\begin{itemize}
\item Exprimer un but : \eng{to / in order to} + BV (même sujet) ; \eng{so that} + proposition (sujets différents).
\item Exprimer une conséquence : \eng{so} + adjectif + \eng{that} ; \eng{such} + (a) + nom + \eng{that}.
\item Nuancer l'obligation : \eng{must} (obligatoire), \eng{mustn't} (interdit), \eng{don't have to} (pas obligé).
\item Proposer : \eng{let's / let's not} + base verbale.
\item Exprimer une préférence : \eng{would rather} + BV ; \eng{prefer} + nom/\eng{-ing} + \eng{to}.
\item Choisir entre prétérit (action datée), present perfect (lien avec le présent) et past perfect (antériorité).
\end{itemize}
\end{aretenir}

\newpage

\coursec{Activité d'intégration}

\coursec{Lesson 5 — Grammar: Expressing Purpose, Result, Obligation, Preference; Let's / Let's Not + Bare Infinitive; Present and Past Perfect}

\courssub{Objectifs de l'activité}

Cette activité d'intégration clôt la leçon de grammaire en faisant travailler les élèves sur \textbf{toutes} les structures étudiées, dans une tâche de communication réelle et motivante : organiser une fête de famille. À la fin de l'activité, l'élève doit être capable de :

\begin{itemize}
\item exprimer le \cle{but} avec \eng{to}, \eng{in order to}, \eng{so as to} et \eng{so that} ;
\item exprimer la \cle{conséquence} avec \eng{so}, \eng{so... that} et \eng{such... that} ;
\item exprimer l'\cle{obligation} et l'\cle{interdiction} avec \eng{must}, \eng{have to}, \eng{mustn't} et \eng{don't have to} ;
\item exprimer la \cle{préférence} avec \eng{would rather... than} et \eng{prefer... to} ;
\item faire des \cle{propositions} avec \eng{let's} et \eng{let's not} + base verbale ;
\item employer correctement le \cle{present perfect} et le \cle{past perfect} pour parler de ce qui est déjà fait ou de ce qui s'était passé avant.
\end{itemize}

\courssub{Situation}

\begin{textebox}[A message from Auntie Ngo Bassa — Yaoundé]
Dear children,

Great news! Your cousin Sandrine is getting married on the 20th of December, and our family has been chosen to organise the traditional reception at home in Mvog-Ada, Yaoundé. We are expecting about one hundred and fifty guests: family, friends, neighbours and colleagues.

Last year's family party was a disaster! The food had finished before half the guests arrived, the generator broke down during the dancing, and we had forgotten to buy enough drinks. This year, everything must be perfect.

I need your help. Please form small committees and send me a clear written plan. Tell me what we have to do, what we mustn't forget, what you would rather do, and what has already been done. The best plan will be followed on the wedding day!

With love,

Auntie Ngo Bassa
\end{textebox}

\textbf{Glossaire :} \emph{wedding} (n. : mariage) ; \emph{reception} (n. : réception) ; \emph{guest} (n. : invité) ; \emph{generator} (n. : groupe électrogène) ; \emph{to break down} (v. : tomber en panne) ; \emph{committee} (n. : comité).

La classe reçoit donc une mission concrète : répondre à Auntie Ngo Bassa en lui proposant un plan d'organisation complet, rédigé en anglais, puis présenté oralement devant la classe (le « conseil de famille »).

\courssub{Consignes}

Travail par \textbf{groupes de 4}. Chaque groupe choisit un rôle pour chaque membre : le \emph{coordinator} (présente le plan), le \emph{secretary} (rédige), le \emph{treasurer} (parle du budget), le \emph{timekeeper} (vérifie le calendrier).

\begin{enumerate}
\item \textbf{Compréhension.} Lisez le message d'Auntie Ngo Bassa. Relevez : (a) deux problèmes de la fête de l'année dernière exprimés au past perfect ; (b) une phrase exprimant l'obligation ; (c) ce que la tante attend de vous.

\item \textbf{Brainstorming lexical.} En 5 minutes, listez dans votre cahier au moins dix mots utiles du champ lexical de la fête : \eng{hall, drinks, DJ, band, chairs, food, decoration, invitation, budget, guests...}

\item \textbf{Préparation grammaticale guidée.} Complétez d'abord ces phrases-cadres avec vos propres idées :
\begin{itemize}
\item \eng{We have to \ldots\ (obligation)}
\item \eng{We mustn't \ldots\ (interdiction)}
\item \eng{We will cook early in order to \ldots\ (but)}
\item \eng{We will hire a big hall so that \ldots\ (but + proposition)}
\item \eng{Last year the food had finished before \ldots, so this year \ldots\ (past perfect + conséquence)}
\item \eng{We would rather \ldots\ than \ldots\ (préférence)}
\item \eng{We have already \ldots\ (present perfect)}
\item \eng{Let's \ldots\ / Let's not \ldots\ (propositions)}
\end{itemize}

\item \textbf{Rédaction.} Rédigez ensemble un plan de \textbf{8 à 10 lignes} intitulé \eng{Our Plan for Sandrine's Wedding}. Votre texte doit contenir \textbf{au moins une occurrence de chaque structure} de la leçon. Soulignez chaque structure utilisée.

\item \textbf{Présentation orale.} Le \emph{coordinator} présente le plan au « conseil de famille » (la classe) en 2 minutes. Les autres membres peuvent ajouter des détails. Parlez lentement, sans lire mot à mot.

\item \textbf{Vote.} La classe écoute chaque groupe et remplit la grille d'évaluation ci-dessous, puis vote pour le meilleur plan. Le groupe gagnant verra son plan « envoyé » à Auntie Ngo Bassa.
\end{enumerate}

\begin{center}
\begin{tabularx}{\linewidth}{|X|c|}\hline
\rowcolor{popblueL} \textbf{Critère} & \textbf{Points} \\ \hline
Toutes les structures de la leçon sont utilisées & 4 \\ \hline
Les structures sont correctes (forme) & 4 \\ \hline
Le plan est réaliste et clair & 2 \\ \hline
La présentation orale est fluide et audible & 2 \\ \hline
\end{tabularx}
\end{center}

\courssub{Ressources à mobiliser}

\begin{aretenir}[Boîte à outils grammaticale]
\begin{itemize}
\item \textbf{But :} \eng{to} / \eng{in order to} / \eng{so as to} + base verbale ; \eng{so that} + sujet + \eng{can/will} + base verbale. \eng{We will cook early so that the guests can eat on time.}
\item \textbf{Conséquence :} \eng{so} + proposition ; \eng{so} + adjectif + \eng{that} ; \eng{such} + (a) + nom + \eng{that}. \eng{The hall was so small that we hired a bigger one.}
\item \textbf{Obligation :} \eng{must} / \eng{have to} + base verbale ; interdiction : \eng{mustn't} ; absence d'obligation : \eng{don't have to}. Attention : \eng{mustn't} $\neq$ \eng{don't have to} !
\item \textbf{Préférence :} \eng{would rather} + base verbale + \eng{than} ; \eng{prefer} + nom/\eng{-ing} + \eng{to}.
\item \textbf{Proposition :} \eng{Let's} / \eng{Let's not} + base verbale (jamais \eng{to} après \eng{let's}).
\item \textbf{Present perfect :} have/has + participe passé, avec \eng{already}, \eng{just}, \eng{yet}, \eng{since}, \eng{for}. \eng{We have already bought the drinks.}
\item \textbf{Past perfect :} had + participe passé, pour l'action antérieure à une autre action passée. \eng{The food had finished before the guests arrived.}
\end{itemize}
\end{aretenir}

\textbf{Lexique utile :} \eng{to book a hall} (réserver une salle), \eng{to hire a band} (engager un orchestre), \eng{caterer} (traiteur), \eng{to decorate} (décorer), \eng{to send invitations}, \eng{drinks} (boissons), \eng{to rent chairs}, \eng{budget} (budget), \eng{on time} (à l'heure).

\courssub{Proposition de production et corrigé commenté}

\begin{exoresolu}[Our Plan for Sandrine's Wedding — modèle de production]
Rédigez un plan de 8 à 10 lignes utilisant toutes les structures de la leçon.
\tcblower
\textbf{Production modèle :}

\eng{Dear Auntie, here is our plan for Sandrine's wedding. First, we have to book the Mvog-Ada community hall before the end of November, and we mustn't forget to pay the deposit. We would rather hire a live band than use a simple speaker, because music is the soul of the party. We will cook the ndolé and the eru early in the morning in order to be ready before noon, so that the guests can eat on time. Last year, the food had finished before half the guests arrived, so this year we will prepare much more. We have already bought the drinks and we have just ordered two hundred chairs. The budget is so small that we must not waste money: let's decorate the house ourselves on Saturday, and let's not spend too much on drinks. We are sure the party will be a great success!}

« Chère tante, voici notre plan pour le mariage de Sandrine. D'abord, nous devons réserver la salle communautaire de Mvog-Ada avant la fin novembre, et nous ne devons pas oublier de payer l'acompte. Nous préférerions engager un orchestre plutôt qu'une simple enceinte, car la musique est l'âme de la fête. Nous cuisinerons le ndolé et l'eru tôt le matin afin d'être prêts avant midi, pour que les invités puissent manger à l'heure. L'année dernière, la nourriture était déjà finie avant que la moitié des invités n'arrive, donc cette année nous en préparerons beaucoup plus. Nous avons déjà acheté les boissons et venons de commander deux cents chaises. Le budget est si petit que nous ne devons pas gaspiller d'argent : décorons la maison nous-mêmes samedi, et ne dépensons pas trop pour les boissons. Nous sommes sûrs que la fête sera un grand succès ! »

\textbf{Commentaire des choix de langue :}
\begin{itemize}
\item \eng{we have to book} : obligation extérieure (il faut réserver tôt, les salles sont demandées) ; \eng{we mustn't forget} : interdiction forte, attention à ne pas confondre avec \eng{don't have to} (pas obligé).
\item \eng{We would rather hire... than use} : préférence, base verbale après \eng{would rather} et après \eng{than} — jamais \eng{to use}.
\item \eng{in order to be ready} : but + base verbale ; \eng{so that the guests can eat} : \eng{so that} + sujet + modal, car le sujet change.
\item \eng{the food had finished before... arrived} : past perfect pour l'action antérieure, simple past pour l'action de référence ; \eng{so this year...} : conséquence.
\item \eng{We have already bought} / \eng{have just ordered} : present perfect avec \eng{already} et \eng{just} pour des actions récentes au résultat présent.
\item \eng{so small that} : structure de conséquence avec adjectif ; \eng{let's decorate} / \eng{let's not spend} : propositions, base verbale sans \eng{to}.
\end{itemize}
\end{exoresolu}

\begin{attention}[Erreurs à éviter dans vos productions]
Ne dites pas \eng{let's to decorate} : après \eng{let's}, base verbale seule. Ne dites pas \eng{we mustn't to book}. Ne confondez pas \eng{mustn't} (interdit) et \eng{don't have to} (pas nécessaire). N'employez pas le present perfect avec une date passée : on dit \eng{we bought the drinks last week}, mais \eng{we have already bought the drinks}. Enfin, \eng{so that} est suivi d'une proposition complète avec sujet, pas d'un infinitif.
\end{attention}

\courssub{Bilan de l'activité}

À l'issue de cette tâche, faites le point collectivement au tableau : demandez à chaque groupe de citer une phrase de son plan et de nommer la structure utilisée. Vérifiez que chacun sait désormais :

\begin{itemize}
\item choisir entre \eng{to / in order to} (même sujet) et \eng{so that} (sujet différent) pour exprimer le but ;
\item construire une conséquence avec \eng{so}, \eng{so... that}, \eng{such... that} ;
\item distinguer \eng{must}, \eng{have to}, \eng{mustn't} et \eng{don't have to} ;
\item utiliser \eng{would rather... than} et \eng{let's / let's not} + base verbale ;
\item situer les actions dans le temps grâce au present perfect (bilan présent) et au past perfect (antériorité dans le passé).
\end{itemize}

Cette activité montre que la grammaire n'est pas une fin en soi : exprimer le but, l'obligation ou la préférence sert à \textbf{agir dans la vie réelle} — ici, organiser avec succès un événement familial, situation que tout élève camerounais connaît bien. Le vote final valorise la coopération, la prise de parole en public et la correction de la langue, trois compétences au cœur du programme de Première.

\newpage

\coursec{Méthodes et exercices résolus}

\coursec{Lesson 5 — Grammar: Expressing purpose, result, obligation, preference, let's/let's not, present/past perfect}

\courssub{Observation : découvrir les structures en contexte}

\begin{textebox}[Dialogue : Planning a family reunion]
\textbf{Context.} Amina and Peter, cousins, are organizing a family reunion in Yaoundé.

Amina: We \textbf{must} book the hall at Mvog-Betsi early. It is compulsory for big events.
Peter: True. But we \textbf{don't have to} pay the full amount today; a deposit is enough.
Amina: Good. \textbf{Let's} pay the deposit tomorrow \textbf{so that} we secure the date.
Peter: \textbf{Let's not} forget to order the fufu corn and eru. Everyone prefers it \textbf{to} rice.
Amina: I \textbf{would rather} cook it myself \textbf{than} hire a caterer; it tastes better.
Peter: You \textbf{have cooked} for the whole family before, \textbf{so} you know the quantities.
Amina: Yes, but last year, by the time the guests \textbf{arrived}, the food \textbf{had already finished}!
Peter: \textbf{So} many people came \textbf{that} we ran out of drinks. \textbf{Therefore}, this year we will buy more.
Amina: It was \textbf{such a busy day} \textbf{that} I \textbf{hadn't even eaten} myself.
Peter: \textbf{Let's} make a checklist \textbf{in order not to} miss anything.
Amina: Agreed. \textbf{Let me} check the guest list first.
\end{textebox}

\begin{exoresolu}[Observation guidée]
\textbf{Énoncé.} Relisez le dialogue et répondez aux questions.
\begin{enumerate}
\item Relevez trois structures exprimant le \textbf{but} (purpose). Classez-les selon que les sujets sont identiques ou différents.
\item Trouvez deux structures exprimant la \textbf{conséquence} (result). Quelle différence de structure notez-vous ?
\item Identifiez : une obligation intérieure, une obligation extérieure, une absence d'obligation, une interdiction (absente ici, mais forme connue).
\item Quelles formes pour exprimer la \textbf{préférence} ? Notez les prépositions ou conjonctions qui les suivent.
\item Repérez \textbf{let's}, \textbf{let's not} et \textbf{let me}. Quelle est la forme verbale qui suit ?
\item Surlignez les verbes au \textbf{present perfect} et au \textbf{past perfect}. Justifiez chaque emploi.
\end{enumerate}
\tcblower
\textbf{Solution détaillée.}
\begin{enumerate}
\item But : \eng{in order to secure} (même sujet, \eng{to/in order to} + BV) ; \eng{so that we secure} (sujets différents \eng{we/we} mais focus sur résultat, \eng{so that} + présent) ; \eng{in order not to miss} (négatif, même sujet, \eng{in order not to} + BV).
\item Conséquence : \eng{so you know} (conjonction de coordination) ; \eng{so... that} (\eng{So many people came that...}) ; \eng{such... that} (\eng{such a busy day that...}) ; \eng{Therefore} (adverbe de liaison formel).
\item Obligation intérieure (règle ressentie) : \eng{must book}. Obligation extérieure (circonstance) : \eng{have to book} (implicite dans le contexte). Absence d'obligation : \eng{don't have to pay}. Interdiction (forme) : \eng{mustn't}.
\item \eng{prefer ... to} (préposition \eng{to} + nom/\eng{-ing}) ; \eng{would rather ... than} (base verbale + \eng{than}).
\item \eng{Let's pay}, \eng{Let's not forget}, \eng{Let me check} $\rightarrow$ toujours \textbf{base verbale} (infinitif sans \eng{to}).
\item Present perfect : \eng{have cooked} (expérience vie, antériorité vague) ; \eng{hadn't eaten} (past perfect, antériorité par rapport au passé \eng{was}). Past perfect : \eng{had already finished} (antériorité par rapport à \eng{arrived}) ; \eng{hadn't even eaten} (antériorité par rapport à \eng{was}).
\end{enumerate}
\end{exoresolu}

\courssub{Fiche méthode 1 : Exprimer le but (purpose)}

\begin{methode}[Comment exprimer le but en anglais]
\begin{enumerate}
\item \textbf{Identifier la relation logique} : la seconde action explique \emph{pourquoi} on fait la première. En français : « pour », « afin de », « dans le but de », « pour que ».
\item \textbf{Choisir la structure selon le sujet} :
\begin{itemize}
\item \textbf{Même sujet} : \cle{to + base verbale} (standard), \cle{in order to / so as to + base verbale} (plus formel, insistance sur le but).
\eng{She works hard to pass her exam.} « Elle travaille dur pour réussir son examen. »
\item \textbf{Sujets différents} : \cle{so that + sujet + can/could/will/would + base verbale}.
\eng{He spoke slowly so that everyone could understand.} « Il a parlé lentement pour que tout le monde puisse comprendre. »
\item \textbf{But négatif} (même sujet) : \cle{so as not to / in order not to + base verbale}. \textbf{Jamais} \eng{not to} seul en début de phrase ou après verbe de mouvement.
\eng{She left early so as not to miss the bus.} « Elle est partie tôt pour ne pas rater le bus. »
\end{itemize}
\item \textbf{Accorder le temps après \eng{so that}} : principale au présent/futur $\rightarrow$ \eng{can/will} ; principale au passé $\rightarrow$ \eng{could/would}.
\item \textbf{Test de vérification} : remplacer mentalement par « pour » / « afin que » en français. Si cela fonctionne, la structure est correcte.
\end{enumerate}
\end{methode}

\begin{exoresolu}[Transformation : exprimer le but]
\textbf{Énoncé.} Rewrite each sentence using the word in brackets without changing the meaning.
\begin{enumerate}
\item Ngozi studies English. She wants to become a journalist. (\eng{in order to})
\item Mr Tabi built a fence. He wanted to protect his cocoa farm from goats. (\eng{so that})
\item They whispered. They did not want to wake the baby. (\eng{so as not to})
\end{enumerate}
\tcblower
\textbf{Solution détaillée.}
\begin{enumerate}
\item Même sujet (\eng{Ngozi}), but positif $\rightarrow$ \eng{in order to + BV}.\\
\eng{Ngozi studies English in order to become a journalist.}\\
« Ngozi étudie l'anglais afin de devenir journaliste. »
\item Sujets différents (\eng{Mr Tabi} / protection ferme), principale au passé (\eng{built}) $\rightarrow$ \eng{so that + would/could}.\\
\eng{Mr Tabi built a fence so that his cocoa farm would be protected from goats.} (passif) ou \eng{Mr Tabi built a fence so that goats could not enter his cocoa farm.}\\
« M. Tabi a construit une clôture pour que les chèvres ne puissent pas entrer dans sa plantation de cacao. »
\item But négatif, même sujet $\rightarrow$ \eng{so as not to + BV}.\\
\eng{They whispered so as not to wake the baby.}\\
« Ils chuchotaient pour ne pas réveiller le bébé. »
\end{enumerate}
\textbf{Conclusion.} Même sujet $\rightarrow$ \eng{to / in order to / so as to} ; sujets différents $\rightarrow$ \eng{so that + modal} ; but négatif $\rightarrow$ \eng{so as not to / in order not to}.
\end{exoresolu}

\courssub{Fiche méthode 2 : Exprimer la conséquence (result)}

\begin{methode}[So, so... that, such... that, therefore]
\begin{enumerate}
\item \textbf{Repérer le lien} : la seconde proposition est la \emph{conséquence} de la première (« si bien que », « tellement... que », « donc », « par conséquent »).
\item \textbf{Choisir la structure} :
\begin{itemize}
\item \cle{so + adjectif/adverbe + that} : \eng{The test was so difficult that many students failed.} « Le test était si difficile que beaucoup d'élèves ont échoué. »
\item \cle{such (a/an) + (adjectif) + nom + that} : \eng{It was such a difficult test that many students failed.} « C'était un test si difficile que... »
\item \cle{so + quantificateur (many, few, much, little) + nom + that} : \eng{So many people came that we ran out of chairs.}
\item \cle{so} (conjonction de coordination) = « donc », « si bien que » : \eng{It rained heavily, so the match was postponed.}
\item \cle{therefore / consequently / as a result} (adverbes de liaison, registre formel, souvent en début de phrase suivis d'une virgule) : \eng{The road was flooded. Therefore, the bendskins stopped working.}
\end{itemize}
\item \textbf{Attention à l'ordre des mots} : \eng{so + adjectif} mais \cle{such + (a/an) + nom}. Comparer : \eng{so tall a boy} (littéraire, rare) vs \eng{such a tall boy} (standard).
\item \textbf{Vérification} : remplacer par « tellement... que » / « si bien que ». Si le sens tient, la structure est bonne.
\end{enumerate}
\end{methode}

\begin{exoresolu}[Complétion : so / such / such a/an]
\textbf{Énoncé.} Complete the sentences with \emph{so}, \emph{such} or \emph{such a/an}.
\begin{enumerate}
\item The market in Mokolo was \dots\ crowded \dots\ we could not find our aunt.
\item It was \dots\ hot day \dots\ we decided to stay at home.
\item The Principal spoke \dots\ clearly \dots\ every student understood the new rules.
\item There were \dots\ many students \dots\ the hall was full.
\end{enumerate}
\tcblower
\textbf{Solution détaillée.}
\begin{enumerate}
\item \eng{crowded} = adjectif $\rightarrow$ \eng{so + adjectif + that}.\\
\eng{The market in Mokolo was so crowded that we could not find our aunt.}\\
« Le marché de Mokolo était tellement bondé que nous n'avons pas pu retrouver notre tante. »
\item \eng{hot day} = adjectif + nom singulier comptable $\rightarrow$ \eng{such a + adjectif + nom + that}.\\
\eng{It was such a hot day that we decided to stay at home.}\\
« Il faisait une journée si chaude que nous avons décidé de rester à la maison. »
\item \eng{clearly} = adverbe $\rightarrow$ \eng{so + adverbe + that}.\\
\eng{The Principal spoke so clearly that every student understood the new rules.}\\
« Le Proviseur parlait si clairement que chaque élève a compris les nouvelles règles. »
\item \eng{many} = quantificateur + nom pluriel $\rightarrow$ \eng{so + many + nom + that}.\\
\eng{There were so many students that the hall was full.}\\
« Il y avait tellement d'élèves que la salle était pleine. »
\end{enumerate}
\textbf{Conclusion.} Devant adjectif/adverbe/quantificateur : \eng{so}. Devant nom (avec/sans adjectif) : \eng{such (a/an)}.
\end{exoresolu}

\courssub{Fiche méthode 3 : Obligation, interdiction et préférence}

\begin{methode}[Must, have to, mustn't, don't have to, prefer, would rather]
\begin{enumerate}
\item \textbf{Obligation} :
\begin{itemize}
\item \cle{must} : obligation intérieure, règle personnelle, ordre émanant du locuteur, déduction forte. \eng{You must wear a uniform.} (Règle de l'école / Mon avis).
\item \cle{have to} : obligation extérieure, contrainte factuelle, loi, circonstance. \eng{I have to wear a uniform.} (C'est la règle, je n'ai pas le choix).
\item \textbf{Au passé et au futur} : seul \eng{have to} se conjugue (\eng{had to}, \eng{will have to}). \eng{Must} n'a pas de passé ni de futur.
\end{itemize}
\item \textbf{Distinction capitale (négation)} :
\begin{itemize}
\item \cle{mustn't} = \textbf{interdiction} (« il est interdit de », « il ne faut pas »). \eng{You mustn't smoke here.} « Il est interdit de fumer ici. »
\item \cle{don't have to} = \textbf{absence d'obligation} (« ce n'est pas la peine de », « il n'est pas nécessaire de »). \eng{You don't have to come early.} « Tu n'as pas besoin de venir tôt / Ce n'est pas la peine. »
\end{itemize}
\item \textbf{Préférence} :
\begin{itemize}
\item \cle{prefer + nom / -ing + to + nom / -ing} (comparaison générale). \eng{I prefer reading to watching television.} « Je préfère lire à regarder la télé. »
\item \cle{would rather ('d rather) + base verbale + than} (préférence spécifique, immédiate). \eng{I would rather stay at home than go out.} « Je préférerais rester à la maison plutôt que sortir. »
\item \cle{would prefer + to + base verbale} (alternative à \eng{would rather}). \eng{I would prefer to walk.} « Je préférerais marcher. »
\end{itemize}
\item \textbf{Vérification des formes} : après \eng{must/mustn't/would rather/let's} $\rightarrow$ \textbf{base verbale sans \eng{to}} ; après \eng{have to/would prefer} $\rightarrow$ \eng{to} + BV ; après \eng{prefer... to} $\rightarrow$ \eng{to} = préposition $\rightarrow$ nom ou \eng{-ing}.
\end{enumerate}
\end{methode}

\begin{exoresolu}[QCM grammatical : obligation et préférence]
\textbf{Énoncé.} Choose the correct option.
\begin{enumerate}
\item You \dots\ wear a uniform at Government High School; it is compulsory. (\eng{must} / \eng{mustn't} / \eng{don't have to})
\item Today is a holiday, so we \dots\ go to school. (\eng{mustn't} / \eng{don't have to} / \eng{had to})
\item I prefer fufu corn \dots\ rice. (\eng{than} / \eng{to} / \eng{from})
\item She would rather \dots\ a taxi than walk in the rain. (\eng{take} / \eng{to take} / \eng{taking})
\end{enumerate}
\tcblower
\textbf{Solution détaillée.}
\begin{enumerate}
\item \textbf{Compulsory} = obligatoire (règle) $\rightarrow$ \eng{must}.\\
\eng{You must wear a uniform at Government High School; it is compulsory.}\\
« Tu dois porter l'uniforme au lycée public ; c'est obligatoire. »
\item Jour férié $\rightarrow$ \textbf{absence d'obligation} $\rightarrow$ \eng{don't have to}. \eng{Mustn't} signifierait « il est interdit d'aller à l'école » (non-sens).\\
\eng{Today is a holiday, so we don't have to go to school.}\\
« Aujourd'hui c'est férié, donc nous n'avons pas besoin d'aller à l'école. »
\item Avec \eng{prefer}, la préposition est \eng{to} (jamais \eng{than}).\\
\eng{I prefer fufu corn to rice.}\\
« Je préfère le fufu de maïs au riz. »
\item Après \eng{would rather}, \textbf{base verbale sans \eng{to}}.\\
\eng{She would rather take a taxi than walk in the rain.}\\
« Elle préférerait prendre un taxi plutôt que de marcher sous la pluie. »
\end{enumerate}
\textbf{Conclusion.} \eng{Mustn't} = interdit ; \eng{don't have to} = pas obligé ; \eng{prefer... to} ; \eng{would rather + BV}.
\end{exoresolu}

\courssub{Fiche méthode 4 : Let's / Let's not + base verbale}

\begin{methode}[Faire une proposition avec let's]
\begin{enumerate}
\item \textbf{Fonction} : \cle{let's} (= \eng{let us}) sert à proposer une activité incluant le locuteur : « Et si nous... ? », « Allons... ».
\item \textbf{Forme affirmative} : \cle{let's + base verbale} (sans \eng{to}). \eng{Let's visit Grandma this weekend.} « Rendons visite à Grand-mère ce week-end. »
\item \textbf{Forme négative} : \cle{let's not + base verbale} (jamais \eng{let's don't}). \eng{Let's not argue about money.} « Ne nous disputons pas pour l'argent. »
\item \textbf{Réponses types} : \eng{Yes, let's!} / \eng{That's a good idea.} / \eng{I'd rather not.} / \eng{Let's not.}
\item \textbf{Autres personnes} : \eng{let's} ne s'utilise qu'à la 1\textsuperscript{re} personne du pluriel. Pour les autres : \cle{let + pronom complément (me, him, her, us, them) + base verbale}. \eng{Let him speak.} « Laisse-le parler. » / \eng{Let me check.} « Laisse-moi vérifier. »
\end{enumerate}
\end{methode}

\begin{exoresolu}[Mise en situation : proposer, accepter, refuser]
\textbf{Énoncé.} Two friends, Amina and Peter, are planning their Saturday in Yaoundé. Complete their dialogue with \emph{let's}, \emph{let's not} or \emph{let + pronoun} and the verbs in brackets.

Amina: The weather is fine. (1) \dots\ (go) to Mvog-Betsi Zoo!\\
Peter: Good idea, but (2) \dots\ (take) a bendskin; it is too expensive. We can walk.\\
Amina: All right. And (3) \dots\ (buy) some groundnuts on the way.\\
Peter: My little brother wants to come too. (4) \dots\ (he / join) us.\\
Amina: Of course! But (5) \dots\ (not / be) late.
\tcblower
\textbf{Solution détaillée.}
\begin{enumerate}
\item Proposition affirmative (nous) $\rightarrow$ \eng{Let's go to Mvog-Betsi Zoo!} « Allons au zoo de Mvog-Betsi ! »
\item Proposition négative (nous) $\rightarrow$ \eng{let's not + BV} : \eng{...but let's not take a bendskin; it is too expensive.} « ...mais ne prenons pas de bendskin ; c'est trop cher. »
\item Nouvelle proposition affirmative (nous) $\rightarrow$ \eng{And let's buy some groundnuts on the way.} « Et achetons des arachides en chemin. »
\item Troisième personne (lui) $\rightarrow$ \eng{let + him + BV} : \eng{Let him join us.} « Laisse-le se joindre à nous. »
\item Proposition négative (nous) $\rightarrow$ \eng{But let's not be late.} « Mais ne soyons pas en retard. »
\end{enumerate}
\textbf{Conclusion.} \eng{Let's / Let's not + BV} pour « nous » ; \eng{let + me/him/her/them + BV} pour les autres personnes.
\end{exoresolu}

\courssub{Fiche méthode 5 : Present perfect et past perfect}

\begin{methode}[Choisir entre present perfect et past perfect]
\begin{enumerate}
\item \textbf{Present Perfect} : \cle{have/has + participe passé (V3)}.
\textbf{Emplois majeurs} :
\begin{itemize}
\item Action passée avec \textbf{résultat présent} : \eng{I have lost my keys} (= je ne les ai plus \textbf{maintenant}).
\item \textbf{Expérience de vie} (souvent avec \eng{ever, never, already, just, yet}) : \eng{Have you ever visited Bamenda?}
\item Action \textbf{commencée dans le passé et qui continue} (avec \eng{for} + durée / \eng{since} + point de départ) : \eng{She has lived in Douala since 2015.}
\end{itemize}
\textbf{Interdiction formelle} : JAMAIS avec un repère de temps passé précis (\eng{yesterday, last year, in 2019, two days ago}) $\rightarrow$ \textbf{Preterit}.
\item \textbf{Past Perfect} : \cle{had + participe passé (V3)}.
\textbf{Emploi} : exprime l'\textbf{antériorité} dans le passé. Une action accomplie \emph{avant} une autre action passée (souvent au preterit).
\eng{When we arrived, the match had already started.}
\item \textbf{Démarche de résolution (frise chronologique)} :
\begin{enumerate}
\item Repérer les deux événements passés.
\item Les placer sur une ligne du temps.
\item L'événement \textbf{le plus ancien} (antérieur) $\rightarrow$ \textbf{Past Perfect}.
\item L'événement \textbf{le plus récent} (point de repère) $\rightarrow$ \textbf{Preterit}.
\end{enumerate}
\item \textbf{Marqueurs usuels} : \eng{before, after, when, by the time, already, just, no sooner... than, hardly... when}.
\end{enumerate}
\end{methode}

\begin{popfigure}
\begin{tikzpicture}[
    event/.style={circle, draw=popblue, fill=popblueL, thick, minimum width=1.5cm, align=center},
    arrow/.style={->, >=Stealth, thick, popdark},
    time/.style={popmuted, font=\small}
]
% Timeline
\draw[popline, thick] (0,0) -- (12,0);
% Events
\node[event, align=center] (E1) at (2, 1.5){Past Perfect\\(had + V3)\\\textbf{Action antérieure}};
\node[event, align=center] (E2) at (8, 1.5){Preterit\\(V2)\\\textbf{Point de repère passé}};
% Arrows
\draw[arrow] (E1) -- (E2) node[midway, above, font=\small, popdark] {before / when / by the time};
% Present Perfect reference
\node[event, fill=popgreenL, draw=popgreen, align=center] (E3) at (5, -2){Present Perfect\\(have/has + V3)\\\textbf{Pasé $\rightarrow$ Présent}};
\node[time] at (5, -3.5) {for / since / just / already / yet / ever / never};
\draw[arrow, popgreen] (E3.north) -- ++(0, 0.5) -| (6,0) node[pos=0.7, right, font=\small, popgreen] {Now};
\end{tikzpicture}
\legende{Frise chronologique : Past Perfect (antériorité passée) vs Present Perfect (lien passé-présent)}
\end{popfigure}

\begin{exoresolu}[Conjugaison : present perfect, past perfect ou preterit]
\textbf{Énoncé.} Put the verbs in brackets into the correct tense (present perfect, past perfect or simple past).
\begin{enumerate}
\item My uncle (live) \dots\ in Buea for twenty years, and he is still there.
\item When the students (arrive) \dots\ , the examination (already / begin) \dots\ .
\item I (never / see) \dots\ the Waza National Park, but I would love to visit it.
\item After she (finish) \dots\ her homework, she (go) \dots\ to bed.
\end{enumerate}
\tcblower
\textbf{Solution détaillée.}
\begin{enumerate}
\item Action commencée il y a 20 ans et qui \textbf{continue encore} (\eng{he is still there}) avec \eng{for twenty years} $\rightarrow$ \textbf{Present Perfect}.\\
\eng{My uncle has lived in Buea for twenty years, and he is still there.}\\
« Mon oncle vit à Buea depuis vingt ans, et il y est toujours. »
\item Deux actions passées : l'examen a commencé \textbf{avant} l'arrivée des élèves $\rightarrow$ antérieur = \textbf{Past Perfect}, récent = \textbf{Preterit}.\\
\eng{When the students arrived, the examination had already begun.}\\
« Quand les élèves sont arrivés, l'examen avait déjà commencé. »
\item Expérience de vie avec \eng{never}, sans date précise $\rightarrow$ \textbf{Present Perfect}.\\
\eng{I have never seen the Waza National Park, but I would love to visit it.}\\
« Je n'ai jamais vu le parc national de Waza, mais j'aimerais beaucoup le visiter. »
\item Avec \eng{after}, l'action terminée en premier = \textbf{Past Perfect} ; la conséquence = \textbf{Preterit}.\\
\eng{After she had finished her homework, she went to bed.}\\
« Après avoir fini ses devoirs, elle s'est couchée. »
\end{enumerate}
\textbf{Conclusion.} Résultat présent, expérience, \eng{for/since} $\rightarrow$ Present Perfect. Antériorité par rapport à un passé $\rightarrow$ Past Perfect. Repère de temps passé précis ($\rightarrow$ Preterit).
\end{exoresolu}

\courssub{Entraînement et production}

\begin{exercice}[Exercice 1 : Choix multiple et justification]
\textbf{Énoncé.} Choose the correct form and justify briefly in French.
\begin{enumerate}
\item If you want to pass, you \_\_\_\_\_ work harder. (\eng{must} / \eng{have to}) \textit{Justification :}
\item You \_\_\_\_\_ drive on the right in the UK. (\eng{mustn't} / \eng{don't have to}) \textit{Justification :}
\item \_\_\_\_\_ go to the cinema tonight? (\eng{Let's} / \eng{Let us}) \textit{Justification :}
\item She spoke very \_\_\_\_\_ \_\_\_\_\_ everyone heard her. (\eng{so / such} \dots\ \eng{that}) \textit{Justification :}
\item By the time I got to the station, the train \_\_\_\_\_. (\eng{left} / \eng{had left}) \textit{Justification :}
\end{enumerate}
\end{exercice}

\begin{corrige}[Exercice 1]
\begin{enumerate}
\item \eng{must} (obligation intérieure, conseil fort) ou \eng{have to} (obligation objective). Les deux possibles, nuance différente.
\item \eng{mustn't} (interdiction : au UK on conduit à gauche, donc interdiction de rouler à droite).
\item \eng{Let's} (proposition usuelle). \eng{Let us} est formel/insistant (« Allow us to... »).
\item \eng{so loudly that} (\eng{loudly} = adverbe $\rightarrow$ \eng{so}).
\item \eng{had left} (antériorité : le train est parti \emph{avant} mon arrivée $\rightarrow$ Past Perfect).
\end{enumerate}
\end{corrige}

\begin{exercice}[Exercice 2 : Transformation de phrases]
\textbf{Énoncé.} Rewrite the sentences using the word in bold without changing the meaning.
\begin{enumerate}
\item It is forbidden to use phones in class. (\eng{MUSTN'T}) $\rightarrow$ Students \dots
\item It isn't necessary to bring food. (\eng{HAVE}) $\rightarrow$ You \dots
\item "Shall we go for a walk?" he said. (\eng{LET'S}) $\rightarrow$ He suggested \dots
\item The coffee was too hot to drink. (\eng{SO}) $\rightarrow$ The coffee was \dots
\item He drove fast because he wanted to arrive early. (\eng{SO THAT}) $\rightarrow$ He drove fast \dots
\item I last saw him in 2010. (\eng{FOR}) $\rightarrow$ I \dots\ 2010.
\end{enumerate}
\end{exercice}

\begin{corrige}[Exercice 2]
\begin{enumerate}
\item Students \eng{mustn't use} phones in class.
\item You \eng{don't have to bring} food.
\item He suggested \eng{let's go} for a walk. (ou \eng{they should go} / \eng{going} mais consigne \eng{LET'S}).
\item The coffee was \eng{so hot that I couldn't drink it} / \eng{so hot that it couldn't be drunk}.
\item He drove fast \eng{so that he would arrive early} / \eng{so that he could arrive early}.
\item I \eng{haven't seen him for} 2010. (Incorrect : \eng{for} attend une durée. Correct : \eng{I haven't seen him since 2010.} Si contrainte \eng{FOR} imposée : \eng{I haven't seen him for fourteen years} (approx). \textbf{Note pédagogique} : \eng{since} pour point de départ, \eng{for} pour durée).
\end{enumerate}
\end{corrige}

\begin{exercice}[Exercice 3 : Production écrite guidée]
\textbf{Énoncé.} Write a short paragraph (80--100 words) describing a family event you organized or attended. You \textbf{must} include at least:
\begin{itemize}
\item One \textbf{purpose} structure (\eng{to / in order to / so that}).
\item One \textbf{result} structure (\eng{so/such... that / therefore}).
\item One \textbf{obligation/preference} structure (\eng{must / have to / would rather / prefer}).
\item \textbf{Let's / Let's not} (reported or direct speech).
\item One \textbf{Present Perfect} and one \textbf{Past Perfect} verb form.
\end{itemize}
\textit{Context suggestion:} A birthday party in Douala, a wedding in Bamenda, a family meeting in the village.
\end{exercice}

\begin{corrige}[Exercice 3 -- Proposition de rédaction]
\textbf{Modèle :}
Last December, my family \textbf{had organized} (Past Perfect) a big reunion in our village near Bafoussam. We \textbf{had to} (obligation extérieure) book the community hall months in advance \textbf{in order to} (but) secure the date. My father \textbf{would rather} (préférence) cook the traditional \emph{ndolé} himself \textbf{than} hire a caterer. On the day, \textbf{so} many relatives came \textbf{that} (conséquence) we ran out of chairs. \textbf{Therefore}, we \textbf{had to} fetch more from the neighbour. My uncle said, "\textbf{Let's not} worry, the food is delicious!" I \textbf{have never seen} (Present Perfect -- expérience) such a joyful atmosphere. Everyone danced until dawn.
\end{corrige}

\begin{attention}[Les 5 erreurs qui coûtent des points au Bac]
\begin{enumerate}
\item \textbf{Confondre \eng{mustn't} et \eng{don't have to}.} \eng{You mustn't be late} = « il est interdit d'être en retard » ; \eng{You don't have to come} = « tu n'es pas obligé de venir ». Traduire « ne pas devoir » par \eng{mustn't} est une faute de sens éliminatoire.
\item \textbf{Mettre \eng{to} après \eng{let's} ou \eng{would rather}.} On écrit \eng{Let's go}, \eng{I'd rather stay}, jamais \eng{let's to go} ni \eng{I'd rather to stay}. La négation est \eng{let's not}, jamais \eng{let's don't}.
\item \textbf{Employer le present perfect avec un repère passé.} \eng{I have seen him yesterday} est incorrect : avec \eng{yesterday, last week, ago}, le preterit s'impose : \eng{I saw him yesterday.}
\item \textbf{Calquer « pour que » par \eng{for to}.} Le but s'exprime par \eng{to / in order to / so as to + BV} ou \eng{so that + proposition}. \eng{He came for to help me} est un calque du français à bannir : \eng{He came to help me.}
\item \textbf{Inverser \eng{so} et \eng{such}.} Devant un adjectif ou un adverbe : \eng{so} (\eng{so tired}). Devant un nom : \eng{such (a/an)} (\eng{such a tiring journey}). Ne jamais écrire \eng{a so tall boy} : dire \eng{such a tall boy}.
\end{enumerate}
\end{attention}

\newpage

\coursec{Exercices}

\courssub{Je vérifie mes connaissances}

\begin{exercice}[QCM : la bonne structure]
Entoure la bonne réponse dans chaque phrase.
\begin{enumerate}
\item My father works hard \textbf{to / for to / so that} pay our school fees.
\item You \textbf{must / mustn't / don't have to} smoke in the hospital. It is forbidden.
\item The soup was \textbf{so / such / too} hot that nobody could eat it.
\item It was \textbf{so / such / too} a long journey that we arrived at midnight.
\item \textbf{Let's / Let we / Let's to} visit grandmother this weekend.
\item I \textbf{prefer / would rather / had better} stay at home tonight. I am tired.
\item You \textbf{mustn't / don't have to / must} bring your own pen; the school provides one.
\item She left early \textbf{in order to / so that / for} she could catch the first bus to Bamenda.
\end{enumerate}
\end{exercice}

\begin{exercice}[Vrai ou faux ?]
Dis si chaque affirmation est \textbf{vraie} ou \textbf{fausse}, puis justifie ta réponse en citant la règle.
\begin{enumerate}
\item After \eng{let's}, we use the infinitive with \eng{to}.
\item \eng{Mustn't} and \eng{don't have to} have exactly the same meaning.
\item \eng{So} is used before an adjective, and \eng{such} before “a/an + adjective + noun”.
\item The past perfect is formed with \eng{had} + past participle.
\item \eng{Would rather} is followed by the \eng{-ing} form of the verb.
\item The present perfect can be used with \eng{yesterday}.
\end{enumerate}
\end{exercice}

\begin{exercice}[Vocabulaire : les mots du but et de la conséquence]
Complète chaque phrase avec le mot ou l'expression qui convient, choisi dans la liste : \eng{in order to — so that — therefore — as a result — would rather — let's not}.
\begin{enumerate}
\item Ngozi saved money \ldots\ldots\ldots\ldots buy a sewing machine.
\item It rained heavily all night; \ldots\ldots\ldots\ldots, the football match was cancelled.
\item Take an umbrella \ldots\ldots\ldots\ldots you don't get wet.
\item \ldots\ldots\ldots\ldots argue about money; let's find a solution together.
\item I \ldots\ldots\ldots\ldots eat fufu than rice this evening.
\item The driver was careless and, \ldots\ldots\ldots\ldots, the bendskin had an accident.
\end{enumerate}
\end{exercice}

\begin{exercice}[Conjugaison : present perfect ou past perfect ?]
Mets le verbe entre parenthèses au \textbf{present perfect} ou au \textbf{past perfect}.
\begin{enumerate}
\item I \ldots\ldots\ldots\ldots (never / see) a white man before I went to Buea.
\item My aunt \ldots\ldots\ldots\ldots (just / arrive) from Douala.
\item When we reached the market, the traders \ldots\ldots\ldots\ldots (already / close) their shops.
\item They \ldots\ldots\ldots\ldots (live) in Yaoundé since 2015.
\item She could not enter the house because she \ldots\ldots\ldots\ldots (lose) her keys.
\item \ldots\ldots\ldots\ldots you \ldots\ldots\ldots\ldots (ever / taste) eru?
\end{enumerate}
\end{exercice}

\courssub{Je m'entraîne}

\begin{exercice}[Traduction : français vers anglais]
Traduis les phrases suivantes en anglais en utilisant la structure indiquée entre parenthèses.
\begin{enumerate}
\item Mon frère étudie beaucoup afin de réussir son examen. (\eng{in order to})
\item Le repas était si bon que nous avons tout fini. (\eng{so\ldots that})
\item Je préférerais rester au village plutôt que d'aller en ville. (\eng{would rather})
\item N'allons pas au marché aujourd'hui ; il pleut. (\eng{let's not})
\item Quand je suis arrivé à la maison, ma mère avait déjà préparé le dîner. (\eng{past perfect})
\end{enumerate}
\end{exercice}

\begin{exercice}[Transformation de phrases]
Réécris chaque phrase en utilisant le mot ou l'expression donné, sans changer le sens.
\begin{enumerate}
\item The boy was very tired. He could not walk any further. (\eng{so\ldots that})
\item It was a very difficult exercise. Nobody could do it. (\eng{such\ldots that})
\item She went to the library because she wanted to borrow a book. (\eng{in order to})
\item I prefer drinking water to drinking beer. (\eng{would rather})
\item We should not waste our time. (\eng{let's not})
\item You are not allowed to use your phone during the exam. (\eng{mustn't})
\end{enumerate}
\end{exercice}

\begin{exercice}[Texte à trous : A family meeting]
Complète le texte suivant en mettant les verbes entre parenthèses au \textbf{prétérit}, au \textbf{present perfect} ou au \textbf{past perfect}, selon le sens.
\begin{textebox}[A family meeting]
Last Saturday, our family (1. hold) \ldots\ldots\ldots\ldots a meeting to discuss the school fees. My father (2. say) \ldots\ldots\ldots\ldots that prices (3. rise) \ldots\ldots\ldots\ldots a lot this year, and that we (4. have) \ldots\ldots\ldots\ldots to save money. My elder brother, who (5. finish) \ldots\ldots\ldots\ldots secondary school two years before, (6. promise) \ldots\ldots\ldots\ldots to help. I (7. never / see) \ldots\ldots\ldots\ldots my father so serious. When the meeting (8. end) \ldots\ldots\ldots\ldots, we all (9. agree) \ldots\ldots\ldots\ldots to work harder. Since that day, our family (10. become) \ldots\ldots\ldots\ldots more united.
\end{textebox}
\end{exercice}

\begin{exercice}[Correction d'erreurs]
Chaque phrase contient \textbf{une erreur} typique des élèves francophones. Souligne l'erreur et réécris la phrase correctement.
\begin{enumerate}
\item I went to Douala for to visit my uncle.
\item Let's to go to the stadium this afternoon.
\item You mustn't to come tomorrow; it is a holiday.
\item It was so a good film that we watched it twice.
\item I have seen him yesterday at the market.
\item When I arrived, she already cooked the food.
\item I would rather to stay at home tonight.
\item He works hardly in order to succeed.
\end{enumerate}
\end{exercice}

\courssub{Je me prépare au Bac}

\begin{exercice}[Type Bac — Compréhension et langue]
Lis le texte suivant, puis réponds aux questions.

\begin{textebox}[A new start for the Ekane family]
The Ekane family had lived in a small village near Kumba for many years. Life there was peaceful, but the parents wanted their children to have a better education, so they decided to move to Douala. Before leaving, they sold their cocoa farm in order to pay for the move and the children's new school.

The first months in the city were so difficult that Mrs Ekane almost lost hope. She had never lived in such a noisy place, and everything was expensive. “Let's not give up,” her husband kept saying. “We must be patient.” Little by little, things improved. Mr Ekane found a job as a driver, and the children, who had never spoken English at school before, made rapid progress.

Today, the family has been in Douala for three years. The children have made new friends, and Mrs Ekane has opened a small restaurant. “I would rather live in the village,” she sometimes admits with a smile, “but I know we made the right choice for our children.”
\end{textebox}

\textbf{A. Comprehension questions} — Answer in complete sentences.
\begin{enumerate}
\item Why did the Ekane family leave their village? (2 marks)
\item What did they do in order to pay for the move? (2 marks)
\item Why did Mrs Ekane almost lose hope at the beginning? Give two reasons. (2 marks)
\item What advice did Mr Ekane give his wife? (2 marks)
\end{enumerate}

\textbf{B. Language work}
\begin{enumerate}
\item Find in the text: one sentence expressing \textbf{purpose}, one expressing \textbf{result}, and one expressing \textbf{preference}. Copy them. (3 marks)
\item Identify two verbs in the \textbf{past perfect} and two in the \textbf{present perfect} in the text. (2 marks)
\item Rewrite: “Life in the city was very difficult. Mrs Ekane almost lost hope.” Use \eng{so\ldots that}. (1 mark)
\item Rewrite: “We should be patient.” Use \eng{must}. (1 mark)
\end{enumerate}
\end{exercice}

\begin{exercice}[Type Bac — Traduction]
Traduis en anglais le passage suivant. Fais attention aux temps (past perfect, present perfect) et aux structures de but et de préférence.

\begin{textebox}[Texte à traduire]
Ma sœur avait toujours rêvé de devenir infirmière. Après avoir obtenu son baccalauréat, elle a travaillé pendant deux ans afin d'économiser de l'argent pour ses études. Le travail était si fatigant qu'elle rentrait épuisée tous les soirs, mais elle n'a jamais abandonné. « Je préférerais souffrir maintenant et réussir demain », disait-elle souvent. Aujourd'hui, elle a terminé sa formation et elle travaille dans un grand hôpital de Yaoundé. Toute la famille est fière d'elle.
\end{textebox}
\end{exercice}

\begin{exercice}[Type Bac — Production écrite]
\textbf{Topic:} \emph{How my family prepares for Christmas.}

Rédige une composition de \textbf{150 à 200 mots} (environ 15 lignes) dans laquelle tu décris comment ta famille se prépare pour Noël : les achats, la cuisine, les voyages, les retrouvailles.

\textbf{Consignes obligatoires :} ta composition doit contenir \textbf{au moins six structures} de la leçon, choisies parmi les suivantes :
\begin{itemize}
\item une expression de but : \eng{to}, \eng{in order to} ou \eng{so that} ;
\item une expression de conséquence : \eng{so\ldots that} ou \eng{such\ldots that} ;
\item une expression d'obligation ou d'interdiction : \eng{must} ou \eng{mustn't} ;
\item une expression de préférence : \eng{would rather} ou \eng{prefer} ;
\item une suggestion avec \eng{let's} ou \eng{let's not} ;
\item un verbe au \textbf{present perfect} et un verbe au \textbf{past perfect}.
\end{itemize}

Souligne chaque structure utilisée dans ta copie. Soigne l'organisation de ton texte (introduction, développement, conclusion) et l'orthographe.
\end{exercice}

\newpage

\coursec{Corrigés des exercices}

\begin{corrige}[Exercice 1 — QCM : la bonne structure]
\begin{enumerate}
\item \textbf{to} — \eng{My father works hard to pay our school fees.} Le but s'exprime avec \eng{to} + base verbale ; \eng{for to} n'existe pas en anglais standard.
\item \textbf{mustn't} — \eng{You mustn't smoke in the hospital.} Il s'agit d'une interdiction (« It is forbidden »). \eng{Don't have to} signifierait « ce n'est pas obligatoire ».
\item \textbf{so} — \eng{The soup was so hot that nobody could eat it.} \eng{So} + adjectif (\eng{hot}) + \eng{that}.
\item \textbf{such} — \eng{It was such a long journey that we arrived at midnight.} \eng{Such} + a/an + adjectif + nom (\eng{a long journey}).
\item \textbf{Let's} — \eng{Let's visit grandmother this weekend.} \eng{Let's} + base verbale, sans \eng{to} ; \eng{let we} est impossible.
\item \textbf{would rather} — \eng{I would rather stay at home tonight.} Expression de préférence suivie de la base verbale.
\item \textbf{don't have to} — \eng{You don't have to bring your own pen; the school provides one.} Absence d'obligation (ce n'est pas une interdiction).
\item \textbf{so that} — \eng{She left early so that she could catch the first bus to Bamenda.} \eng{So that} + sujet + verbe conjugué ; \eng{in order to} est suivi d'une base verbale, pas d'une proposition.
\end{enumerate}
\end{corrige}

\begin{corrige}[Exercice 2 — Vrai ou faux ?]
\begin{enumerate}
\item \textbf{Faux.} Après \eng{let's}, on utilise la \emph{base verbale} (infinitif sans \eng{to}) : \eng{Let's go.} et non *\eng{Let's to go.}
\item \textbf{Faux.} \eng{Mustn't} exprime l'\emph{interdiction} (« il est interdit de »), tandis que \eng{don't have to} exprime l'\emph{absence d'obligation} (« ce n'est pas la peine de »).
\item \textbf{Vrai.} \eng{So} + adjectif (\eng{so hot}) ; \eng{such} + a/an + adjectif + nom (\eng{such a hot day}).
\item \textbf{Vrai.} Past perfect = \eng{had} + participe passé : \eng{She had left.} (« elle était partie »).
\item \textbf{Faux.} \eng{Would rather} est suivi de la \emph{base verbale} : \eng{I would rather stay.} et non *\eng{I would rather staying.}
\item \textbf{Faux.} Le present perfect est incompatible avec un repère de temps passé et terminé comme \eng{yesterday} ; on utilise le prétérit : \eng{I saw him yesterday.}
\end{enumerate}
\end{corrige}

\begin{corrige}[Exercice 3 — Vocabulaire : les mots du but et de la conséquence]
\begin{enumerate}
\item \eng{Ngozi saved money \textbf{in order to} buy a sewing machine.} (but + base verbale)
\item \eng{It rained heavily all night; \textbf{therefore}, the football match was cancelled.} (conséquence logique, « par conséquent » ; noter la virgule après \eng{therefore})
\item \eng{Take an umbrella \textbf{so that} you don't get wet.} (but + proposition avec sujet et verbe)
\item \eng{\textbf{Let's not} argue about money; let's find a solution together.} (suggestion négative + base verbale)
\item \eng{I \textbf{would rather} eat fufu than rice this evening.} (préférence ; noter la construction \eng{would rather\ldots than})
\item \eng{The driver was careless and, \textbf{as a result}, the bendskin had an accident.} (conséquence, « en conséquence »)
\end{enumerate}
\end{corrige}

\begin{corrige}[Exercice 4 — Conjugaison : present perfect ou past perfect ?]
\begin{enumerate}
\item \eng{I \textbf{had never seen} a white man before I went to Buea.} Past perfect : action antérieure à un moment passé (\eng{before I went}).
\item \eng{My aunt \textbf{has just arrived} from Douala.} Present perfect avec \eng{just} : action récente liée au présent.
\item \eng{When we reached the market, the traders \textbf{had already closed} their shops.} Past perfect : la fermeture est antérieure à notre arrivée.
\item \eng{They \textbf{have lived} in Yaoundé since 2015.} Present perfect avec \eng{since} : action commencée dans le passé et qui continue jusqu'à aujourd'hui.
\item \eng{She could not enter the house because she \textbf{had lost} her keys.} Past perfect : la perte des clés précède le moment où elle ne peut pas entrer.
\item \eng{\textbf{Have} you \textbf{ever tasted} eru?} Present perfect avec \eng{ever} : expérience de vie, sans date précise.
\end{enumerate}
\textbf{Rappel :} present perfect = \eng{have/has} + participe passé ; past perfect = \eng{had} + participe passé. Attention aux participes passés irréguliers : \eng{see–saw–seen}, \eng{lose–lost–lost}, \eng{go–went–gone}.
\end{corrige}

\begin{corrige}[Exercice 5 — Traduction : français vers anglais]
\begin{enumerate}
\item \eng{My brother studies hard in order to pass his exam.} (« afin de » + verbe = \eng{in order to} + base verbale.)
\item \eng{The meal was so good that we finished everything.} (\eng{so} + adjectif + \eng{that} ; prétérit \eng{finished}.)
\item \eng{I would rather stay in the village than go to town.} (Base verbale après \eng{would rather} et après \eng{than}, sans \eng{to}.)
\item \eng{Let's not go to the market today; it is raining.} (\eng{Let's not} + base verbale ; pas de \eng{to}.)
\item \eng{When I arrived home, my mother had already prepared dinner.} (Past perfect \eng{had prepared} pour l'action antérieure ; \eng{already} se place entre \eng{had} et le participe passé.)
\end{enumerate}
\end{corrige}

\begin{corrige}[Exercice 6 — Transformation de phrases]
\begin{enumerate}
\item \eng{The boy was \textbf{so} tired \textbf{that} he could not walk any further.} (\eng{so} + adjectif + \eng{that}.)
\item \eng{It was \textbf{such} a difficult exercise \textbf{that} nobody could do it.} (\eng{such} + a + adjectif + nom + \eng{that}.)
\item \eng{She went to the library \textbf{in order to} borrow a book.} (\eng{in order to} + base verbale remplace la subordonnée de but introduite par \eng{because she wanted to}.)
\item \eng{I \textbf{would rather} drink water \textbf{than} (drink) beer.} (Base verbale des deux côtés de \eng{than}.)
\item \eng{\textbf{Let's not} waste our time.} (Suggestion négative + base verbale.)
\item \eng{You \textbf{mustn't} use your phone during the exam.} (\eng{Mustn't} = interdiction, équivalent de \eng{not be allowed to}.)
\end{enumerate}
\end{corrige}

\begin{corrige}[Exercice 7 — Texte à trous : A family meeting]
\begin{enumerate}
\item \eng{held} — prétérit (\eng{last Saturday}, repère passé daté ; verbe irrégulier \eng{hold–held–held}).
\item \eng{said} — prétérit (\eng{say–said–said}).
\item \eng{had risen} — past perfect : la hausse des prix est antérieure au discours du père (\eng{rise–rose–risen}).
\item \eng{had} — prétérit (\eng{have to} au passé : « nous devions »).
\item \eng{had finished} — past perfect : action antérieure (\eng{two years before}).
\item \eng{promised} — prétérit, action du récit.
\item \eng{had never seen} — past perfect : expérience antérieure au moment du récit.
\item \eng{ended} — prétérit (\eng{when} + prétérit pour l'action-repère).
\item \eng{agreed} — prétérit.
\item \eng{has become} — present perfect avec \eng{since that day} : résultat qui dure jusqu'à aujourd'hui (\eng{become–became–become}).
\end{enumerate}
\textbf{Méthode :} repérer d'abord les marqueurs temporels (\eng{last Saturday}, \eng{before}, \eng{since that day}) ; ils indiquent le temps attendu. Le prétérit est le temps de base du récit ; le past perfect marque l'antériorité ; le present perfect relie le passé au présent.
\end{corrige}

\begin{corrige}[Exercice 8 — Correction d'erreurs]
\begin{enumerate}
\item Erreur : \eng{for to}. Correction : \eng{I went to Douala \textbf{to} visit my uncle.} (ou \eng{in order to visit}). Le but s'exprime avec \eng{to} seul ; \eng{for} est suivi d'un \emph{nom} (\eng{for a visit}).
\item Erreur : \eng{to go}. Correction : \eng{\textbf{Let's go} to the stadium this afternoon.} Base verbale après \eng{let's}.
\item Erreur : \eng{to come}. Correction : \eng{You \textbf{mustn't come} tomorrow; it is a holiday.} \eng{Must} est un modal : il est suivi de la base verbale sans \eng{to}.
\item Erreur : \eng{so a good film}. Correction : \eng{It was \textbf{such a good film} that we watched it twice.} Devant « a + adjectif + nom », on utilise \eng{such}, pas \eng{so}.
\item Erreur : \eng{have seen\ldots yesterday}. Correction : \eng{I \textbf{saw} him yesterday at the market.} Avec un repère passé terminé (\eng{yesterday}), le prétérit est obligatoire.
\item Erreur : \eng{already cooked}. Correction : \eng{When I arrived, she \textbf{had already cooked} the food.} L'action antérieure exige le past perfect.
\item Erreur : \eng{to stay}. Correction : \eng{I would rather \textbf{stay} at home tonight.} Base verbale après \eng{would rather}.
\item Erreur : \eng{hardly}. Correction : \eng{He works \textbf{hard} in order to succeed.} Faux ami : \eng{hard} = « dur » ; \eng{hardly} = « à peine ».
\end{enumerate}
\end{corrige}

\begin{corrige}[Exercice 9 — Type Bac : compréhension et langue]
\textbf{Barème indicatif :} A = 8 points (2 par question) ; B = 7 points (3 + 2 + 1 + 1).

\textbf{A. Comprehension questions — réponses modèles}
\begin{enumerate}
\item \eng{They left their village because the parents wanted their children to have a better education.} (2 pts : raison exacte + phrase complète)
\item \eng{They sold their cocoa farm in order to pay for the move and the children's new school.} (2 pts)
\item \eng{She almost lost hope because life in the city was very difficult and everything was expensive; she had also never lived in such a noisy place.} (2 pts : deux raisons exigées)
\item \eng{He advised her not to give up and told her that they had to be patient.} (2 pts)
\end{enumerate}

\textbf{B. Language work}
\begin{enumerate}
\item Purpose : \eng{“Before leaving, they sold their cocoa farm in order to pay for the move\ldots”} — Result : \eng{“The first months in the city were so difficult that Mrs Ekane almost lost hope.”} — Preference : \eng{“I would rather live in the village\ldots”} (1 pt par phrase correctement copiée)
\item Past perfect : \eng{had lived}, \eng{had never lived}, \eng{had never spoken} (deux au choix). Present perfect : \eng{has been}, \eng{have made}, \eng{has opened} (deux au choix). (0,5 pt par verbe)
\item \eng{Life in the city was so difficult that Mrs Ekane almost lost hope.} (1 pt)
\item \eng{We must be patient.} (1 pt)
\end{enumerate}

\textbf{Points de vigilance :} répondre par des phrases complètes (le correcteur sanctionne les réponses d'un seul mot) ; recopier les citations \emph{sans les modifier} ; ne pas confondre \eng{had lived} (past perfect) et \eng{has been} (present perfect).
\end{corrige}

\begin{corrige}[Exercice 10 — Type Bac : traduction]
\textbf{Barème indicatif :} 10 points (2 pts par segment correct, moins 0,5 pt par faute de grammaire ou de vocabulaire).

\textbf{Traduction modèle :}
\end{corrige}

\begin{textebox}[Model translation]
My sister had always dreamt of becoming a nurse. After passing her baccalauréat, she worked for two years in order to save money for her studies. The work was so tiring that she came home exhausted every evening, but she never gave up. “I would rather suffer now and succeed tomorrow,” she often said. Today, she has finished her training and she works in a big hospital in Yaoundé. The whole family is proud of her.
\end{textebox}

\begin{corrige}[Exercice 10 — Points de vigilance]
\begin{itemize}
\item « avait toujours rêvé » = past perfect : \eng{had always dreamt} (ou \eng{dreamed}) ; \eng{dream of} + \emph{-ing} : \eng{dreamt of becoming}.
\item « après avoir eu son baccalauréat » = \eng{after passing her baccalauréat} (\eng{pass an exam} = « réussir un examen » ; \eng{after} + \emph{-ing}).
\item « afin d'économiser » = \eng{in order to save} (base verbale, jamais *\eng{for to save}).
\item « si fatigant que\ldots » = \eng{so tiring that} ; attention à \eng{tiring} (« fatigant ») et \eng{tired} (« fatigué »).
\item « elle n'a jamais abandonné » : le contexte est un récit passé, prétérit accepté : \eng{she never gave up} (phrasal verb \eng{give up} = « abandonner »).
\item « Je préférerais\ldots » = \eng{I would rather suffer\ldots} (base verbale, sans \eng{to}).
\item « elle a terminé sa formation » (résultat présent) = present perfect : \eng{she has finished her training}.
\item « fière d'elle » = \eng{proud of her} (préposition \eng{of}).
\end{itemize}
\end{corrige}

\begin{corrige}[Exercice 11 — Type Bac : production écrite]
\textbf{Barème indicatif :} 10 points — respect du sujet et de la longueur (2 pts) ; présence des six structures exigées, soulignées (3 pts) ; correction grammaticale et orthographe (3 pts) ; organisation et cohérence (2 pts).

\textbf{Composition modèle (environ 170 mots) ; les structures attendues sont mises en évidence :}
\end{corrige}

\begin{textebox}[Model composition — How my family prepares for Christmas]
Christmas is the most important celebration in my family, and we start preparing for it several weeks in advance. At the beginning of December, my mother saves money \textbf{in order to} buy new clothes for all the children. My father, who \textbf{had never celebrated} Christmas in the city before he married my mother, now enjoys every moment of it.

A few days before Christmas, we travel to the village \textbf{so that} we can spend the feast with our grandparents. The journey is sometimes \textbf{so} long \textbf{that} we arrive completely exhausted. Once there, everybody has a task: the women cook, the men clean the compound, and the children decorate the house. We \textbf{mustn't} forget anything, because the whole family is coming. My mother always says, “\textbf{Let's not} argue during Christmas; let's be united.”

Personally, I \textbf{would rather} help in the kitchen than go to the farm, because I love the smell of the food. This year, we \textbf{have bought} a goat and many presents already. I am sure it will be a wonderful Christmas.
\end{textebox}

\begin{corrige}[Exercice 11 — Points de vigilance]
\begin{itemize}
\item Vérifier que chaque structure imposée apparaît au moins une fois et qu'elle est \emph{soulignée} sur la copie.
\item \eng{Let's not} + base verbale ; \eng{would rather} + base verbale ; \eng{mustn't} + base verbale : aucun \eng{to} après ces expressions.
\item Present perfect pour les actions récentes liées au présent (\eng{we have bought}) ; past perfect pour l'antériorité (\eng{had never celebrated\ldots before he married}).
\item Éviter les calques du français : « fêter Noël » = \eng{celebrate Christmas}, « en avance » = \eng{in advance}, « toute la famille » = \eng{the whole family}.
\item Compter les mots : rester entre 150 et 200 mots ; une copie trop courte perd des points.
\end{itemize}
\end{corrige}

\newpage

\coursec{Fiche bilan}

\courssub{Carte mentale de la leçon}

\begin{popfigure}
\begin{tikzpicture}[mindmap, grow cyclic, every node/.style={concept}, text=white,
  root concept/.append style={concept color=popblue, font=\bfseries\small},
  level 1/.append style={level distance=3.2cm, sibling angle=60, font=\small},
  level 1 concept/.append style={minimum size=2.2cm}]
\node[root concept, align=center]{Grammar\\ ENA05}
  child[concept color=poporange] { node[align=center]{Purpose\\ \eng{to, in order to, so as to, so that}} }
  child[concept color=popgreen] { node[align=center]{Result\\ \eng{so, so...that, such...that, therefore}} }
  child[concept color=poppurple] { node[align=center]{Obligation\\ \eng{must, have to, mustn't}} }
  child[concept color=poppink] { node[align=center]{Preference\\ \eng{prefer, would rather, would prefer}} }
  child[concept color=popgold] { node[align=center]{Suggestions\\ \eng{let's / let's not + BV}} }
  child[concept color=popdark] { node[align=center]{Perfect tenses\\ \eng{present / past perfect}} };
\end{tikzpicture}
\legende{Carte mentale : les six points de grammaire de la leçon ENA05.}
\end{popfigure}

\begin{bilan}
\textbf{1. Lexique clé (English / Français)}
\begin{itemize}
  \item \eng{purpose} (n.) : but, objectif ; \eng{result} (n.) : résultat, conséquence
  \item \eng{in order to} : afin de ; \eng{so that} : afin que, pour que
  \item \eng{therefore} : par conséquent ; \eng{as a result} : en conséquence
  \item \eng{obligation} (n.) : obligation ; \eng{compulsory} (adj.) : obligatoire
  \item \eng{preference} (n.) : préférence ; \eng{to forbid} : interdire
  \item \eng{bare infinitive} : base verbale (infinitif sans \eng{to})
\end{itemize}

\textbf{2. Règles à connaître par cœur}
\begin{center}
\begin{tabularx}{\linewidth}{|l|X|X|}\hline
\rowcolor{popblueL} \textbf{Notion} & \textbf{Structure} & \textbf{Exemple} \\ \hline
But & \eng{to / in order to / so as to + BV} ; négation : \eng{in order not to} & \eng{She works hard in order to succeed.} \\ \hline
But (2 sujets) & \eng{so that + sujet + will/would/can/could + BV} & \eng{He spoke slowly so that we could understand.} \\ \hline
Conséquence & \eng{so + adj/adv + that} ; \eng{such (a) + nom + that} & \eng{It was so hot that we stayed inside.} \\ \hline
Obligation & \eng{must / have to + BV} ; interdiction : \eng{mustn't + BV} & \eng{Students must wear a uniform.} \\ \hline
Absence d'obligation & \eng{don't have to / needn't + BV} & \eng{You don't have to come early.} \\ \hline
Préférence & \eng{prefer + V-ing / would rather + BV / would prefer + to + BV} & \eng{I would rather stay at home.} \\ \hline
Suggestion & \eng{let's / let's not + BV} & \eng{Let's not waste time.} \\ \hline
Present perfect & \eng{have/has + participe passé} & \eng{I have just finished my homework.} \\ \hline
Past perfect & \eng{had + participe passé} & \eng{The bus had left when we arrived.} \\ \hline
\end{tabularx}
\end{center}

\textbf{3. Méthodes express}
\begin{itemize}
  \item \textbf{But ou conséquence ?} Le but répond à « pourquoi ? » (\eng{to, so that}) ; la conséquence répond à « avec quel résultat ? » (\eng{so...that, therefore}).
  \item \textbf{Deux sujets différents ?} Impossible d'utiliser \eng{to} : il faut \eng{so that + proposition}.
  \item \textbf{Chronologie passée :} l'action la plus ancienne se met au past perfect (\eng{had + p.p.}), la plus récente au preterit.
  \item \textbf{Present perfect :} toujours un lien avec le présent (résultat visible, expérience, action récente avec \eng{just, already, yet}).
\end{itemize}
\end{bilan}

\begin{aretenir}[Checklist avant le Bac]
\begin{itemize}
  \item $\square$ Je sais exprimer le but avec \eng{to / in order to / so as to + base verbale}.
  \item $\square$ Je sais former la négation : \eng{in order not to / so as not to}.
  \item $\square$ J'utilise \eng{so that + sujet + modal} quand les deux sujets sont différents.
  \item $\square$ Je distingue \eng{so + adjectif + that} et \eng{such (a) + nom + that}.
  \item $\square$ Je ne confonds pas \eng{mustn't} (interdiction) et \eng{don't have to} (absence d'obligation).
  \item $\square$ Je sais construire \eng{would rather + BV} et \eng{prefer + V-ing}.
  \item $\square$ J'utilise \eng{let's / let's not + base verbale} sans \eng{to}.
  \item $\square$ Je forme le present perfect : \eng{have/has + participe passé} (attention aux participes irréguliers : \eng{gone, written, eaten}).
  \item $\square$ Je place correctement \eng{just, already, yet, ever, never} avec le present perfect.
  \item $\square$ J'utilise le past perfect (\eng{had + p.p.}) pour l'action antérieure à une autre action passée.
  \item $\square$ Je n'écris jamais \eng{for to} pour traduire « pour » + infinitif : c'est simplement \eng{to}.
  \item $\square$ Je vérifie la concordance des temps : \eng{so that + would/could} après un passé.
\end{itemize}
\end{aretenir}

\findecours

\end{document}
